\documentclass[a4paper,11pt]{amsart}

% =========================================================================
%  Electromagnetismo I — Tema 1: Leyes básicas de la electrostática
%  Apuntes formalizados. Compilar con: pdflatex + biber + pdflatex ×2
%  (o latexmk -pdf main.tex).  Todo el preámbulo está en preambulo.tex.
% =========================================================================
% =========================================================================
%  Electromagnetismo I — Tema 1: Leyes básicas de la electrostática
%  Estilo: artículo científico minimalista, heredado del informe «Efecto
%  Doppler con ultrasonidos» (amsart, Palatino, versalitas, enunciados
%  numerados por sección, paleta de azules y grises sobrios), con la
%  notación vectorial de «Cálculo integral multivariable» (L. Navas) y las
%  expresiones del formulario del curso.
% =========================================================================

% --- IDIOMA Y CODIFICACIÓN ------------------------------------------------
\usepackage[utf8]{inputenc}
\usepackage[T1]{fontenc}
\usepackage[spanish,es-tabla,es-nodecimaldot]{babel}

% --- TIPOGRAFÍA (Palatino, como el documento de referencia) ---------------
\usepackage{mathpazo}
\linespread{1.05}
\emergencystretch=2em
\usepackage[final]{microtype}

% --- MATEMÁTICAS -------------------------------------------------------------
\usepackage{amsmath, amssymb, mathtools}
\usepackage{bm}
\usepackage{mathrsfs}
\usepackage{esint}

% --- PÁGINA, COLORES Y GRÁFICOS --------------------------------------------
\usepackage[a4paper,margin=2.6cm]{geometry}
\usepackage{graphicx}
\usepackage{xcolor}
% Paleta base (idéntica a la del informe de ecuaciones diferenciales)
\definecolor{azulprof}{HTML}{1F3A5F}    % azul oscuro (EDO)
\definecolor{azulacero}{HTML}{3B6EA5}   % azul medio
\definecolor{azulclaro}{HTML}{8FB3D9}   % azul claro
\definecolor{azulfondo}{HTML}{F3F7FB}   % fondo tenue de definiciones
\definecolor{grispizarra}{HTML}{6B7280} % gris
\definecolor{grisclaro}{HTML}{D4D8DE}   % gris claro
\definecolor{acento}{HTML}{B5651D}      % acento (ocre quemado)
% Paleta ampliada para campos escalares y vectoriales (magnitud creciente:
% índigo -> púrpura -> carmín -> naranja -> ámbar)
\definecolor{cm0}{RGB}{38,30,112}
\definecolor{cm1}{RGB}{112,36,140}
\definecolor{cm2}{RGB}{190,52,118}
\definecolor{cm3}{RGB}{236,104,66}
\definecolor{cm4}{RGB}{252,184,42}
% Colores de apoyo de los diagramas geométricos (tonos sobrios)
\definecolor{rojoteo}{HTML}{B5651D}     % trayectorias, circuitos, resultados
\definecolor{verdeteo}{HTML}{2F6F68}    % versores y normales
\definecolor{ocreteo}{HTML}{8FB3D9}     % superficies y volúmenes (relleno)

\usepackage{tikz}
\usepackage{tikz-3dplot}
\usetikzlibrary{babel, arrows.meta, calc, positioning, decorations.pathmorphing,
                decorations.markings, patterns, 3d, intersections, angles, quotes,
                fadings, shapes.geometric}
\usepackage{pgfplots}
\pgfplotsset{compat=1.17}
\usepgfplotslibrary{colormaps, patchplots, fillbetween}
\pgfplotsset{colormap={campo}{
  rgb255=(38,30,112) rgb255=(112,36,140) rgb255=(190,52,118)
  rgb255=(236,104,66) rgb255=(252,184,42) rgb255=(255,226,110)}}
\pgfplotsset{colormap={campoclaro}{
  rgb255=(232,236,240) rgb255=(183,203,227) rgb255=(59,110,168) rgb255=(27,54,93)}}
\pgfplotsset{colormap={verdeamarillo}{
  rgb255=(236,243,250) rgb255=(143,179,217) rgb255=(59,110,165) rgb255=(31,58,95)}}
\usepackage{booktabs}
\usepackage{needspace}
\usepackage{array}
\usepackage{enumitem}
\setlist[itemize]{leftmargin=1.6em, itemsep=2pt, topsep=4pt}
\setlist[enumerate]{leftmargin=1.8em, itemsep=2pt, topsep=4pt}

% --- ENLACES Y REFERENCIAS CRUZADAS ----------------------------------------
\usepackage[colorlinks=true, linkcolor=azulacero, citecolor=azulacero,
            urlcolor=azulacero, plainpages=false, pdfusetitle=false,
            pdftitle={Tema 1. Leyes básicas de la electrostática},
            pdfauthor={Daniel Martín Guerrero}]{hyperref}
\usepackage[spanish,nameinlink]{cleveref}

% --- BIBLIOGRAFÍA -----------------------------------------------------------
\usepackage{csquotes}
\usepackage[backend=biber, style=numeric, sorting=none, giveninits=true,
            maxbibnames=6, minbibnames=3, doi=true, url=true,
            isbn=false]{biblatex}
\addbibresource{referencias.bib}
\DeclareFieldFormat[article,techreport,unpublished,book]{title}{#1\isdot}
\DeclareFieldFormat{url}{\url{#1}}
\renewcommand*{\bibfont}{\small}
\setlength{\bibitemsep}{3pt}

% --- CABECERAS ---------------------------------------------------------------
\pagestyle{headings}
\makeatletter
\def\ps@headings{%
  \let\@mkboth\markboth
  \def\@evenhead{\small\scshape\normalfont\thepage\hfil Electromagnetismo I}%
  \def\@oddhead{\small\scshape\normalfont Tema 1. Leyes básicas de la electrostática\hfil\thepage}%
  \let\@evenfoot\@empty \let\@oddfoot\@empty}
\makeatother
\numberwithin{equation}{section}
% Subsecciones como títulos en bloque (no «run-in»), con aire antes y después
\usepackage{titlesec}
\titleformat{\subsection}[block]{\normalfont\bfseries\color{azulprof}}{\thesubsection.}{0.6em}{}
\titlespacing*{\subsection}{0pt}{3ex plus 1ex minus .5ex}{1.6ex plus .4ex}
\setcounter{tocdepth}{2}
\makeatletter
\def\@@and{y}
\renewcommand{\keywordsname}{Palabras clave}
\makeatother

% =========================================================================
%               NOTACIÓN Y MACROS DE USO FRECUENTE
% =========================================================================
% --- Diferencial, derivadas parciales y números -----------------------------
\newcommand{\dif}{\mathrm{d}}
\newcommand{\pd}[2]{\frac{\partial #1}{\partial #2}}
\newcommand{\pdd}[2]{\frac{\partial^{2} #1}{\partial #2^{2}}}
\newcommand{\pdm}[3]{\frac{\partial^{2} #1}{\partial #2\,\partial #3}}
\newcommand{\R}{\mathbb{R}}
\newcommand{\N}{\mathbb{N}}
\newcommand{\Z}{\mathbb{Z}}
\newcommand{\eul}{\mathrm{e}}

% --- Vectores, versores y operador nabla ---------------------------------
\newcommand{\vc}[1]{\vec{#1}}                       % vector con flecha
\newcommand{\uh}[1]{\hat{u}_{#1}}                   % versor û_#1
\newcommand{\ux}{\uh{x}}\newcommand{\uy}{\uh{y}}\newcommand{\uz}{\uh{z}}
\newcommand{\ur}{\uh{r}}\newcommand{\uth}{\uh{\theta}}\newcommand{\uph}{\uh{\varphi}}
\newcommand{\urho}{\uh{\rho}}\newcommand{\uR}{\uh{R}}\newcommand{\un}{\hat{n}}
\newcommand{\nab}{\vec{\nabla}}                     % ∇⃗
\newcommand{\nabp}{\vec{\nabla}'}                   % ∇⃗′ (coordenadas fuente)
\newcommand{\Lap}{\nabla^{2}}                       % laplaciano
\newcommand{\Lapp}{\nabla'^{2}}                      % laplaciano en coordenadas fuente
\newcommand{\norm}[1]{\lVert #1\rVert}
\newcommand{\abs}[1]{\lvert #1\rvert}
\newcommand{\prodesc}[2]{\langle #1,#2\rangle}
\newcommand{\rp}{\vc{r}\,'}                         % r⃗′ punto fuente
\newcommand{\dl}{\dif\vc{\ell}}                     % elemento de línea dirigido
\newcommand{\dS}{\dif\vc{S}}                        % elemento de área dirigido
\newcommand{\dsc}{\dif S}                           % elemento de área escalar
\newcommand{\dV}{\dif V}                            % elemento de volumen
\newcommand{\defeq}{\mathrel{\overset{\mathrm{def}}{=}}}
\newcommand{\bordeS}{\partial S}
\newcommand{\bordeV}{\partial V}
\DeclareMathOperator{\grad}{grad}
\DeclareMathOperator{\dive}{div}
\DeclareMathOperator{\rot}{rot}
\DeclareMathOperator{\sen}{sen}
\DeclareMathOperator{\senh}{senh}
\DeclareMathOperator{\sgn}{sgn}
\DeclareMathOperator{\Tr}{Tr}
\renewcommand{\sin}{\sen}                            % «sen» en todo el texto

% Integrales cerradas de superficie (∯) y de volumen cerrado (∰): paquete esint, como en L. Navas

% Tensor de Levi-Civita y delta de Kronecker
\newcommand{\eps}{\varepsilon}
\newcommand{\dpart}[1]{\partial_{#1}}

% --- Notación específica del Tema 1 (electrostática) -------------------------
\newcommand{\eo}{\varepsilon_{0}}                    % permitividad del vacío
\newcommand{\qp}{q\,'}                              % carga fuente q'
\newcommand{\pot}{\phi}                              % potencial electrostático
\newcommand{\flujo}{\Phi_{\!E}}                      % flujo eléctrico
\newcommand{\dOm}{\dif\Omega}                       % elemento de ángulo sólido
\newcommand{\Qenc}{Q_{\mathrm{enc}}}                 % carga encerrada
\newcommand{\Ke}{\frac{1}{4\pi\eo}}                  % constante de Coulomb 1/(4 pi eps0)
\newcommand{\vE}{\vc{E}}
\newcommand{\vF}{\vc{F}}
\newcommand{\dlp}{\dif l\,'}                        % elemento de longitud fuente
\newcommand{\dSp}{\dif S\,'}                        % elemento de superficie fuente
\newcommand{\dVp}{\dif V\,'}                        % elemento de volumen fuente
\newcommand{\ent}[1]{\bigl\lvert #1\bigr\rvert}    % módulo grande
% Estilos de cargas puntuales para los diagramas
\tikzset{
  cargapos/.style={circle, fill=azulacero, draw=azulprof, line width=0.4pt, inner sep=0pt,
                   minimum size=9pt, text=white, font=\scriptsize\bfseries},
  carganeg/.style={circle, fill=acento, draw=acento!70!black, line width=0.4pt, inner sep=0pt,
                   minimum size=9pt, text=white, font=\scriptsize\bfseries},
  cargapru/.style={circle, fill=white, draw=azulprof, line width=0.6pt, inner sep=0pt,
                   minimum size=9pt, text=azulprof, font=\scriptsize\bfseries},
}

% Texto de relleno discreto para datos pendientes
\newcommand{\pend}[1]{\textcolor{grispizarra}{\colorbox{grisclaro}{\strut\,#1\,}}}

% =========================================================================
%        ENTORNOS FORMALES (enunciados numerados, estilo artículo)
% =========================================================================
\usepackage{amsthm}
\usepackage[most]{tcolorbox}

% Enunciados en cursiva (resultados generales)
\newtheoremstyle{cienplano}{10pt}{10pt}{\itshape}{}%
  {\bfseries\color{azulprof}}{.}{0.6em}%
  {\thmname{#1}\thmnumber{ #2}\thmnote{\normalfont\upshape\ (#3)}}
% Definiciones, cálculos y ejemplos (cuerpo en redonda)
\newtheoremstyle{ciendef}{10pt}{10pt}{\normalfont}{}%
  {\bfseries\color{azulprof}}{.}{0.6em}%
  {\thmname{#1}\thmnumber{ #2}\thmnote{\normalfont\upshape\ (#3)}}
% Observaciones y notas (encabezado en cursiva)
\newtheoremstyle{cienobs}{8pt}{8pt}{\normalfont}{}%
  {\itshape\color{azulprof}}{.}{0.6em}%
  {\thmname{#1}\thmnumber{ #2}\thmnote{\normalfont\upshape\ (#3)}}

\theoremstyle{cienplano}
\newtheorem{proposicion}{Proposición}[section]
\newtheorem{teorema}[proposicion]{Teorema}
\newtheorem{lema}[proposicion]{Lema}
\newtheorem{corolario}[proposicion]{Corolario}
\theoremstyle{ciendef}
\newtheorem{notacion}[proposicion]{Notación}
\newtheorem{ejemplo}[proposicion]{Ejemplo}
\newtheorem{ejercicio}[proposicion]{Ejercicio}
% Definiciones: único entorno enmarcado (como en el tema de EDO): caja fina con
% fondo tenue y título sobre el borde. Comparte contador con los enunciados.
\newcounter{definicion}
\renewcommand{\thedefinicion}{\theproposicion}
\newtcolorbox{definicionbox}[1]{%
  enhanced, breakable, sharp corners,
  colback=azulfondo, colframe=azulclaro, boxrule=0.4pt,
  left=9pt, right=9pt, top=5pt, bottom=5pt,
  before skip=12pt, after skip=12pt,
  attach boxed title to top left={yshift*=-\tcboxedtitleheight/2, xshift=8pt},
  boxed title style={colback=white, colframe=white, sharp corners, boxrule=0pt,
                     left=3pt, right=3pt, top=1pt, bottom=1pt},
  fonttitle=\bfseries\color{azulprof}, coltitle=azulprof,
  title={Definición~\thedefinicion\ifx\relax#1\relax\else\ (#1)\fi}}
\newenvironment{definicion}[1][]
  {\par\stepcounter{proposicion}\refstepcounter{definicion}\begin{definicionbox}{#1}}
  {\end{definicionbox}}

\theoremstyle{cienobs}
\newtheorem{observacion}[proposicion]{Observación}

% Resolución de ejercicios: misma maquinaria que «proof», con otro rótulo
\renewcommand{\qedsymbol}{\ensuremath{\blacksquare}}
\newenvironment{resolucion}[1][Resolución]
  {\begin{proof}[\itshape\color{azulprof}#1]}{\end{proof}}

\crefname{proposicion}{proposición}{proposiciones}
\Crefname{proposicion}{Proposición}{Proposiciones}
\crefname{teorema}{teorema}{teoremas}
\Crefname{teorema}{Teorema}{Teoremas}
\crefname{lema}{lema}{lemas}
\Crefname{lema}{Lema}{Lemas}
\crefname{corolario}{corolario}{corolarios}
\Crefname{corolario}{Corolario}{Corolarios}
\crefname{definicion}{definición}{definiciones}
\Crefname{definicion}{Definición}{Definiciones}
\crefname{notacion}{notación}{notaciones}
\Crefname{notacion}{Notación}{Notaciones}
\crefname{ejemplo}{ejemplo}{ejemplos}
\Crefname{ejemplo}{Ejemplo}{Ejemplos}
\crefname{ejercicio}{ejercicio}{ejercicios}
\Crefname{ejercicio}{Ejercicio}{Ejercicios}
\crefname{observacion}{observación}{observaciones}
\Crefname{observacion}{Observación}{Observaciones}
\crefname{section}{sección}{secciones}
\Crefname{section}{Sección}{Secciones}
\crefname{subsection}{sección}{secciones}
\Crefname{subsection}{Sección}{Secciones}
\crefname{table}{tabla}{tablas}
\Crefname{table}{Tabla}{Tablas}
\crefname{figure}{figura}{figuras}
\Crefname{figure}{Figura}{Figuras}
\crefname{equation}{ecuación}{ecuaciones}
\Crefname{equation}{Ecuación}{Ecuaciones}
\crefformat{equation}{ecuación~(#2#1#3)}
\Crefformat{equation}{Ecuación~(#2#1#3)}
\crefrangeformat{equation}{ecuaciones~(#3#1#4)--(#5#2#6)}
\crefmultiformat{equation}{ecuaciones~(#2#1#3)}{ y~(#2#1#3)}{, (#2#1#3)}{ y~(#2#1#3)}

% Pies de figura y tabla, en la línea del documento de referencia
\usepackage[font=small, labelfont={sc,color=azulprof}, labelsep=period,
            margin=0.6cm, justification=justified]{caption}

% =========================================================================
%        ESTILOS GRÁFICOS COMUNES (TikZ / pgfplots)
% =========================================================================
% =========================================================================
%  Estilos y macros gráficas comunes
%  - \mapcolor{t}: escala cromática «campo» (t en [0,1]); define el color mapc
%  - \flechas{u}{v}{lista x}{lista y}{escala}{|F|max}: campo vectorial 2D
%  - \flechasIII{u}{v}{w}{...}: campo vectorial 3D (con tikz-3dplot)
% =========================================================================

% --- Escala cromática (5 tramos entre cm0 y cm4) ------------------------
\newcommand{\mapcolor}[1]{%
  \pgfmathsetmacro{\mt}{min(max(#1,0),0.9999)*4}%
  \pgfmathtruncatemacro{\mi}{floor(\mt)}%
  \pgfmathtruncatemacro{\mj}{\mi+1}%
  \pgfmathsetmacro{\mf}{100*(\mt-\mi)}%
  \colorlet{mapc}{cm\mj!\mf!cm\mi}}

% --- Estilos TikZ generales ------------------------------------------------
\tikzset{
  >=Stealth,
  eje/.style={->, line width=0.6pt, grispizarra},
  ejeprof/.style={->, line width=0.7pt, azulprof},
  vect/.style={->, line width=1pt, azulprof},
  vectrojo/.style={->, line width=1pt, rojoteo},
  vectverde/.style={->, line width=1pt, verdeteo},
  fino/.style={line width=0.4pt, grispizarra},
  punteado/.style={densely dotted, line width=0.5pt, grispizarra},
  etiq/.style={font=\small, text=black!85},
  cuadro/.style={draw=grispizarra!60, line width=0.4pt, fill=white, inner sep=2pt},
  trayecto/.style={line width=1.1pt, rojoteo},
  orienta/.style={decoration={markings, mark=at position 0.55 with {\arrow{Stealth[length=5pt]}}}, postaction={decorate}},
}

% --- Marco con marcas para gráficas de campos 2D ---------------------------
% \marco2D{xmin}{xmax}{ymin}{ymax}{xticks}{yticks}  (unidades TikZ = unidades del campo)
\newcommand{\marcoD}[6]{%
  \draw[grispizarra, line width=0.5pt] (#1,#3) rectangle (#2,#4);
  \foreach \t in {#5}{%
     \draw[grispizarra, line width=0.4pt] (\t,#3) -- ++(0,0.08);
     \draw[grispizarra, line width=0.4pt] (\t,#4) -- ++(0,-0.08);
     \node[below=1pt, font=\scriptsize] at (\t,#3) {\pgfmathprintnumber[fixed,precision=1,use comma,zerofill=false]{\t}};}%
  \foreach \t in {#6}{%
     \draw[grispizarra, line width=0.4pt] (#1,\t) -- ++(0.08,0);
     \draw[grispizarra, line width=0.4pt] (#2,\t) -- ++(-0.08,0);
     \node[left=1pt, font=\scriptsize] at (#1,\t) {\pgfmathprintnumber[fixed,precision=1,use comma,zerofill=false]{\t}};}%
}

% --- Campo vectorial plano: flechas centradas en cada nudo de la malla ---
% #1,#2: componentes u(x,y), v(x,y) (en términos de \x e \y)
% #3,#4: listas de abscisas y ordenadas; #5: factor de escala de longitud
% #6: módulo máximo (normaliza el color)
\newcommand{\flechas}[6]{%
  \foreach \x in {#3}{\foreach \y in {#4}{%
    \pgfmathsetmacro{\fu}{#1}\pgfmathsetmacro{\fv}{#2}%
    \pgfmathsetmacro{\fm}{sqrt(\fu*\fu+\fv*\fv)}%
    \pgfmathsetmacro{\ft}{\fm/(#6)}%
    \mapcolor{\ft}%
    \ifdim\fm pt>0.001pt
      \pgfmathsetmacro{\hl}{2.2+4.2*\ft}%
      \pgfmathsetmacro{\lw}{0.35+0.5*\ft}%
      \draw[-{Stealth[length=\hl pt,width=\hl pt]}, line width=\lw pt, mapc]
        (\x-0.5*#5*\fu,\y-0.5*#5*\fv) -- ++(#5*\fu,#5*\fv);
    \else
      \fill[mapc] (\x,\y) circle (0.7pt);
    \fi}}}

% --- Campo vectorial espacial (tikz-3dplot): lo mismo en 3D ----------------
\newcommand{\flechasIII}[8]{%
  \foreach \x in {#4}{\foreach \y in {#5}{\foreach \z in {#6}{%
    \pgfmathsetmacro{\fu}{#1}\pgfmathsetmacro{\fv}{#2}\pgfmathsetmacro{\fw}{#3}%
    \pgfmathsetmacro{\fm}{sqrt(\fu*\fu+\fv*\fv+\fw*\fw)}%
    \pgfmathsetmacro{\ft}{\fm/(#8)}%
    \mapcolor{\ft}%
    \ifdim\fm pt>0.001pt
      \pgfmathsetmacro{\hl}{2.0+3.6*\ft}%
      \pgfmathsetmacro{\lw}{0.35+0.5*\ft}%
      \draw[-{Stealth[length=\hl pt,width=\hl pt]}, line width=\lw pt, mapc]
        (\x-0.5*#7*\fu,\y-0.5*#7*\fv,\z-0.5*#7*\fw) -- ++(#7*\fu,#7*\fv,#7*\fw);
    \else
      \fill[mapc] (\x,\y,\z) circle (0.7pt);
    \fi}}}}

% --- Estilo base de ejes pgfplots (como en la figura del efecto Doppler) --
\pgfplotsset{
  tema/.style={
    tick label style={font=\scriptsize}, label style={font=\small},
    axis lines=left, axis line style={grispizarra},
    grid=major, grid style={grisclaro, line width=0.4pt},
    every tick/.style={grispizarra},
    legend style={font=\scriptsize, draw=none, fill=white, fill opacity=0.85,
                  text opacity=1},
    /pgf/number format/use comma,
  },
  mapa/.style={
    tick label style={font=\scriptsize}, label style={font=\small},
    axis line style={grispizarra}, every tick/.style={grispizarra},
    /pgf/number format/use comma,
    colormap name=campo,
  },
  superficie/.style={
    tick label style={font=\scriptsize}, label style={font=\small},
    axis line style={grispizarra}, every tick/.style={grispizarra},
    grid=none, /pgf/number format/use comma,
    colormap name=campo,
  },
}

% Rótulos de ejes para marcoD:  \ejesD{xmin}{xmax}{ymin}{ymax}
\newcommand{\ejesD}[4]{%
  \node[font=\small, anchor=north] at ($({(#1+#2)/2},#3)+(0,-13pt)$) {$x$ (u.\,arb.)};
  \node[font=\small, rotate=90, anchor=south] at ($(#1,{(#3+#4)/2})+(-17pt,0)$) {$y$ (u.\,arb.)};}

% Recorta el contenido #5 al rectángulo [#1,#2]x[#3,#4]
\newcommand{\clipD}[5]{\begin{scope}\clip (#1,#3) rectangle (#2,#4); #5 \end{scope}}

% Versor dibujado en el punto (#1,#2,#3) con desplazamiento #4 (coordenadas 3D)
\newcommand{\uvect}[5]{%  versor en P: origen (#1,#2,#3), dirección (#4) y rótulo (#5)
  \draw[-{Stealth[length=4.5pt,width=4pt]}, rojoteo, line width=1.1pt] (#1,#2,#3) -- ++#4 node[#5] {};}

% Mapa de calor de f=exp(-r^2) (anillos concéntricos) con curvas de nivel y marco; #1 = extras
\newcommand{\mapaGauss}[1]{%
  \begin{scope}
    \clip (-2,-2) rectangle (2,2);
    \foreach \k in {0,...,90}{%
      \pgfmathsetmacro{\rk}{3.0*(1-\k/91)}%
      \pgfmathsetmacro{\fk}{exp(-\rk*\rk)}%
      \mapcolor{\fk}%
      \fill[mapc] (0,0) circle (\rk);}
    \foreach \lv in {0.9,0.8,0.7,0.6,0.5,0.4,0.3,0.2,0.1}{%
      \pgfmathsetmacro{\rl}{sqrt(-ln(\lv))}%
      \draw[black!75, line width=0.45pt] (0,0) circle (\rl);}
    #1
  \end{scope}
  \marcoD{-2}{2}{-2}{2}{-2,-1,0,1,2}{-2,-1,0,1,2}
  \ejesD{-2}{2}{-2}{2}}




% =========================================================================
%                     VARIABLES PARA LA PORTADA
% =========================================================================
\newcommand{\Universidad}{Universidad de Salamanca}
\newcommand{\Facultad}{Facultad de Ciencias}
\newcommand{\Titulacion}{Grado en Física}
\newcommand{\Asignatura}{Electromagnetismo I}
\newcommand{\NumTema}{Tema 1}
\newcommand{\TituloTema}{Leyes básicas\\de la electrostática}
\newcommand{\tipotrabajo}{Apuntes formalizados con seminarios prácticos}
\newcommand{\miNombre}{Daniel Martín Guerrero}
\newcommand{\miProfesor}{Luis López Díaz}
\newcommand{\SubtituloTema}{Ley de Coulomb, campo eléctrico, ley de Gauss,\\potencial y condiciones de contorno}
\newcommand{\miUbicacion}{Salamanca, curso 2026/2027}

% =========================================================================
%                     DATOS DEL ARTÍCULO (cabecera)
% =========================================================================
\title{Leyes básicas de la electrostática en el vacío}
\author{Daniel Martín Guerrero}
\address{Grado en Física, Facultad de Ciencias, Universidad de Salamanca,
         Salamanca, España}
\date{}

% =========================================================================
%                                DOCUMENTO
% =========================================================================
\begin{document}
% --- PORTADA ----------------------------------------------------------------
\pagenumbering{alph}
%%%%%%%%%%%%%%%%%%%%%%%%%%%%%%%%%%%%%%%%%%%%%%%%%%%%%%%%%%%%%%%%%%%%%%%%
% Portada formal (estilo tesis): logotipos, identificación institucional,
% título y datos de autoría. Banda vertical a la derecha, como en el
% «Cálculo integral multivariable» de L. Navas. Requiere en main.tex las
% macros \Universidad, \Facultad, \Titulacion, \Asignatura, \NumTema,
% \TituloTema, \SubtituloTema, \tipotrabajo, \miNombre, \miProfesor,
% \miUbicacion.
%%%%%%%%%%%%%%%%%%%%%%%%%%%%%%%%%%%%%%%%%%%%%%%%%%%%%%%%%%%%%%%%%%%%%%%%

\begin{titlepage}
\thispagestyle{empty}
\pagecolor{white}

\begin{tikzpicture}[remember picture, overlay]
  \coordinate (O) at (current page.south west);

  % Velo muy tenue en la mitad superior (cabecera institucional)
  \shade[top color=azulfondo, bottom color=white]
    ($(O)+(0,18.5cm)$) rectangle ($(O)+(21cm,29.7cm)$);

  % Banda vertical azul, a la derecha del bloque de texto
  \fill[azulacero] ($(O)+(19.0cm,2.6cm)$) rectangle ($(O)+(19.14cm,27.4cm)$);

  % Logotipos
  \node[anchor=north west, inner sep=0] at ($(O)+(2.4cm,27.7cm)$)
    {\includegraphics[height=1.55cm]{logo_usal.png}};
  \node[anchor=north east, inner sep=0] at ($(O)+(17.4cm,27.7cm)$)
    {\includegraphics[height=1.55cm]{logo_facultad.png}};
  \draw[grispizarra!60, line width=0.4pt] ($(O)+(2.4cm,25.7cm)$) -- ($(O)+(17.4cm,25.7cm)$);

  % Identificación institucional
  \node[anchor=north, align=center, inner sep=0, text width=15cm] at ($(O)+(9.9cm,25.2cm)$)
    {{\color{azulprof}\scshape\Large\textls[120]{\Universidad}}\\[5pt]
     {\color{grispizarra}\scshape\large\textls[80]{\Facultad}}\\[3pt]
     {\color{grispizarra}\scshape\large\textls[80]{\Titulacion}}};

  % Asignatura y título
  \node[anchor=north, align=center, inner sep=0, text width=15cm] at ($(O)+(9.9cm,19.2cm)$)
    {{\color{azulacero}\scshape\Large\textls[160]{\Asignatura}}\\[4pt]
     {\color{grispizarra}\scshape\large\textls[100]{\NumTema}}\\[22pt]
     {\color{azulprof}\fontsize{34}{40}\selectfont\bfseries\TituloTema}\\[14pt]
     {\color{azulacero}\rule{3.2cm}{0.8pt}}\\[14pt]
     {\color{black!80}\itshape\large\SubtituloTema}};

  % Epígrafe: las leyes básicas del tema
  \node[anchor=north, align=center, inner sep=0, text width=15cm] at ($(O)+(9.9cm,10.4cm)$)
    {\color{grispizarra}\small
     $\displaystyle\vc{F}_{12}=\frac{1}{4\pi\eo}\,\frac{q_1q_2}{R_{12}^{2}}\,\uR\qquad
      \oiint_{S}\vc{E}\cdot\dS=\frac{Q_{\mathrm{enc}}}{\eo}\qquad
      \vc{E}=-\nab\phi$};

  % Autoría y dirección
  \draw[grispizarra!60, line width=0.4pt] ($(O)+(2.4cm,7.6cm)$) -- ($(O)+(17.4cm,7.6cm)$);
  \node[anchor=north west, inner sep=0, text width=7cm, align=left] at ($(O)+(2.4cm,7.2cm)$)
    {{\scshape\small\color{azulacero}\textls[80]{Autor}}\\[4pt]
     {\large\color{black!85}\miNombre}\\[2pt]
     {\small\color{grispizarra}Grado en Física}};
  \node[anchor=north west, inner sep=0, text width=7cm, align=left] at ($(O)+(10.4cm,7.2cm)$)
    {{\scshape\small\color{azulacero}\textls[80]{Profesor}}\\[4pt]
     {\large\color{black!85}\miProfesor}\\[2pt]
     {\small\color{grispizarra}\Asignatura}};

  % Tipo de documento, lugar y curso
  \node[anchor=south, align=center, inner sep=0, text width=15cm] at ($(O)+(9.9cm,2.9cm)$)
    {{\color{azulprof}\scshape\large\textls[100]{\tipotrabajo}}\\[4pt]
     {\color{grispizarra}\scshape\small\textls[80]{\miUbicacion}}};
\end{tikzpicture}

\end{titlepage}

\clearpage
\pagenumbering{arabic}
\setcounter{page}{1}

\begin{abstract}
Formalizamos las leyes básicas de la electrostática en el vacío a partir de
dos datos experimentales: la ley de Coulomb y el principio de superposición.
Tras situar el problema en su contexto histórico y en el marco de las teorías
de campos, definimos el campo eléctrico de distribuciones discretas y continuas
de carga, demostramos la ley de Gauss mediante el ángulo sólido, la expresamos
en forma diferencial con ayuda de la delta de Dirac, y probamos que el campo
electrostático es irrotacional, lo que permite introducir el potencial y las
ecuaciones de Poisson y de Laplace. El tema se cierra con las condiciones de
contorno de los campos al atravesar una superficie cargada. Cada resultado se
enuncia con rigor, se demuestra o se justifica geométricamente, y se acompaña de
diagramas y gráficas; los seminarios finales resuelven, con todo detalle,
los doce problemas de la lista del tema.
\end{abstract}

\maketitle
\tableofcontents

% =========================================================================
\section{Introducción y contexto histórico}\label{sec:introduccion}
% =========================================================================
La electrostática estudia las interacciones entre cargas eléctricas \emph{en
reposo}. Es el escalón más sencillo del electromagnetismo y, a la vez, el lugar
donde aparecen por primera vez las ideas que sostienen toda la teoría: el
\emph{campo} como intermediario de la interacción, las \emph{fuentes} que lo
generan, y la pareja de ecuaciones, una para la divergencia y otra para el
rotacional, que lo determinan por completo. Este tema parte de un único dato
experimental (la ley de Coulomb) y de un principio (la superposición), y de ellos
deduce, con las herramientas del Tema~0~\cite{TEMA0}, la ley de Gauss, el potencial
electrostático y las condiciones que cumplen los campos en una frontera. Seguimos
el programa de los apuntes de la asignatura~\cite{APUNTES-T1} y los desarrollos de
la clase, y nos apoyamos en los textos clásicos~\cite{GRIFFITHS,WANGSNESS,JACKSON,PURCELL}. Las
figuras reproducen, como diagramas en TikZ y gráficas en pgfplots, las del material
de figuras del curso~\cite{FIGURAS-T1} (preparadas con Mathematica) y las
completan con esquemas propios de la teoría y de los problemas~\cite{PROBLEMAS-T1}.

% -------------------------------------------------------------------------
\subsection{El electromagnetismo clásico y sus límites de validez}\label{sec:limites}
% -------------------------------------------------------------------------
Electromagnetismo~I y~II se enmarcan en la \emph{física clásica}. A comienzos del
siglo~XX se comprobó que la mecánica de Newton resulta «inadecuada» fuera de su
dominio natural: para objetos que se mueven a velocidades próximas a la de la luz
hay que sustituirla por la relatividad especial; para objetos muy pequeños, a escala
atómica, por la mecánica cuántica; y cuando ambas circunstancias se dan a la vez, por
la teoría cuántica de campos (\cref{fig:teorias}).

% Figura: las cuatro teorías de la física y el dominio de la electrodinámica clásica.
\begin{figure}[ht]
\centering
\begin{tikzpicture}[
    caja/.style={draw=azulclaro, fill=azulfondo, line width=0.5pt, text width=4.6cm,
                 align=center, inner sep=6pt, font=\small},
    caja0/.style={caja, draw=azulacero, line width=0.9pt},
    flecha/.style={->, line width=0.7pt, azulacero, shorten >=2pt, shorten <=2pt},
    rotulo/.style={font=\scriptsize, text=grispizarra, align=center, fill=white, inner sep=1.5pt}]
  \node[caja0] (A) at (0,0) {\textbf{Mecánica clásica}\\[1pt]{\footnotesize Newton}};
  \node[caja]  (B) at (8.8,0) {\textbf{Mecánica cuántica}\\[1pt]{\footnotesize Bohr, Heisenberg, Schrödinger}};
  \node[caja]  (C) at (0,-2.5) {\textbf{Relatividad especial}\\[1pt]{\footnotesize Einstein}};
  \node[caja]  (D) at (8.8,-2.5) {\textbf{Teoría cuántica de campos}\\[1pt]{\footnotesize Dirac, Pauli, Feynman}};
  \draw[flecha] (A) -- node[rotulo] {objetos muy pequeños\\$S\sim\hbar$} (B);
  \draw[flecha] (A) -- node[rotulo] {velocidades $v\sim c$} (C);
  \draw[flecha] (C) -- node[rotulo] {objetos muy pequeños\\$S\sim\hbar$} (D);
  \draw[flecha] (B) -- node[rotulo] {velocidades $v\sim c$} (D);
\end{tikzpicture}
\caption{Dominio de validez de las teorías de la física. El recuadro resaltado es la
física clásica, que incluye el electromagnetismo de Maxwell; este, a diferencia de la
mecánica de Newton, ya es compatible con la relatividad especial y se funde con la
cuántica en la electrodinámica cuántica.}
\label{fig:teorias}
\end{figure}


La pregunta natural es entonces si tiene sentido emplear la descripción
\emph{ondulatoria} de Maxwell. Hoy sabemos que toda radiación electromagnética
admite una descripción corpuscular (fotones) y otra ondulatoria, y que cuál de las
dos resulta más útil depende de la energía de los fotones. A frecuencias bajas, un
fotón transporta tan poca energía que cualquier detector responde a un número
enorme de ellos y solo registra su promedio: el campo se comporta como una onda
continua. Al subir la frecuencia, la naturaleza corpuscular se vuelve evidente
(por ejemplo, en los experimentos con rayos~X).

\begin{ejemplo}[Cuántos fotones hay en una onda clásica]\label{ej:fotones}
El flujo de fotones a una distancia $d$ de una fuente isótropa de potencia $P$ y
frecuencia $f$ es $\Phi_\gamma=P/(4\pi d^{2}\,hf)$, con $h=6{,}63\cdot10^{-34}\,\mathrm{J\,s}$.
\begin{itemize}
  \item \emph{Antena de FM}: $f=200\,\mathrm{MHz}$, $P=10\,\mathrm{kW}$, $d=10\,\mathrm{km}$.
        Con $hf\simeq1{,}3\cdot10^{-25}\,\mathrm{J}$ resulta
        $\Phi_\gamma\sim10^{19}$--$10^{20}\ \mathrm{fotones/(m^{2}\,s)}$.
  \item \emph{Bombilla}: $P=100\,\mathrm{W}$, luz visible ($\lambda\sim500\,\mathrm{nm}$,
        $f\sim10^{14}\,\mathrm{Hz}$, $hf\simeq4\cdot10^{-19}\,\mathrm{J}$), a $d=1\,\mathrm{m}$:
        $\Phi_\gamma\sim2\cdot10^{19}\ \mathrm{fotones/(m^{2}\,s)}$.
\end{itemize}
En ambos casos el sensor recibe una cantidad tan grande de fotones que responde a su
promedio, como lo haría ante una onda continua, y el planteamiento ondulatorio es
legítimo.
\end{ejemplo}

Dos hechos explican por qué la descripción clásica sigue siendo esencial en la física
moderna. Primero, la relatividad especial se apoyó en experimentos electromagnéticos y las
ecuaciones de Maxwell resultaron ser ya compatibles con ella. Segundo, las correcciones
cuánticas a la fuerza electromagnética son muy pequeñas incluso a distancias del orden de
$10^{-10}\,\mathrm{cm}$, cien veces menores que el tamaño de un átomo: las mismas leyes
describen la atracción y la repulsión entre partículas dentro del átomo y entre cuerpos
macroscópicos (aunque luego sea preciso recurrir a la mecánica cuántica para predecir cómo
se mueven esas partículas bajo tales fuerzas) y, a distancias aún menores, la teoría
electromagnética se funde con la cuántica en la electrodinámica cuántica. Por «clásico»
entendemos, por tanto, «no cuántico»: todas las magnitudes (posiciones, velocidades,
campos) se consideran medibles con precisión ilimitada.

% -------------------------------------------------------------------------
\subsection{La interacción electromagnética entre las cuatro fuerzas de la naturaleza}\label{sec:fuerzas}
% -------------------------------------------------------------------------
La mecánica dice cómo responde un sistema a una fuerza dada; la pregunta de fondo es de
dónde vienen las fuerzas. Todas las que conocemos se reducen a cuatro interacciones
fundamentales, que ordenamos por intensidad relativa, dando el valor~1 a la más intensa
(\cref{tab:fuerzas}).

\begin{table}[ht]
\centering
\caption{Las cuatro interacciones fundamentales. La intensidad es la relativa entre
dos protones, tomando como unidad la fuerza nuclear fuerte.}
\label{tab:fuerzas}
\small
\renewcommand{\arraystretch}{1.3}
\begin{tabular}{@{}lccl@{}}
\toprule
Interacción & Intensidad relativa & Alcance & Magnitud a la que se acopla\\
\midrule
Fuerte          & $1$          & corto ($\sim10^{-15}\,\mathrm{m}$) & carga de color\\
Electromagnética& $10^{-2}$    & largo                              & carga eléctrica\\
Débil           & $10^{-5}$    & muy corto                          & carga débil\\
Gravitatoria    & $10^{-38}$   & largo                              & masa-energía\\
\bottomrule
\end{tabular}
\end{table}

Conviene darse cuenta de la amplitud del dominio electromagnético. La fricción, la fuerza
normal que sostiene los cuerpos sobre el suelo, la que mantiene unidos los enlaces
moleculares o la que actúa en el choque de dos bolas de billar son, todas, fuerzas
electromagnéticas. La fuerza nuclear fuerte mantiene unidos los protones en el núcleo, pero
es de muy corto alcance: en un núcleo con muchos protones (como el uranio, con 92) el equilibrio es
frágil y el núcleo puede romperse. La interacción débil, responsable de ciertos procesos de
desintegración nuclear, es también de corto alcance y mucho menos intensa. La
electromagnética, asociada a la carga eléctrica, es de largo alcance y domina la escala
atómica en cuanto se supera la nuclear. La gravitatoria, asociada a la masa, es también de
largo alcance, pero extraordinariamente débil en comparación.

\begin{observacion}[Electricidad frente a gravedad]\label{obs:gravedad}
La repulsión eléctrica entre dos electrones es unas $4\cdot10^{42}$ veces mayor que su
atracción gravitatoria; el cociente $e^{2}/(4\pi\eo\,Gm_e^{2})$ no depende de la distancia, pues ambas
fuerzas decaen como $1/r^{2}$. Si la gravedad, y no la electricidad, mantuviese ligado al
electrón en el átomo de hidrógeno, el análogo del radio de Bohr,
$a_G=\hbar^{2}/(Gm_e^{2}m_p)\approx10^{29}\,\mathrm{m}$, sería mayor que el radio del
universo observable ($\sim4\cdot10^{26}\,\mathrm{m}$).
\end{observacion}

% -------------------------------------------------------------------------
\subsection{Breve historia: de la piedra de imán a la teoría de Maxwell}\label{sec:historia}
% -------------------------------------------------------------------------
Una exposición correcta del electromagnetismo pasa por conocer su historia, y el
planteamiento de este curso, que va de lo más sencillo a lo más complicado, es parecido a su
desarrollo histórico (\cref{fig:historia})~\cite{WHITTAKER,HEILBRON}.

% Figura: línea temporal (no a escala) del electromagnetismo hasta 1905.
\begin{figure}[ht]
\centering
\begin{tikzpicture}[x=1.52cm, y=1cm]
  \draw[line width=1.1pt, azulacero, -{Stealth[length=6pt]}] (-0.4,0) -- (9.5,0);
  \tikzset{ev/.style={font=\scriptsize, align=center, text width=2.55cm, text=black!85},
           anio/.style={font=\scriptsize\bfseries, text=azulprof}}
  \foreach \i/\an/\txt/\pos in {
      0/{s.\,VI a.\,C.}/{Tales de Mileto:\\el ámbar frotado}/1,
      1/{1600}/{Gilbert,\\\emph{De Magnete}}/-1,
      2/{1771}/{Cavendish:\\ley del cuadrado}/1,
      3/{1785}/{Coulomb:\\balanza de torsión}/-1,
      4/{1800}/{Volta:\\la pila}/1,
      5/{1820}/{Oersted y Ampère}/-1,
      6/{1831}/{Faraday:\\inducción y campo}/1,
      7/{1873}/{Maxwell:\\\emph{Tratado}}/-1,
      8/{1888}/{Hertz:\\ondas e.\,m.}/1,
      9/{1900--05}/{Planck,\\Einstein}/-1}{
    \fill[azulprof] (\i,0) circle (2.2pt);
    \ifnum\pos=1
      \draw[grispizarra!60, line width=0.4pt] (\i,0.1) -- (\i,0.55);
      \node[anio, anchor=south] at (\i,0.52) {\an};
      \node[ev, anchor=south] at (\i,0.95) {\txt};
    \else
      \draw[grispizarra!60, line width=0.4pt] (\i,-0.1) -- (\i,-0.55);
      \node[anio, anchor=north] at (\i,-0.52) {\an};
      \node[ev, anchor=north] at (\i,-0.95) {\txt};
    \fi}
  \node[font=\tiny, text=grispizarra, anchor=north east] at (9.45,-0.12) {no a escala};
\end{tikzpicture}
\caption{Hitos del electromagnetismo, de la Antigüedad al cambio de siglo. El eje
no está a escala: casi todo ocurre entre 1771 y 1888.}
\label{fig:historia}
\end{figure}


Las raíces de las palabras \emph{electricidad} y \emph{magnetismo} se remontan a la
Antigüedad. Tales de Mileto (siglo~VI a.\,C.) observó que el ámbar (en griego
$\ddot{\eta}\lambda\epsilon\kappa\tau\rho o\nu$, \emph{élektron}), al ser frotado, atraía
pajitas. El conocimiento del magnetismo nace de la observación de las fuerzas entre un mineral,
la «piedra de imán», y trozos de hierro; el origen de la palabra es oscuro. Lucrecio la
hace derivar de Magnesia, en Asia Menor, y Plinio refiere la leyenda del pastor Magnes, cuyo
cayado de hierro quedó prendido a una roca cercana. En 1600, William Gilbert separó con
cuidado los fenómenos eléctricos de los magnéticos en su \emph{De Magnete}.

El salto cuantitativo llega en el siglo~XVIII. En 1771 Henry Cavendish demostró, con un
experimento de esferas concéntricas (\cref{sec:coulomb-exp}), que la fuerza eléctrica
decae con el cuadrado de la distancia; sus resultados no se publicaron hasta que
Maxwell los editó en 1879~\cite{CAVENDISH-ELEC}. En 1785 Charles-Augustin de
Coulomb publicó con una balanza de torsión la ley que lleva su nombre~\cite{COULOMB}.

Electricidad y magnetismo fueron dos ciencias independientes hasta 1820, cuando Hans Christian
Oersted descubrió (casi por accidente) que una corriente eléctrica desvía una brújula: nacía el
electromagnetismo. Pocos meses después, André-Marie Ampère desarrolló una teoría matemática de
la fuerza entre elementos de corriente. Michael Faraday introdujo una idea nueva, distinta de
la acción a distancia de la mecánica newtoniana: el concepto de \emph{campo}, y en 1831
descubrió las corrientes de inducción que aparecen en conductores sometidos a un flujo magnético
variable, base de generadores, motores y transformadores. A partir de los experimentos de
Faraday, James Clerk Maxwell reunió en 1873, en su \emph{Tratado de electricidad y
magnetismo}~\cite{MAXWELL}, todo el formalismo matemático, incluidas las ecuaciones que hoy
llevan su nombre: un campo magnético variable crea un campo eléctrico sin necesidad de soporte
material, la corriente de desplazamiento y la teoría electromagnética de la luz. Su predicción de
las ondas electromagnéticas fue confirmada por Heinrich Hertz en 1888, y abrió el camino de las
telecomunicaciones (Marconi, 1896). Ni la teoría de los cuantos (Planck, 1900) ni la relatividad
especial (Einstein, 1905) han empañado el brillo de las ecuaciones de Maxwell.

% -------------------------------------------------------------------------
\subsection{De la acción a distancia al concepto de campo}\label{sec:campo-idea}
% -------------------------------------------------------------------------
La ley de Coulomb es un ejemplo de \emph{acción a distancia}: dos cargas separadas se
atraen o repelen sin que nada ocupe el espacio entre ellas. Esa idea es heredera de la gravitación
newtoniana. Sin embargo, esta interacción (como otras) se estudia con mucha más comodidad
como el encadenamiento de dos procesos: en el primero, un grupo de cargas, las \emph{fuentes},
perturban el espacio que las rodea, dotándolo de propiedades que no tenía antes; en el segundo,
otra carga sufre una fuerza neta por interactuar con ese campo previamente creado. El campo es,
pues, una propiedad de cada punto del espacio, que el teórico puede estudiar sin saber quién va
a ser la carga que lo sienta. Lo que hace a esta descripción tan fecunda es que relacionar
campos con fuentes es un problema mucho más rico que calcular fuerzas.

El plan del tema sigue una cadena lógica sencilla (\cref{fig:esquema}). La \emph{ley de
Coulomb} junto con el \emph{principio de superposición} determinan el campo eléctrico
$\vE(\vc{r})$. De la forma $1/R^{2}$ de la ley se deduce la \emph{divergencia} del campo
(ley de Gauss) y, de su carácter central, su \emph{rotacional} (nulo), de modo que $\vE$ deriva
de un potencial. Como el teorema de Helmholtz asegura que un campo queda determinado cuando se
conocen su divergencia y su rotacional (con condiciones de regularidad en el infinito), esas dos
ecuaciones, junto con las condiciones de contorno, contienen \emph{toda} la información de la
electrostática en el vacío.

% Figura: esquema lógico del tema (equivalente al esquema de los apuntes manuscritos).
\begin{figure}[ht]
\centering
\begin{tikzpicture}[
    caja/.style={draw=azulclaro, fill=azulfondo, line width=0.5pt, align=center,
                 inner sep=5pt, font=\small, minimum height=1.0cm},
    final/.style={draw=azulacero, fill=white, line width=0.8pt, align=center,
                  inner sep=5pt, font=\small, minimum height=1.0cm},
    fl/.style={->, line width=0.7pt, azulacero, shorten >=2pt, shorten <=2pt},
    nota/.style={font=\scriptsize, text=grispizarra, align=center}]
  \node[caja] (L) at (0,0) {Ley de Coulomb\\[-1pt]{\footnotesize (experimental)}};
  \node[caja] (S) at (0,-1.75) {Principio de\\superposición};
  \node[caja] (E) at (3.7,-0.87) {Campo eléctrico\\$\vE(\vc{r})$};
  \node[caja] (D) at (7.5,0.35) {$\nab\cdot\vE=\rho/\eo$\\[-1pt]{\footnotesize (ley de Gauss)}};
  \node[caja] (R) at (7.5,-2.1) {$\nab\times\vE=\vc{0}$\\[-1pt]{\footnotesize $\vE=-\nab\phi$}};
  \node[final] (H) at (11.5,-0.87) {Condiciones\\de contorno};
  \draw[fl] (L.east) -- (E.north west);
  \draw[fl] (S.east) -- (E.south west);
  \draw[fl] (E.east) -- (D.west);
  \draw[fl] (E.east) -- (R.west);
  \draw[fl] (D.east) -- (H.north west);
  \draw[fl] (R.east) -- (H.south west);
  \node[nota] at (7.5,-0.87) {teorema de\\Helmholtz};
\end{tikzpicture}
\caption{Estructura lógica del tema. De la ley de Coulomb y la superposición se deduce
el campo; su divergencia (ley de Gauss) y su rotacional (campo conservativo, potencial)
lo determinan por el teorema de Helmholtz, y, aplicadas a una frontera, dan las
condiciones de contorno.}
\label{fig:esquema}
\end{figure}


% -------------------------------------------------------------------------
\subsection{Notación y convenios del tema}\label{sec:notacion}
% -------------------------------------------------------------------------
Mantenemos la notación del Tema~0~\cite{TEMA0,NAVAS} y las expresiones del
formulario~\cite{FORMULARIO}. Los vectores llevan flecha, los versores sombrero, y con
$\vc{r}$ designamos el punto en el que se evalúa el campo (\emph{punto campo}) y con
$\vc{r}\,'$ el punto donde está la carga que lo crea (\emph{punto fuente}). Como en el
Tema~0, usamos $\sen$ para el seno.

\begin{table}[ht]
\centering
\caption{Notación específica del Tema~1.}
\label{tab:notacion1}
\small
\renewcommand{\arraystretch}{1.35}
\begin{tabular}{@{}p{4.4cm}p{9.4cm}@{}}
\toprule
Símbolo & Significado\\
\midrule
$q,\ \qp$ & Carga de prueba y carga fuente (en C).\\
$\vc{r},\ \vc{r}\,'$ & Posición del punto campo y del punto fuente.\\
$\vc{R}=\vc{r}-\vc{r}\,',\ R=\norm{\vc{R}},\ \uR=\vc{R}/R$ &
  Vector de posición relativa (de la fuente al punto campo), su módulo y su versor.\\
$\vE(\vc{r}),\ \vF$ & Campo eléctrico (V/m) y fuerza (N).\\
$\rho,\ \sigma,\ \lambda$ & Densidades volúmica (C/m$^{3}$), superficial (C/m$^{2}$) y lineal (C/m) de carga.\\
$\phi(\vc{r})$ & Potencial electrostático (V).\\
$\flujo=\oiint_{S}\vE\cdot\dS$ & Flujo eléctrico a través de una superficie cerrada.\\
$\dOm$ & Elemento de ángulo sólido (sr).\\
$\eo,\ \mu_{0},\ c$ & Permitividad y permeabilidad del vacío; velocidad de la luz.\\
$\delta(x),\ \delta^{3}(\vc{r})$ & Delta de Dirac en una y en tres dimensiones.\\
\bottomrule
\end{tabular}
\end{table}

\begin{notacion}[Convenios del tema]\label{not:convenios1}
Trabajamos en el vacío, con cargas en reposo (estática) y campos de clase $C^{1}$ salvo en
las superficies y curvas de carga, que se tratan como idealizaciones (\cref{sec:contorno}). Las
superficies cerradas se orientan con la normal exterior $\un$ y escribimos $\dS=\un\,\dsc$.
Las unidades son las del Sistema Internacional (SI).
\end{notacion}

% =========================================================================
\section{La ley de Coulomb}\label{sec:coulomb}
% =========================================================================
% -------------------------------------------------------------------------
\subsection{La carga eléctrica}\label{sec:carga}
% -------------------------------------------------------------------------
La materia está formada por cargas eléctricas positivas y negativas que, en la materia
ordinaria, se compensan: los cuerpos son neutros. Recoge este equilibrio el siguiente
ejemplo, que da una idea de la enorme intensidad de la fuerza eléctrica.

\begin{ejemplo}[¿Cuánta carga sobra si se rompe la neutralidad?]\label{ej:neutralidad}
Una persona de $70\,\mathrm{kg}$ contiene unos $2\cdot10^{28}$ electrones. Si cada una de
dos personas situadas a $1\,\mathrm{m}$ tuviese un 1\,\% más de electrones que de protones,
cada una llevaría una carga $Q\approx2\cdot10^{26}\times1{,}6\cdot10^{-19}\,\mathrm{C}\approx3\cdot10^{7}\,\mathrm{C}$,
y la fuerza de repulsión sería
\[
  F\approx\frac{1}{4\pi\eo}\frac{Q^{2}}{d^{2}}\approx9\cdot10^{9}\cdot\frac{(3\cdot10^{7})^{2}}{1^{2}}\,\mathrm{N}\sim10^{25}\,\mathrm{N},
\]
del orden del peso de toda la Tierra ($\sim6\cdot10^{25}\,\mathrm{N}$).
\end{ejemplo}

Se puede suministrar o extraer carga a un cuerpo neutro, por ejemplo frotando unos materiales
con otros (vidrio con seda): el vidrio queda cargado positivamente (con defecto de electrones) y la
seda negativamente (con exceso). Pero la carga total del sistema se \emph{conserva}: la carga no se
crea ni se destruye. Además, está \emph{cuantizada}: cualquier carga $q$ es un múltiplo entero
de la carga elemental, $e=1{,}6\cdot10^{-19}\,\mathrm{C}$ (carga del electrón, $-e$). Entre dos cuerpos cargados aparecen
fuerzas que dependen de la cantidad de carga y de cómo está distribuida sobre el volumen o la superficie
de cada cuerpo. Por eso es útil la idealización de \emph{carga puntual}: toda la carga concentrada
en un punto geométrico. Es una buena aproximación cuando la distancia entre los cuerpos es mucho mayor que
sus dimensiones.

% -------------------------------------------------------------------------
\subsection{Evidencia experimental: de la balanza de torsión a Cavendish}\label{sec:coulomb-exp}
% -------------------------------------------------------------------------
La primera investigación cuantitativa de la dependencia de esta fuerza con la distancia fue
realizada por Coulomb en 1785~\cite{COULOMB} con una \emph{balanza de torsión}
(\cref{fig:torsion}). Una aguja aislante, colgada de un hilo fino, lleva una esfera cargada en un extremo
y un contrapeso en el otro; otra esfera cargada, fija, repele a la primera. En el equilibrio, el par
de torsión del hilo, proporcional al ángulo girado, iguala al par de la fuerza electrostática: midiendo ángulos se miden
fuerzas. Es un experimento muy delicado, sujeto a pérdidas de carga y a fuerzas de inducción sobre las esferas, y de él
solo se obtiene con precisión modesta la ley $1/r^{2}$.

% Figura: balanza de torsión de Coulomb (esquema).
\begin{figure}[ht]
\centering
\tikzset{esf/.style={circle, inner sep=0pt, minimum size=9pt, ball color=azulacero!75, draw=azulprof, line width=0.4pt},
         esffija/.style={circle, inner sep=0pt, minimum size=9pt, ball color=acento!80, draw=acento!60!black, line width=0.4pt}}
\begin{minipage}[t]{0.56\textwidth}\centering
\begin{tikzpicture}[baseline=(current bounding box.north), x=1cm, y=1cm]
  % (a) vista en perspectiva: cilindro de vidrio, hilo, aguja y esferas
  \def\rx{1.55}\def\ry{0.42}
  \fill[azulclaro!30] (-\rx,0) -- (-\rx,3.2) arc (180:360:{\rx} and \ry) -- (\rx,0) arc (0:-180:{\rx} and \ry) -- cycle;
  \draw[azulprof, line width=0.8pt] (-\rx,0) -- (-\rx,3.2) (\rx,0) -- (\rx,3.2);
  \draw[azulprof, line width=0.8pt, fill=azulclaro!45] (0,3.2) ellipse ({\rx} and \ry);
  \draw[azulprof, line width=0.8pt] (-\rx,0) arc (180:360:{\rx} and \ry);
  \draw[azulprof, line width=0.5pt, densely dashed] (-\rx,0) arc (180:0:{\rx} and \ry);
  % escala graduada alrededor del cilindro (marcas en la mitad frontal)
  \draw[azulprof!70, line width=0.5pt] (-\rx,1.1) arc (180:360:{\rx} and \ry);
  \foreach \a in {190,200,...,350}{\draw[azulprof!70, line width=0.35pt] ({\rx*cos(\a)},{1.1+\ry*sin(\a)}) -- ({\rx*cos(\a)},{1.1+\ry*sin(\a)+0.1});}
  % cabezal micrométrico
  \draw[azulprof, line width=0.7pt, fill=azulclaro!70] (-0.45,3.52) rectangle (0.45,3.95);
  \draw[azulprof, line width=0.7pt, fill=azulclaro!40] (0,3.95) ellipse (0.45 and 0.13);
  \draw[acento, line width=0.9pt] (0,3.95) -- (0.3,3.99);
  \draw[azulprof, line width=0.7pt, fill=azulclaro!55] (0,3.52) ellipse (0.45 and 0.13);
  % hilo de torsión
  \draw[grispizarra, line width=0.6pt] (0,3.38) -- (0,1.45);
  % aguja (en perspectiva) con esfera móvil y contrapeso
  \draw[grispizarra, line width=1.3pt] (-1.15,1.28) -- (1.15,1.62);
  \fill[azulprof] (0,1.45) circle (1.3pt);
  \node[esf] at (1.15,1.62) {}; \node[esf] at (-1.15,1.28) {};
  % esfera fija sujeta desde la tapa
  \draw[grispizarra, line width=0.9pt] (1.02,3.05) -- (1.02,1.35);
  \node[esffija, minimum size=8pt] at (1.02,1.30) {};
  % rótulos
  \draw[grispizarra, line width=0.3pt] (0.5,3.8) -- (2.3,3.8); \node[font=\footnotesize, anchor=west] at (2.35,3.8) {cabezal micrométrico};
  \draw[grispizarra, line width=0.3pt] (0.02,2.6) -- (2.3,2.6); \node[font=\footnotesize, anchor=west] at (2.35,2.6) {hilo de torsión};
  \draw[grispizarra, line width=0.3pt] (1.25,1.7) -- (2.3,1.95); \node[font=\footnotesize, anchor=west] at (2.35,1.95) {esfera móvil};
  \draw[grispizarra, line width=0.3pt] (1.12,1.18) -- (2.3,1.05); \node[font=\footnotesize, anchor=west] at (2.35,1.05) {esfera fija};
  \draw[grispizarra, line width=0.3pt] (1.5,0.62) -- (2.3,0.35); \node[font=\footnotesize, anchor=west] at (2.35,0.3) {escala graduada};
  \node[font=\footnotesize, anchor=east] at (-1.75,1.1) {aguja};
  \draw[grispizarra, line width=0.3pt] (-1.7,1.15) -- (-1.25,1.27);
\end{tikzpicture}\\[2pt]{\small\scshape (a)}
\end{minipage}\hfill
\begin{minipage}[t]{0.40\textwidth}\centering
\begin{tikzpicture}[baseline=(current bounding box.north), x=1cm, y=1cm]
  % (b) vista superior: escala graduada y ángulo de torsión
  \def\Rg{1.65}
  \draw[azulprof, line width=0.8pt] (0,0) circle (\Rg);
  \foreach \a in {0,6,...,354}{\draw[azulprof!70, line width=0.3pt] (\a:\Rg) -- (\a:{\Rg-0.09});}
  \foreach \a in {0,30,...,330}{\draw[azulprof, line width=0.5pt] (\a:\Rg) -- (\a:{\Rg-0.17});}
  \draw[verdeteo, densely dashed, line width=0.5pt] (0,0) -- (0:1.3);
  \draw[grispizarra, line width=1.3pt] (200:1.2) -- (22:1.2);
  \fill[azulprof] (0,0) circle (1.4pt) node[below left, font=\scriptsize] {$O$};
  \node[esf] at (22:1.2) {}; \node[esf] at (200:1.2) {};
  \node[esffija] at (0:1.3) {};
  \draw[verdeteo, line width=0.6pt] (0:0.7) arc (0:22:0.7);
  \node[font=\small, text=verdeteo] at (11:0.95) {$\theta$};
  \node[font=\scriptsize, text=acento, anchor=west] at (1.85,-0.35) {fija};
  \node[font=\scriptsize, text=azulprof, anchor=south west] at (1.35,1.25) {móvil};
\end{tikzpicture}\\[2pt]{\small\scshape (b)}
\end{minipage}
\caption{Balanza de torsión de Coulomb (esquema). (a)~Una aguja aislante suspendida de un hilo fino lleva una esfera móvil en un extremo y un contrapeso en el otro; una segunda esfera
fija se introduce desde la tapa. Un cabezal micrométrico permite girar el extremo superior del hilo, y una escala graduada rodea el cilindro de vidrio. (b)~Vista desde arriba: al cargar ambas esferas, la repulsión desvía la aguja un
ángulo $\theta$ medido sobre la escala; en el equilibrio, el par recuperador del hilo, proporcional a $\theta$, iguala al de la fuerza electrostática. Girando el cabezal se varía la distancia entre las esferas y se comprueba cómo escala la fuerza con ella.}
\label{fig:torsion}
\end{figure}


Mucho antes que Coulomb, y con mayor precisión, Cavendish había contrastado el exponente
por un camino indirecto (1771--73), que merece la pena entender porque es una de las pruebas más
limpias de la ley (\cref{fig:cavendish}). El argumento descansa en un resultado geométrico: una
cáscara esférica cargada uniformemente ejerce una fuerza nula en su interior \emph{si y solo si} la fuerza
decae como $1/r^{2}$. Este hecho aparece ya, sin demostración, en Franklin y Priestley, y es el análogo
eléctrico del teorema de Newton para la gravedad.

\begin{proposicion}[Argumento de Cavendish]\label{prop:cavendish}
Sea una cáscara esférica con densidad superficial de carga uniforme $\sigma$ y supongamos que
la fuerza entre cargas puntuales es central, repulsiva (para cargas iguales) y de módulo
$K\,q_1q_2/R^{\,n}$. Una carga $q$ en cualquier punto $P$ interior a la cáscara sufre fuerza neta nula
si y solo si $n=2$.
\end{proposicion}
\begin{proof}
Tracemos desde $P$ un cono estrecho de ángulo sólido $\dOm$ y su prolongación opuesta
(\cref{fig:cavendish}b). Los dos conos recortan en la cáscara sendos elementos
$\dif S_1$ y $\dif S_2$ a distancias $r_1$ y $r_2$ de $P$. El ángulo $\alpha$ entre la cuerda y la
normal a la esfera (el radio) es el mismo en ambos extremos, porque el triángulo formado por los dos
radios y la cuerda es isósceles. Por tanto, $\dif S_i=r_i^{2}\,\dOm/\cos\alpha$ y las fuerzas
sobre $q$ actúan a lo largo de la cuerda y en sentidos opuestos,
\[
  \dif F_i=K\,q\,\frac{\sigma\,\dif S_i}{r_i^{\,n}}=\frac{K q\sigma\,\dOm}{\cos\alpha}\;r_i^{\,2-n},
  \qquad
  \dif F_{\mathrm{neta}}=\frac{K q\sigma\,\dOm}{\cos\alpha}\left(r_2^{\,2-n}-r_1^{\,2-n}\right).
\]
Para $n=2$ esta cantidad se anula para cada par de elementos opuestos, y la fuerza total, suma
(integral) sobre todos ellos, es nula. Si $n\neq2$ y $P$ no es el centro, $r_1\neq r_2$ y
$r_2^{\,2-n}\neq r_1^{\,2-n}$, de modo que la contribución de cada pareja no es nula y la
fuerza total es, en general, distinta de cero.
\end{proof}

% Figura: experimento de Cavendish y argumento de la cáscara esférica.
\begin{figure}[ht]
\centering
\tikzset{
  parche/.style={draw=azulprof, line width=0.6pt, fill=azulclaro!70, fill opacity=0.9},
  parche2/.style={draw=acento, line width=0.6pt, fill=acento!25, fill opacity=0.9},
  vn/.style={-{Stealth[length=5pt,width=4.5pt]}, line width=1.1pt}}
\begin{minipage}[t]{0.50\textwidth}\centering
\begin{tikzpicture}[baseline=(current bounding box.north), x=0.82cm, y=0.82cm]
  % (a) cáscara de dos hemisferios con esfera interior unida por un hilo conductor
  \def\Re{1.9}\def\Ri{0.6}
  \shade[ball color=azulclaro!55, opacity=0.35] (0,0) circle (\Re);
  \draw[azulprof, line width=0.8pt] (0,0) circle (\Re);
  \draw[azulprof, line width=0.8pt] (-\Re,0) arc (180:360:{\Re} and 0.45);
  \draw[azulprof, line width=0.5pt, densely dashed] (-\Re,0) arc (180:0:{\Re} and 0.45);
  \shade[ball color=azulacero!70, opacity=0.95] (0,0) circle (\Ri);
  \draw[azulprof, line width=0.7pt] (0,0) circle (\Ri);
  % hilo conductor vertical con collar aislante en la cáscara y terminal exterior
  \draw[grispizarra, line width=1.0pt] (0,\Ri) -- (0,2.55);
  \draw[azulprof, line width=0.6pt, fill=white] (0,\Re) ellipse (0.13 and 0.05);
  \fill[acento] (0,2.55) circle (2.0pt);
  % rótulos con líneas de llamada rectas
  \node[font=\footnotesize, anchor=west] at (0.35,2.55) {hilo conductor};
  \draw[grispizarra, line width=0.3pt] (-2.65,0.45) -- (-1.85,0.45); \node[font=\footnotesize, anchor=east] at (-2.7,0.45) {hemisferios};
  \draw[grispizarra, line width=0.3pt] (2.7,-0.9) -- (1.5,-0.9); \node[font=\footnotesize, anchor=west] at (2.75,-0.9) {cáscara};
  \draw[grispizarra, line width=0.3pt] (2.7,0.0) -- (0.62,0.0); \node[font=\footnotesize, anchor=west] at (2.75,0.0) {esfera interior};
  \draw[verdeteo, <->, line width=0.5pt] (0,0) -- (-35:\Ri) node[pos=0.7, below=1pt, font=\scriptsize, text=verdeteo] {$a$};
\end{tikzpicture}\\[2pt]{\small\scshape (a)}
\end{minipage}\hfill
\begin{minipage}[t]{0.48\textwidth}\centering
\begin{tikzpicture}[baseline=(current bounding box.north), x=1.05cm, y=1.05cm]
  % (b) cono desde un punto interior P que recorta dos elementos opuestos de la cáscara
  \def\Rr{2.0}
  \draw[azulprof, line width=0.9pt, name path=circ, fill=azulfondo, fill opacity=0.5] (0,0) circle (\Rr);
  \coordinate (O) at (0,0); \fill[azulprof] (O) circle (1.2pt); \node[font=\scriptsize, anchor=north east] at (O) {$C$};
  \coordinate (P) at (0.7,-0.25);
  \def\dd{24}\def\ap{9}
  \foreach \dir/\nm in {\dd-\ap/a,\dd+\ap/b,\dd/m}{
    \path[name path=ray\nm] (P) -- ++({\dir}:5);
    \path[name path=rev\nm] (P) -- ++({\dir+180}:5);
    \path[name intersections={of=circ and ray\nm, by=A\nm}];
    \path[name intersections={of=circ and rev\nm, by=B\nm}];}
  \fill[azulclaro, opacity=0.40] (P) -- (Aa) -- (Ab) -- cycle;
  \fill[azulclaro, opacity=0.40] (P) -- (Ba) -- (Bb) -- cycle;
  \draw[azulacero, line width=0.45pt] (Aa) -- (Ba) (Ab) -- (Bb);
  \draw[azulprof, line width=0.8pt] (Bm) -- (Am);
  \draw[acento, line width=2.8pt] (Aa) -- (Ab);
  \draw[acento, line width=2.8pt] (Ba) -- (Bb);
  \draw[verdeteo, line width=0.6pt, densely dashed] (0,0) -- (Am) (0,0) -- (Bm);
  \fill[azulprof] (P) circle (1.5pt) node[below=2pt, font=\small] {$P$};
  \node[font=\scriptsize, text=azulprof, anchor=south east] at ($(P)!0.55!(Am)$) {$r_2$};
  \node[font=\scriptsize, text=azulprof, anchor=north west] at ($(P)!0.55!(Bm)$) {$r_1$};
  \node[font=\small, text=acento, anchor=south west] at ($(Am)+(0.1,0.12)$) {$\dif S_2$};
  \node[font=\small, text=acento, anchor=north east] at ($(Bm)+(-0.12,-0.2)$) {$\dif S_1$};
  % ángulo alpha entre la cuerda y el radio en cada extremo (arcos dibujados a mano)
  \draw[verdeteo, line width=0.55pt] let \p1=(Am), \n1={atan2(\y1,\x1)+180} in
     (Am) ++(\n1:0.45) arc[start angle=\n1, end angle=\dd+180, radius=0.45];
  \draw let \p1=(Am), \n1={atan2(\y1,\x1)+180} in
     node[font=\scriptsize, text=verdeteo] at ($(Am)+({(\n1+\dd+180)/2}:0.78)$) {$\alpha$};
  \draw[verdeteo, line width=0.55pt] let \p1=(Bm), \n1={atan2(\y1,\x1)+180} in
     (Bm) ++(\n1:0.45) arc[start angle=\n1, end angle=\dd, radius=0.45];
  \draw let \p1=(Bm), \n1={atan2(\y1,\x1)+180} in
     node[font=\scriptsize, text=verdeteo] at ($(Bm)+({(\n1+\dd)/2}:0.78)$) {$\alpha$};
\end{tikzpicture}\\[2pt]{\small\scshape (b)}
\end{minipage}
\caption{(a)~Experimento de Cavendish: una esfera conductora interior se conecta, mediante un hilo, a una cáscara formada por dos hemisferios, y se cargan; al retirar el hilo y separar los hemisferios,
la esfera interior resulta descargada. (b)~Argumento geométrico: un cono de ángulo sólido $\dOm$ con vértice en un punto interior $P$ recorta dos elementos, $\dif S_1$ y $\dif S_2$,
a distancias $r_1$ y $r_2$ de $P$; ambos forman con la cuerda el mismo ángulo $\alpha$ (el triángulo formado por los dos radios y la cuerda es isósceles).}
\label{fig:cavendish}
\end{figure}


Cavendish consiguió la configuración de la \cref{fig:cavendish}a: una esfera conductora dentro de
una cáscara, unidas por un hilo, de modo que ambas forman un único conductor. Se cargan, se retira
el hilo y se abre la cáscara. Si la ley fuese exactamente $1/r^{2}$, la carga residiría
toda en la cáscara y la esfera interior quedaría descargada, como en efecto observó: con electrómetros de
la época concluyó que el exponente era $2\pm0{,}02$. Los experimentos posteriores, siguiendo la misma idea, han
reducido la incertidumbre en muchos órdenes de magnitud. Si escribimos el exponente como $2+\varepsilon$:
\begin{center}\small
\renewcommand{\arraystretch}{1.25}
\begin{tabular}{@{}lcc@{}}
\toprule
Experimento & Año & $\lvert\varepsilon\rvert<$\\
\midrule
Cavendish~\cite{CAVENDISH-ELEC} & 1773 & $0{,}02$\\
Maxwell~\cite{MAXWELL} & 1873 & $5\cdot10^{-5}$\\
Plimpton y Lawton~\cite{PLIMPTON} & 1936 & $2\cdot10^{-9}$\\
Williams, Faller y Hill~\cite{WILLIAMS} & 1971 & $\sim3\cdot10^{-16}$\\
\bottomrule
\end{tabular}
\end{center}
El exponente 2 es, así, una de las constantes mejor comprobadas de la física~\cite{FEYNMAN}; sin embargo, la ley de Coulomb
sigue siendo una ley \emph{experimental}, que no admite demostración desde otros principios dentro de
la electrostática.

% -------------------------------------------------------------------------
\subsection{Enunciado de la ley}\label{sec:coulomb-ley}
% -------------------------------------------------------------------------
Elegido un sistema de coordenadas y dos cargas puntuales $q_1$ y $q_2$ en reposo, en el vacío, situadas
en $\vc{r}_1$ y $\vc{r}_2$ (\cref{fig:coulomb}a), la ley de Coulomb afirma que la fuerza es proporcional al
producto de las cargas, inversamente proporcional al cuadrado de la distancia que las separa, está dirigida
a lo largo de la recta que las une, y es atractiva si las cargas son de signos opuestos y repulsiva si
tienen el mismo signo. Matemáticamente:

\begin{definicion}[Ley de Coulomb]\label{def:coulomb}
La fuerza que la carga puntual $q_1$ ejerce sobre la carga puntual $q_2$, en reposo y en el vacío, es
\begin{equation}\label{eq:coulomb}
  \boxed{\;\vc{F}_{12}=K\,\frac{q_1q_2}{R_{12}^{2}}\,\uh{R_{12}}
  =\frac{1}{4\pi\eo}\,\frac{q_1q_2}{\norm{\vc{R}_{12}}^{3}}\,\vc{R}_{12},\qquad
  \vc{R}_{12}=\vc{r}_2-\vc{r}_1\;}
\end{equation}
con $\uh{R_{12}}=\vc{R}_{12}/R_{12}$ y $R_{12}=\norm{\vc{R}_{12}}$. Si $q_1q_2>0$ la fuerza es
repulsiva y, si $q_1q_2<0$, atractiva.
\end{definicion}

% Figura: ley de Coulomb (geometría) y principio de superposición.
\begin{figure}[ht]
\centering
\tdplotsetmaincoords{72}{125}
\begin{minipage}[t]{0.49\textwidth}\centering
\begin{tikzpicture}[tdplot_main_coords, scale=1.0, baseline=(current bounding box.north)]
  \draw[eje] (0,0,0) -- (3.2,0,0) node[below left] {\small $x$};
  \draw[eje] (0,0,0) -- (0,3.4,0) node[right] {\small $y$};
  \draw[eje] (0,0,0) -- (0,0,3.2) node[above] {\small $z$};
  \coordinate (Q1) at (2.6,0.4,0.6);
  \coordinate (Q2) at (0.4,2.7,2.3);
  \draw[->, azulclaro!80!azulacero, line width=0.7pt] (0,0,0) -- (Q1) node[pos=0.5, below, font=\small, text=azulprof] {$\vc{r}_1$};
  \draw[->, azulclaro!80!azulacero, line width=0.7pt] (0,0,0) -- (Q2) node[pos=0.45, left, font=\small, text=azulprof] {$\vc{r}_2$};
  \draw[vect, azulprof, line width=1.2pt] (Q1) -- (Q2) node[pos=0.62, above left, font=\small, text=azulprof] {$\vc{R}_{12}$};
  \draw[vectrojo] (Q2) -- ($(Q2)+(-0.45,0.7,0.65)$) node[above, font=\small, text=acento] {$\vc{F}_{12}$};
  \draw[vectrojo] (Q1) -- ($(Q1)+(0.45,-0.7,-0.65)$) node[left, font=\small, text=acento] {$\vc{F}_{21}$};
  \node[cargapos] at (Q1) {$+$};
  \node[cargapos] at (Q2) {$+$};
  \node[font=\small, anchor=south west] at ($(Q1)+(0.1,0.1,0.1)$) {$q_1$};
  \node[font=\small, anchor=south east] at ($(Q2)+(-0.1,0,0.15)$) {$q_2$};
\end{tikzpicture}\\[2pt]{\small\scshape (a)}
\end{minipage}\hfill
\begin{minipage}[t]{0.49\textwidth}\centering
\begin{tikzpicture}[tdplot_main_coords, scale=0.95, baseline=(current bounding box.north)]
  \draw[eje] (0,0,0) -- (3.0,0,0) node[below left] {\small $x$};
  \draw[eje] (0,0,0) -- (0,3.2,0) node[right] {\small $y$};
  \draw[eje] (0,0,0) -- (0,0,3.0) node[above] {\small $z$};
  \coordinate (P) at (1.3,1.6,1.1);
  \foreach \n/\a/\b/\c/\l in {1/2.9/0.0/0.3/{q_1}, 2/2.5/1.5/2.4/{q_2}, 3/0.2/3.1/1.8/{q_3}}{
     \coordinate (S\n) at (\a,\b,\c);
     \draw[vect, azulclaro!90!black, line width=0.7pt] (S\n) -- (P);
     \node[cargapos] at (S\n) {$+$};
     \node[font=\small, anchor=south east] at ($(S\n)+(-0.1,0,0.1)$) {$\l$};}
  \draw[vectrojo, line width=1.3pt] (P) -- ($(P)+(-0.9,0.6,0.45)$) node[above left, font=\small, text=acento] {$\vc{F}_{q}$};
  \node[cargapru] at (P) {$q$};
  \node[font=\scriptsize, text=azulprof, anchor=north] at ($(S1)!0.5!(P)$) {$\vc{R}_1$};
  \node[font=\scriptsize, text=azulprof, anchor=east] at ($(S2)!0.5!(P)$) {$\vc{R}_2$};
  \node[font=\scriptsize, text=azulprof, anchor=west] at ($(S3)!0.45!(P)$) {$\vc{R}_3$};
\end{tikzpicture}\\[2pt]{\small\scshape (b)}
\end{minipage}
\caption{(a)~Ley de Coulomb para dos cargas puntuales del mismo signo:
$\vc{R}_{12}=\vc{r}_2-\vc{r}_1$ apunta de $q_1$ a $q_2$, y la fuerza
$\vc{F}_{12}$ sobre $q_2$ es paralela a $\vc{R}_{12}$ (repulsión); la fuerza sobre $q_1$ es la
opuesta, $\vc{F}_{21}=-\vc{F}_{12}$. (b)~Superposición: la fuerza sobre la carga $q$ es la suma
vectorial de las ejercidas por cada $q_i$, con $\vc{R}_i$ el vector que une $q_i$ con $q$.}
\label{fig:coulomb}
\end{figure}


\begin{observacion}[Sobre la ley de Coulomb]\label{obs:coulomb}
\begin{enumerate}[label=(\alph*)]
  \item \emph{Acción y reacción.} Al intercambiar las cargas, $\vc{R}_{12}=-\vc{R}_{21}$ y por tanto
        $\vc{F}_{21}=-\vc{F}_{12}$: se cumple la tercera ley de Newton.
  \item \emph{La constante.} En el SI, $K=8{,}988\cdot10^{9}\,\mathrm{N\,m^{2}/C^{2}}$ y la unidad de carga,
        el culombio, se define a partir de la unidad de corriente (amperio): es la carga que atraviesa en
        un segundo la sección de un conductor por el que circula $1\,\mathrm{A}$. Se escribe habitualmente
        \[
          K=\frac{1}{4\pi\eo},\qquad
          \eo=\frac{1}{4\pi K}=8{,}85\cdot10^{-12}\,\frac{\mathrm{C^{2}}}{\mathrm{N\,m^{2}}}\;\Bigl(=\frac{\mathrm{F}}{\mathrm{m}}\Bigr),
        \]
        la \emph{permitividad del vacío}, muy próxima a la del aire seco, que es el medio donde Coulomb
        experimentó. Más adelante, en las ecuaciones de Maxwell, aparecerá la relación
        $c^{2}=1/(\eo\mu_0)$ con $c\simeq2{,}998\cdot10^{8}\,\mathrm{m/s}$ la velocidad de la luz y
        $\mu_0\approx4\pi\cdot10^{-7}\,\mathrm{H/m}$ la permeabilidad magnética del vacío~\cite{CODATA}.
  \item \emph{Alcance de validez.} La ley es experimental y su exponente está comprobado con enorme precisión
        (\cref{sec:coulomb-exp}). Aunque se estableció en condiciones macroscópicas, se aplica hasta distancias
        de $10^{-15}\,\mathrm{m}$.
  \item \emph{Tres restricciones.} (i)~\emph{Cargas puntuales}: es una idealización; las cargas deben ser pequeñas
        frente a su separación. (ii)~\emph{Cargas en reposo}: si se mueven aparece además una interacción
        magnética; el estudio de las cargas en reposo es la \emph{electrostática}. Para que se mantengan en
        reposo han de aparecer fuerzas mecánicas tales que $\sum\vF=\vc{0}$. (iii)~\emph{Cargas en el vacío}: en
        presencia de materiales aparecerían fuerzas adicionales (en los conductores, por redistribución de
        sus cargas por influencia; en los dieléctricos, por su polarización).
\end{enumerate}
\end{observacion}

% -------------------------------------------------------------------------
\subsection{Principio de superposición}\label{sec:superposicion}
% -------------------------------------------------------------------------
La fuerza entre dos cargas no cambia por la presencia de una tercera. Este hecho, resultado de la observación
experimental, es la base del \emph{principio de superposición}.

\begin{definicion}[Principio de superposición]\label{def:superposicion}
La fuerza que un sistema de cargas puntuales $q_1,\dots,q_N$ ejerce sobre una carga puntual $q$ es la suma
vectorial de las fuerzas que cada una de ellas ejerce por separado:
\begin{equation}\label{eq:superposicion}
  \vc{F}_q=\sum_{i=1}^{N}\vc{F}_{iq}=\frac{q}{4\pi\eo}\sum_{i=1}^{N}\frac{q_i}{R_i^{3}}\,\vc{R}_i
  =\frac{q}{4\pi\eo}\sum_{i=1}^{N}\frac{q_i}{R_i^{2}}\,\uh{R_i},
\end{equation}
donde $\vc{R}_i$ es el vector que une cada carga $q_i$ con la carga $q$ (\cref{fig:coulomb}b).
\end{definicion}

\begin{observacion}[Validez]\label{obs:superposicion}
La superposición se ha comprobado válida para campos tan intensos como $10^{21}\,\mathrm{V/m}$ y distancias
tan pequeñas como $10^{-10}\,\mathrm{m}$. Expresa la \emph{linealidad} del medio: vale en el vacío, pero
no necesariamente en todos los materiales. Es, junto con la ley de Coulomb, el segundo dato de partida que usaremos:
con \cref{eq:coulomb,eq:superposicion} basta para deducir el resto de la electrostática en el vacío.
\end{observacion}

% =========================================================================
\section{El campo eléctrico}\label{sec:campo}
% =========================================================================
La ley de Coulomb dice qué fuerza sufre una carga, pero no cómo repartir el trabajo entre «quien crea» y
«quien siente». Para separar ambos papeles, llamamos \emph{carga fuente} $\qp$ a la que crea la perturbación,
situada en el \emph{punto fuente} $\vc{r}\,'$, y \emph{carga de prueba} $q$ (que tomamos positiva) a la que
experimenta la fuerza, situada en el \emph{punto campo} $\vc{r}$. La carga de prueba no se incluye entre las
fuentes: ninguna carga ejerce fuerza sobre sí misma (\cref{fig:campo}a).

% -------------------------------------------------------------------------
\subsection{Definición del campo eléctrico}\label{sec:def-campo}
% -------------------------------------------------------------------------
Dividiendo la fuerza de Coulomb por la carga de prueba, desaparece toda dependencia de $q$ y queda algo asociado
únicamente a las fuentes.

\begin{definicion}[Campo eléctrico]\label{def:campo}
El \emph{campo eléctrico} en el punto $\vc{r}$ es la fuerza por unidad de carga que experimentaría una carga de prueba
en ese punto:
\begin{equation}\label{eq:campo-def}
  \vE(\vc{r})=\lim_{q\to0}\frac{\vF_q}{q}.
\end{equation}
Para una carga fuente puntual $\qp$ situada en $\vc{r}\,'$,
\begin{equation}\label{eq:campo-puntual}
  \boxed{\;\vE(\vc{r})=\frac{1}{4\pi\eo}\,\frac{\qp}{R^{2}}\,\uR
  =\frac{1}{4\pi\eo}\,\frac{\qp}{\norm{\vc{R}}^{3}}\,\vc{R},\qquad\vc{R}=\vc{r}-\vc{r}\,'\;}
\end{equation}
Su unidad en el SI es el N/C, aunque se usa con más frecuencia el V/m (\cref{sec:potencial-signif}).
\end{definicion}

\begin{observacion}[Por qué el límite $q\to0$]\label{obs:limite}
Con $\vE=\vF_q/q$ bastaría para una carga fuente «rígida», pero si la fuente es, por ejemplo, una esfera
conductora, acercar una carga de prueba apreciable redistribuye la carga $\qp$ sobre la esfera y el campo que
se quiere medir cambia, haciéndose dependiente de la propia carga de prueba. El límite $q\to0$ corrige esto: la
carga de prueba debe ser tan pequeña que no perturbe la distribución que crea el campo. Es una idealización, pues
existe una carga mínima, la del electrón.
\end{observacion}

Con esta definición, el campo es una función $\vE:\R^{3}\to\R^{3}$ (un \emph{campo vectorial}) y la fuerza
sobre cualquier carga $q$ colocada en $\vc{r}$ se obtiene como $\vF_q=q\,\vE(\vc{r})$.

% Figura: carga fuente, punto campo y distribuciones continuas de carga.
\begin{figure}[ht]
\centering
\begin{minipage}[t]{0.40\textwidth}\centering
\tdplotsetmaincoords{72}{125}
\begin{tikzpicture}[tdplot_main_coords, scale=0.95, baseline=(current bounding box.north)]
  \draw[eje] (0,0,0) -- (3.0,0,0) node[below left] {\small $x$};
  \draw[eje] (0,0,0) -- (0,3.2,0) node[right] {\small $y$};
  \draw[eje] (0,0,0) -- (0,0,3.0) node[above] {\small $z$};
  \coordinate (S) at (2.4,0.3,0.4);
  \coordinate (P) at (0.4,2.6,2.2);
  \draw[->, azulclaro!80!azulacero, line width=0.7pt] (0,0,0) -- (S) node[pos=0.5, below, font=\small, text=azulprof] {$\vc{r}\,'$};
  \draw[->, azulclaro!80!azulacero, line width=0.7pt] (0,0,0) -- (P) node[pos=0.45, left, font=\small, text=azulprof] {$\vc{r}$};
  \draw[vect, azulprof, line width=1.2pt] (S) -- (P) node[pos=0.55, right=2pt, font=\small, text=azulprof] {$\vc{R}$};
  \draw[vectrojo, line width=1.2pt] (P) -- ($(P)+(-0.45,0.7,0.65)$) node[above, font=\small, text=acento] {$\vE(\vc{r})$};
  \node[cargapos] at (S) {$+$};
  \node[cargapru] at (P) {$q$};
  \node[font=\small, anchor=south west] at ($(S)+(0.1,0.1,0.1)$) {$\qp$};
\end{tikzpicture}\\[2pt]{\small\scshape (a)}
\end{minipage}\hfill
\begin{minipage}[t]{0.56\textwidth}\centering
\begin{tikzpicture}[baseline=(current bounding box.north), x=1cm, y=1cm]
  % (b) volumen
  \begin{scope}[shift={(0,0)}]
    \draw[azulprof, fill=azulfondo, line width=0.6pt] (0,0.2) .. controls (-0.2,1.2) and (0.4,1.9) .. (1.1,1.8) .. controls (1.9,1.9) and (2.2,0.9) .. (1.7,0.3) .. controls (1.2,-0.3) and (0.4,-0.2) .. cycle;
    \fill[azulacero] (1.0,1.0) circle (1.3pt);
    \draw[azulacero, fill=azulclaro!60, line width=0.4pt] (0.95,0.98) rectangle (1.12,1.13);
    \node[font=\scriptsize, anchor=south east] at (0.95,1.1) {$\dif V\,'$};
    \node[font=\scriptsize, text=azulprof] at (0.4,0.35) {$V'$};
    \node[font=\scriptsize, anchor=north] at (1.0,-0.5) {$\dif q'=\rho\,\dVp$};
  \end{scope}
  % (c) superficie
  \begin{scope}[shift={(2.8,0)}]
    \draw[azulprof, fill=azulfondo, line width=0.6pt] (0,0.1) -- (0.3,1.5) .. controls (1.0,1.9) and (1.5,1.4) .. (2.0,1.7) -- (1.8,0.4) .. controls (1.2,0.0) and (0.8,0.4) .. cycle;
    \draw[azulacero, fill=azulclaro!60, line width=0.4pt] (0.9,0.8) -- (1.15,0.85) -- (1.1,1.1) -- (0.85,1.05) -- cycle;
    \node[font=\scriptsize, anchor=south east] at (0.85,1.1) {$\dSp$};
    \node[font=\scriptsize, text=azulprof] at (0.4,0.45) {$S'$};
    \node[font=\scriptsize, anchor=north] at (1.0,-0.5) {$\dif q'=\sigma\,\dSp$};
  \end{scope}
  % (d) línea
  \begin{scope}[shift={(5.6,0)}]
    \draw[azulprof, line width=1.1pt] (0,0.3) .. controls (0.5,1.5) and (1.0,-0.4) .. (1.6,1.0) .. controls (1.9,1.5) and (2.0,1.6) .. (2.1,1.7);
    \draw[azulacero, line width=2.8pt] (0.78,0.64) -- (0.99,0.62);
    \node[font=\scriptsize, anchor=south] at (0.85,0.95) {$\dlp$};
    \node[font=\scriptsize, text=azulprof] at (0.05,0.7) {$l'$};
    \node[font=\scriptsize, anchor=north] at (1.0,-0.5) {$\dif q'=\lambda\,\dlp$};
  \end{scope}
\end{tikzpicture}\\[2pt]{\small\scshape (b)}
\end{minipage}
\caption{(a)~Carga fuente $\qp$ en $\vc{r}\,'$ y carga de prueba $q$ en el punto campo $\vc{r}$;
$\vc{R}=\vc{r}-\vc{r}\,'$ va de la fuente al punto campo, y el campo $\vE(\vc{r})$ de una
fuente positiva apunta en su dirección. (b)~Distribuciones continuas: de volumen $V'$
($\dif q'=\rho\,\dVp$), de superficie $S'$ ($\dif q'=\sigma\,\dSp$) y de línea $l'$ ($\dif q'=\lambda\,\dlp$);
cada elemento actúa como una carga puntual $\dif q'$ en su propia posición $\vc{r}\,'$.}
\label{fig:campo}
\end{figure}


% -------------------------------------------------------------------------
\subsection{Campo de una distribución discreta de cargas}\label{sec:campo-discreto}
% -------------------------------------------------------------------------
\begin{proposicion}[Campo de un sistema de cargas puntuales]\label{prop:campo-discreto}
El campo eléctrico creado en $\vc{r}$ por las cargas puntuales $\qp_1,\dots,\qp_N$ situadas en
$\vc{r}\,'_1,\dots,\vc{r}\,'_N$ es
\begin{equation}\label{eq:campo-discreto}
  \vE(\vc{r})=\frac{1}{4\pi\eo}\sum_{i=1}^{N}\frac{\qp_i}{R_i^{2}}\,\uh{R_i}
  =\frac{1}{4\pi\eo}\sum_{i=1}^{N}\frac{\qp_i}{\norm{\vc{R}_i}^{3}}\,\vc{R}_i,\qquad
  \vc{R}_i=\vc{r}-\vc{r}\,'_i .
\end{equation}
\end{proposicion}
\begin{proof}
Por el principio de superposición (\cref{eq:superposicion}), la fuerza sobre una carga de prueba $q$ es la
suma de las fuerzas de cada $\qp_i$. Dividiendo por $q$ y tomando el límite $q\to0$ se obtiene
\cref{eq:campo-discreto}.
\end{proof}

% -------------------------------------------------------------------------
\subsection{Distribuciones continuas de carga}\label{sec:campo-continuo}
% -------------------------------------------------------------------------
Cuando las cargas están muy próximas entre sí en comparación con las demás distancias de interés, resulta conveniente
tratarlas como un continuo y definir densidades de carga.

\begin{definicion}[Densidades de carga]\label{def:densidades}
Sea $\dif q'$ la carga contenida en un elemento de volumen $\dVp$ (respectivamente, de superficie $\dSp$ o de
línea $\dlp$) situado en $\vc{r}\,'$. Se definen las densidades \emph{volúmica}, \emph{superficial} y \emph{lineal}:
\begin{align}
  \rho(\vc{r}\,')&=\frac{\dif q'}{\dVp}, & Q&=\int_{V'}\rho(\vc{r}\,')\,\dVp,\label{eq:rho}\\
  \sigma(\vc{r}\,')&=\frac{\dif q'}{\dSp}, & Q&=\int_{S'}\sigma(\vc{r}\,')\,\dSp,\label{eq:sigma}\\
  \lambda(\vc{r}\,')&=\frac{\dif q'}{\dlp}, & Q&=\int_{l'}\lambda(\vc{r}\,')\,\dlp.\label{eq:lambda}
\end{align}
\end{definicion}

\begin{observacion}[Elemento macroscópico]\label{obs:macro}
En \cref{eq:rho}, $\dVp$ debe ser un volumen \emph{pequeño a escala macroscópica}, de modo que permita seguir las
variaciones de la distribución y tratar $\dif q'$ como una carga puntual, pero \emph{grande a escala microscópica},
de modo que contenga muchos átomos o moléculas. Si $\dVp$ fuese menor que un átomo, en casi todas las posiciones no
habría carga y $\rho$ sería nula, salvo cuando $\dVp$ incluyese un electrón, y esto produciría fluctuaciones enormes en
$\rho$ que lo harían inservible como concepto. Así, $\rho$ representa un valor \emph{medio} que varía de forma suave
y que podemos considerar una función continua de la posición.
\end{observacion}

\begin{proposicion}[Campo de una distribución continua]\label{prop:campo-continuo}
El campo eléctrico creado en $\vc{r}$ por una distribución continua de carga es
\begin{align}
  \vE(\vc{r})&=\frac{1}{4\pi\eo}\int_{V'}\frac{\dif q'}{R^{2}}\,\uR
  =\frac{1}{4\pi\eo}\int_{V'}\frac{\rho(\vc{r}\,')}{\norm{\vc{R}}^{3}}\,\vc{R}\;\dVp,\label{eq:campo-vol}\\
  \vE(\vc{r})&=\frac{1}{4\pi\eo}\int_{S'}\frac{\sigma(\vc{r}\,')}{\norm{\vc{R}}^{3}}\,\vc{R}\;\dSp,\qquad
  \vE(\vc{r})=\frac{1}{4\pi\eo}\int_{l'}\frac{\lambda(\vc{r}\,')}{\norm{\vc{R}}^{3}}\,\vc{R}\;\dlp,\label{eq:campo-sup-lin}
\end{align}
con $\vc{R}=\vc{r}-\vc{r}\,'$.
\end{proposicion}
\begin{proof}
Cada elemento $\dif q'$ de la distribución se comporta como una carga puntual en $\vc{r}\,'$, que crea un campo
elemental $\dif\vE=\frac{1}{4\pi\eo}\frac{\dif q'}{R^{3}}\vc{R}$ por \cref{eq:campo-puntual}. El principio de
superposición, con la suma sobre elementos convertida en integral (el límite de las sumas de
Riemann de \cref{eq:campo-discreto}), da \cref{eq:campo-vol,eq:campo-sup-lin}, sin más que expresar
$\dif q'$ como $\rho\,\dVp$, $\sigma\,\dSp$ o $\lambda\,\dlp$.
\end{proof}

\begin{observacion}[Cuidado con las variables]\label{obs:variables}
Dentro de la integral, $\vc{r}\,'$ y, por tanto, $\vc{R}$ \emph{varían} para cada elemento de la distribución;
$\vc{r}$, en cambio, es el punto campo y permanece fijo. Si la distribución tiene la simetría adecuada, $\vE$ puede
calcularse por integración directa de las expresiones anteriores (veremos varios casos en los
seminarios, \cref{sec:seminarios}), pero esto no es lo habitual: en la mayoría de los casos la integral es difícil, y
obtendremos mucho más partido de las ecuaciones diferenciales que se deducen en las secciones siguientes.
\end{observacion}

% =========================================================================
\section{La ley de Gauss}\label{sec:gauss}
% =========================================================================
Comenzamos a estudiar las propiedades del campo electrostático. La naturaleza de \emph{cuadrado inverso} de la ley de
Coulomb permite conocer el flujo de $\vE$ a través de cualquier superficie cerrada y, de ahí, su divergencia. La herramienta
geométrica que lo hace posible es el ángulo sólido.

% -------------------------------------------------------------------------
\subsection{El ángulo sólido}\label{sec:angulo-solido}
% -------------------------------------------------------------------------
La razón entre las áreas totales de dos esferas concéntricas y los cuadrados de sus radios es constante:
$4\pi r_1^{2}/r_1^{2}=4\pi r_2^{2}/r_2^{2}=4\pi$. Lo mismo ocurre con áreas elementales vistas desde el centro
común $O$ (\figpanel{fig:angulo}{a}): si dos elementos $\dif S_1$ y $\dif S_2$ de esferas de radios $r_1$ y $r_2$ están
recortados por el mismo cono con vértice en $O$, entonces $\dif S_1/r_1^{2}=\dif S_2/r_2^{2}$. Esa constante es el
ángulo sólido, la versión tridimensional del ángulo en radianes (que es la razón entre un arco y su radio).

\begin{definicion}[Ángulo sólido]\label{def:angulo-solido}
El \emph{ángulo sólido elemental} bajo el que se ve desde un punto $O$ un elemento de superficie esférica
$\dif S$ de radio $r$ centrado en $O$ es
\begin{equation}\label{eq:angulo-esfera}
  \dOm=\frac{\dif S_1}{r_1^{2}}=\frac{\dif S_2}{r_2^{2}}
\end{equation}
(\emph{estereorradianes}, sr). Con $r=1$, es el área de la porción de esfera unidad bajo la cual se ve el elemento.
Para un elemento $\dS=\un\,\dsc$ \emph{cualquiera}, situado en un punto cuyo vector de posición relativa a $O$ es
$\vc{R}$, se proyecta sobre la esfera de radio $R$ centrada en $O$ (\figpanel{fig:angulo}{b}): el área proyectada
es $\dif S'=\dsc\cos\theta$, con $\theta$ el ángulo entre $\dS$ y $\vc{R}$, y
\begin{equation}\label{eq:angulo-general}
  \dOm=\frac{\dif S'}{R^{2}}=\frac{\dsc\cos\theta}{R^{2}}=\frac{\dS\cdot\vc{R}}{R^{3}}=\frac{\dS\cdot\uR}{R^{2}}.
\end{equation}
El signo de $\dOm$ es el de $\cos\theta$: positivo si $\vc{R}$ atraviesa el elemento «de dentro hacia fuera»
(según $\un$) y negativo en caso contrario.
\end{definicion}

% Figura: ángulo sólido. (a) esferas concéntricas, (b) elemento oblicuo y su proyección,
% (c) punto exterior: dos elementos opuestos recortados por el mismo cono.
\begin{figure}[ht]
\centering
\tikzset{
  parche/.style={draw=azulprof, line width=0.6pt, fill=azulclaro!70, fill opacity=0.9},
  parche2/.style={draw=acento, line width=0.6pt, fill=acento!25, fill opacity=0.9},
  cono/.style={azulacero, line width=0.45pt},
  vn/.style={-{Stealth[length=5pt,width=4.5pt]}, line width=1.1pt}}
\begin{minipage}[t]{0.50\textwidth}\centering
\begin{tikzpicture}[baseline=(current bounding box.north), x=1cm, y=1cm]
  % (a) dos esferas concéntricas
  \shade[ball color=azulclaro!55, opacity=0.30] (0,0) circle (2.3);
  \draw[azulprof, line width=0.8pt] (0,0) circle (2.3);
  \shade[ball color=azulclaro!70, opacity=0.35] (0,0) circle (1.15);
  \draw[azulacero, line width=0.8pt] (0,0) circle (1.15);
  % ecuadores (círculos máximos en perspectiva)
  \draw[acento, line width=0.8pt] (-2.3,0) arc (180:360:{2.3} and 0.5);
  \draw[acento, line width=0.5pt, densely dashed] (-2.3,0) arc (180:0:{2.3} and 0.5);
  \draw[acento, line width=0.8pt] (-1.15,0) arc (180:360:{1.15} and 0.25);
  \draw[acento, line width=0.5pt, densely dashed] (-1.15,0) arc (180:0:{1.15} and 0.25);
  % cono hacia los dos parches (dirección 52 grados)
  \def\dirr{48}
  \def\ra{0.86}\def\rb{1.72}
  \draw[cono] (0,0) -- ($(\dirr+12:\rb)$) (0,0) -- ($(\dirr-12:\rb)$);
  \begin{scope}[shift={(\dirr:\ra)}, rotate=\dirr+90]
    \path[parche] (0,0) ellipse (0.19 and 0.075);
  \end{scope}
  \begin{scope}[shift={(\dirr:\rb)}, rotate=\dirr+90]
    \path[parche] (0,0) ellipse (0.38 and 0.15);
  \end{scope}
  \fill[azulprof] (0,0) circle (1.6pt);
  \node[font=\small, anchor=north east] at (-0.05,-0.12) {$O$};
  \node[font=\small, text=azulprof, anchor=east] at ($(\dirr:\ra)+(-0.22,0.1)$) {$\dif S_1$};
  \node[font=\small, text=azulprof, anchor=south west] at ($(\dirr:\rb)+(0.2,0.12)$) {$\dif S_2$};
  % radios sobre el plano ecuatorial
  \draw[vn, verdeteo, line width=0.9pt] (0,0) -- (194:1.08) node[pos=0.55, above, font=\scriptsize, text=verdeteo] {$r_1$};
  \draw[vn, verdeteo, line width=0.9pt] (0,0) -- (-12:2.2) node[pos=0.72, above, font=\scriptsize, text=verdeteo] {$r_2$};
\end{tikzpicture}\\[2pt]{\small\scshape (a)}
\end{minipage}\hfill
\begin{minipage}[t]{0.49\textwidth}\centering
\begin{tikzpicture}[baseline=(current bounding box.north), x=1cm, y=1cm]
  % (b) superficie cerrada S, esfera unidad y esfera S' de radio R
  \draw[azulprof, line width=0.9pt, fill=azulfondo, fill opacity=0.6]
    plot[smooth cycle, tension=0.75] coordinates {(-0.4:2.7) (30:3.0) (68:3.4) (120:3.0) (172:2.6) (220:3.0) (268:2.6) (318:3.1)};
  \draw[azulacero, line width=0.5pt] (0,0) circle (0.9);
  \draw[azulacero, line width=0.9pt] (6:2.9) arc (6:70:2.9);
  \draw[azulacero, line width=0.9pt, densely dashed] (70:2.9) arc (70:92:2.9);
  \draw[azulacero, line width=0.9pt, densely dashed] (6:2.9) arc (6:-26:2.9);
  \fill[azulprof] (0,0) circle (1.6pt); \node[font=\small, anchor=north east] at (-0.05,-0.1) {$O$};
  \node[font=\scriptsize, text=azulprof] at (0.0,0.42) {$1$};
  \node[font=\small, text=azulprof] at (150:2.2) {$S$};
  \node[font=\small, text=azulacero] at (-20:2.45) {$S'$};
  % cono
  \def\dirr{30}
  \draw[cono] (0,0) -- ($(\dirr+9:2.9)$) (0,0) -- ($(\dirr-9:2.9)$);
  \begin{scope}[shift={(\dirr:0.9)}, rotate=\dirr+90] \path[parche] (0,0) ellipse (0.14 and 0.05); \end{scope}
  % dS' (proyección sobre S') y dS (inclinado un ángulo theta)
  \begin{scope}[shift={(\dirr:2.9)}, rotate=\dirr+90] \path[parche] (0,0) ellipse (0.46 and 0.15); \end{scope}
  \begin{scope}[shift={(\dirr:2.9)}, rotate=\dirr+90+30] \path[parche2] (0,0) ellipse (0.52 and 0.16); \end{scope}
  \draw[vn, azulacero] (\dirr:2.9) -- ($(\dirr:2.9)+(\dirr:0.95)$) node[right, font=\small, text=azulacero] {$\dif\vc{S}'$};
  \draw[vn, acento] (\dirr:2.9) -- ($(\dirr:2.9)+(\dirr+30:0.95)$) node[left=1pt, font=\small, text=acento] {$\dif\vc{S}$};
  \draw[verdeteo, line width=0.5pt] ($(\dirr:2.9)+(\dirr:0.55)$) arc (\dirr:\dirr+30:0.55);
  \node[font=\scriptsize, text=verdeteo] at ($(\dirr+15:2.9)+(\dirr+15:0.85)$) {$\theta$};
  \draw[vn, azulprof, line width=0.9pt] (0,0) -- ($(\dirr:2.9)-(\dirr:0.2)$) node[pos=0.55, below right=-1pt, font=\small, text=azulprof] {$\vc{R}$};
\end{tikzpicture}\\[2pt]{\small\scshape (b)}
\end{minipage}\\[10pt]
\begin{minipage}[t]{0.8\textwidth}\centering
\begin{tikzpicture}[baseline=(current bounding box.north), x=1cm, y=1cm]
  % (c) O fuera de S: el cono recorta dos elementos en lados opuestos
  \def\ang{7.5}
  \fill[azulprof] (0,0) circle (1.6pt); \node[font=\small, anchor=east] at (-0.1,0) {$O$};
  \draw[azulprof, line width=0.9pt, fill=azulfondo, fill opacity=0.6]
    plot[smooth cycle, tension=0.8] coordinates {(4.2,0.7) (5.0,1.25) (6.2,1.3) (7.3,0.95) (7.9,-0.1) (7.2,-1.2) (5.9,-1.3) (4.7,-0.9)};
  \draw[cono] (0,0) -- (7.55,{7.55*tan(\ang)}) (0,0) -- (7.55,{-7.55*tan(\ang)});
  \draw[azulprof, densely dashed, line width=0.6pt] (0,0) -- (6.9,0);
  % parche de la esfera unidad
  \draw[azulacero, line width=0.5pt] (0.9,0.6) arc (36:-36:1.0);
  \begin{scope}[shift={(1.0,0)}] \path[parche] (0,0) ellipse (0.045 and 0.135); \end{scope}
  \node[font=\scriptsize, text=azulprof, anchor=south] at (0.55,0.07) {$1$};
  % elemento cercano dS2 (entrada) y lejano dS1 (salida)
  \begin{scope}[shift={(4.85,0)}, rotate=20] \path[parche2] (0,0) ellipse (0.14 and 0.46); \end{scope}
  \begin{scope}[shift={(7.35,0)}, rotate=-10] \path[parche] (0,0) ellipse (0.18 and 0.62); \end{scope}
  \draw[vn, acento] (4.85,0) -- ++(150:0.95) node[above left=-1pt, font=\small, text=acento] {$\dif\vc{S}_2$};
  \draw[vn, azulacero] (7.35,0) -- ++(25:0.95) node[above right, font=\small, text=azulacero] {$\dif\vc{S}_1$};
  \draw[verdeteo, line width=0.5pt] (5.3,0) arc (0:150:0.45);
  \node[font=\scriptsize, text=verdeteo] at (4.55,0.55) {$\theta_2$};
  \draw[verdeteo, line width=0.5pt] (7.75,0) arc (0:25:0.4);
  \node[font=\scriptsize, text=verdeteo, anchor=west] at (7.78,0.12) {$\theta_1$};
  \node[font=\scriptsize, text=azulprof, anchor=north] at (2.6,-0.05) {$\vc{R}_2$};
  \node[font=\scriptsize, text=azulprof, anchor=north] at (6.15,-0.05) {$\vc{R}_1$};
  \node[font=\small, text=azulprof] at (6.3,0.85) {$S$};
\end{tikzpicture}\\[2pt]{\small\scshape (c)}
\end{minipage}
\caption{Ángulo sólido. (a)~Dos elementos de esferas concéntricas, $\dif S_1$ y $\dif S_2$, recortados por el mismo cono de vértice $O$: $\dif S_1/r_1^{2}=\dif S_2/r_2^{2}=\dOm$.
(b)~Para un elemento $\dS$ cualquiera de una superficie cerrada $S$ que rodea a $O$ basta proyectarlo sobre la esfera $S'$ de radio $R$: $\dif S'=\dsc\cos\theta$ y $\dOm=\dif S'/R^{2}$.
(c)~Si $O$ queda fuera de $S$, el mismo cono recorta dos elementos en lados opuestos de $S$; el de entrada ($\theta_2>\pi/2$) aporta $-\dOm$ y el de salida ($\theta_1<\pi/2$) aporta $+\dOm$, y se anulan.}
\label{fig:angulo}
\end{figure}


\begin{lema}[Ángulo sólido total de una superficie cerrada]\label{lema:omega-4pi}
Sea $S$ una superficie cerrada, orientada con la normal exterior. Entonces
\begin{equation}\label{eq:omega}
  \oint_{S}\dOm=\oint_{S}\frac{\dS\cdot\uR}{R^{2}}=\begin{cases}4\pi,&\text{si $O$ está en el interior de $S$,}\\[2pt]0,&\text{si $O$ está en el exterior de $S$.}\end{cases}
\end{equation}
\end{lema}
\begin{proof}
\emph{(a) $O$ en el interior.} Cada semirrecta que sale de $O$ corta a $S$ un número impar de veces, y la
normal exterior forma ángulos alternadamente agudos (salida) y obtusos (entrada) con la dirección de la semirrecta:
$+\dOm_0,-\dOm_0,+\dOm_0,\dots$, donde $\dOm_0$ es el ángulo sólido del cono. Como el último cruce es una
salida, la contribución neta de la dirección es $+\dOm_0$, y sumando sobre todas las direcciones se obtiene el ángulo sólido
total del espacio. Para una superficie estrellada respecto de $O$ (un único cruce por semirrecta) basta proyectar sobre una esfera
$S'$ de radio $R$ centrada en $O$, en la que $R$ es constante:
\[
  \Omega=\oint_{S}\dOm=\oint_{S'}\frac{\dif S'}{R^{2}}=\frac{1}{R^{2}}\oint_{S'}\dif S'=\frac{4\pi R^{2}}{R^{2}}=4\pi .
\]
\emph{(b) $O$ en el exterior.} Ahora cada semirrecta que corta a $S$ lo hace un número \emph{par} de veces, con entradas y
salidas alternadas (\figpanel{fig:angulo}{c}). Para un elemento de entrada, $\theta_2>\pi/2$ y por tanto
$\dS_2\cdot\vc{R}_2/R_2^{3}=\dsc_2\cos\theta_2/R_2^{2}=-\dOm_0$, mientras que para el de salida
$\dS_1\cdot\vc{R}_1/R_1^{3}=\dsc_1\cos\theta_1/R_1^{2}=+\dOm_0$. Se anulan por parejas y la integral es nula.
\end{proof}

% -------------------------------------------------------------------------
\subsection{Ley de Gauss en forma integral}\label{sec:gauss-integral}
% -------------------------------------------------------------------------
Calculemos ahora el flujo del campo eléctrico $\flujo=\oint_{S}\vE\cdot\dS$ en distintos supuestos. Observemos de entrada que
para una carga puntual en $\vc{r}\,'$, con $\vc{R}=\vc{r}-\vc{r}\,'$, el integrando contiene exactamente el ángulo sólido:
\[
  \vE\cdot\dS=\frac{\qp}{4\pi\eo}\,\frac{\uR\cdot\dS}{R^{2}}=\frac{\qp}{4\pi\eo}\,\dOm .
\]

\begin{teorema}[Ley de Gauss, forma integral]\label{teo:gauss-integral}
El flujo del campo eléctrico a través de una superficie cerrada cualquiera $S$ es igual a la carga
encerrada por ella dividida entre $\eo$:
\begin{equation}\label{eq:gauss-int}
  \boxed{\;\flujo=\oiint_{S}\vE\cdot\dS=\frac{\Qenc}{\eo}\;}
\end{equation}
con independencia de que la carga sea una distribución discreta o continua. Si hay una distribución volúmica,
$\Qenc=\int_{V}\rho\,\dV$, con $V$ el volumen contenido en $S$.
\end{teorema}
\begin{proof}
\emph{Carga puntual fuera de $S$.} Usando el \cref{lema:omega-4pi}(b),
\[
  \flujo=\oiint_{S}\frac{\qp}{4\pi\eo}\frac{\uR\cdot\dS}{R^{2}}=\frac{\qp}{4\pi\eo}\oiint_{S}\dOm=\frac{\qp}{4\pi\eo}\,(\Omega-\Omega)=0 .
\]
\emph{Carga puntual dentro de $S$.} Por el \cref{lema:omega-4pi}(a),
\[
  \flujo=\frac{\qp}{4\pi\eo}\oiint_{S}\dOm=\frac{\qp}{4\pi\eo}\,4\pi=\frac{\qp}{\eo}.
\]
\emph{Sistema de cargas puntuales.} Por superposición, el flujo es la suma de los flujos de cada carga; solo las
interiores aportan:
\[
  \flujo=\sum_{\text{dentro}}\frac{\qp_i}{\eo}+\sum_{\text{fuera}}0=\frac{\Qenc}{\eo}.
\]
\emph{Distribución continua.} Se aplica el resultado anterior a cada elemento $\dif q'=\rho\,\dVp$ y se integra:
$\flujo=\frac{1}{\eo}\int_{V}\rho\,\dV$.
\end{proof}

\begin{observacion}[Alcance de la ley de Gauss]\label{obs:alcance-gauss}
\begin{enumerate}[label=(\alph*)]
  \item El flujo de $\vE$ a través de una superficie cerrada depende \emph{solo} de la carga encerrada.
  \item El valor del campo puede variar de un punto a otro de la superficie, ser nulo o no, aunque el flujo sea nulo: en un
        dipolo, el flujo a través de una superficie que lo rodee es nulo, pero el campo no se anula en ningún punto de ella.
  \item Las cargas exteriores no afectan al flujo, pero sí al valor del campo en la superficie.
  \item La posición de las cargas interiores no afecta al flujo, pero sí al valor del campo.
\end{enumerate}
\end{observacion}

% Figura: flujo de una carga puntual dentro y fuera de una superficie cerrada.
\begin{figure}[ht]
\centering
\begin{minipage}[t]{0.48\textwidth}\centering
\begin{tikzpicture}[baseline=(current bounding box.north), scale=0.95]
  \draw[azulprof, line width=0.8pt, fill=azulfondo, fill opacity=0.55] plot[smooth cycle, tension=0.8] coordinates {(1.9,0.2) (1.3,1.4) (-0.2,1.7) (-1.6,1.1) (-1.9,-0.3) (-0.9,-1.5) (0.6,-1.3)};
  \node[font=\small, text=azulprof] at (1.35,-1.05) {$S$};
  \foreach \a in {0,30,...,330}{\draw[-{Stealth[length=4pt]}, cm1, line width=0.6pt] (0.2,0.0) ++(\a:0.15) -- ++(\a:2.5);}
  \node[cargapos] at (0.2,0.0) {$+$};
  \node[font=\small, anchor=south west] at (0.3,0.1) {$\qp$};
\end{tikzpicture}\\[2pt]{\small\scshape (a)}
\end{minipage}\hfill
\begin{minipage}[t]{0.48\textwidth}\centering
\begin{tikzpicture}[baseline=(current bounding box.north), scale=0.95]
  \draw[azulprof, line width=0.8pt, fill=azulfondo, fill opacity=0.55, name path=blob] plot[smooth cycle, tension=0.8] coordinates {(1.9,0.2) (1.3,1.4) (-0.2,1.7) (-1.6,1.1) (-1.9,-0.3) (-0.9,-1.5) (0.6,-1.3)};
  \node[font=\small, text=azulprof] at (1.35,-1.05) {$S$};
  \coordinate (Q) at (-3.4,0.1);
  \foreach \a in {-34,-22,-10,2,14,26}{
    \path[name path=r\a] (Q) -- ++(\a:7.0);
    \path[name intersections={of=blob and r\a, name=I\a, total=\t}];
    \draw[-{Stealth[length=4pt]}, cm1, line width=0.6pt] (Q) ++(\a:0.15) -- ++(\a:6.0);}
  \foreach \a in {-10,2,14}{
    \path[name path=r\a] (Q) -- ++(\a:7.0);
    \path[name intersections={of=blob and r\a, sort by=r\a, name=J\a}];
    \draw[azulprof, fill=white, line width=0.5pt] (J\a-1) circle (1.8pt);
    \fill[azulprof] (J\a-2) circle (1.8pt);}
  \node[cargapos] at (Q) {$+$};
  \node[font=\small, anchor=south east] at ($(Q)+(-0.1,0.1)$) {$\qp$};
\end{tikzpicture}\\[2pt]{\small\scshape (b)}
\end{minipage}
\caption{Flujo de una carga puntual. (a)~Carga interior: todas las líneas de campo que parten de $\qp$
atraviesan $S$ hacia fuera y el flujo vale $\qp/\eo$ (equivale a $4\pi$ de ángulo sólido). (b)~Carga exterior:
cada línea que entra por $S$ ($\circ$) sale después por otro punto ($\bullet$); el flujo de entrada compensa
el de salida y el flujo neto es nulo, aunque el campo no lo sea en $S$.}
\label{fig:flujo-gauss}
\end{figure}


% -------------------------------------------------------------------------
\subsection{Ley de Gauss en forma diferencial}\label{sec:gauss-diferencial}
% -------------------------------------------------------------------------
La forma integral es una afirmación sobre regiones extensas; la forma diferencial es su traducción local, punto a punto. Para
pasar de una a otra usamos el teorema de la divergencia del Tema~0 (\cref{eq:div-teorema} más abajo).

\begin{teorema}[Teorema de la divergencia o de Gauss~\cite{TEMA0}]\label{teo:divergencia}
Si $\vc{A}$ es un campo vectorial de clase $C^{1}$ en un entorno de un sólido acotado $V$ de borde $S=\partial V$ orientado con
la normal exterior,
\begin{equation}\label{eq:div-teorema}
  \iiint_{V}\nab\cdot\vc{A}\,\dV=\oiint_{\partial V}\vc{A}\cdot\dS.
\end{equation}
\end{teorema}

\begin{teorema}[Ley de Gauss, forma diferencial: primera ecuación de Maxwell]\label{teo:gauss-dif}
En el vacío, el campo electrostático generado por una densidad volúmica de carga $\rho$ cumple
\begin{equation}\label{eq:gauss-dif}
  \boxed{\;\nab\cdot\vE(\vc{r})=\frac{\rho(\vc{r})}{\eo}\;}
\end{equation}
\end{teorema}
\begin{proof}
Sea $V$ un volumen cualquiera. Por la forma integral (\cref{eq:gauss-int}) y por el teorema de la divergencia
(\cref{eq:div-teorema}),
\[
  \iiint_{V}\nab\cdot\vE\,\dV=\oiint_{\partial V}\vE\cdot\dS=\frac{1}{\eo}\iiint_{V}\rho\,\dV .
\]
Igualando, $\iiint_{V}\bigl(\nab\cdot\vE-\rho/\eo\bigr)\dV=0$ para \emph{todo} volumen $V$. Si el integrando fuese no nulo en
algún punto, por continuidad tendría signo constante en un entorno suyo y la integral sobre ese entorno no sería cero;
por tanto, el integrando es idénticamente nulo.
\end{proof}

Esta ecuación pone de manifiesto que las \emph{fuentes} del campo eléctrico son las cargas, cuya densidad es $\rho$: donde
$\rho>0$ el campo «diverge» (nacen líneas de campo) y donde $\rho<0$ «converge». Es una consecuencia de la ley de Coulomb
y de su dependencia con el inverso del cuadrado de la distancia, y es válida siempre que estemos en el vacío. Falta aún
resolver el caso de las cargas puntuales: en el punto donde está la carga, $\rho$ diverge, y necesitamos una herramienta
matemática para escribir la ecuación. Es la delta de Dirac.

% -------------------------------------------------------------------------
\subsection{La delta de Dirac}\label{sec:delta}
% -------------------------------------------------------------------------
\begin{definicion}[Delta de Dirac]\label{def:delta}
La \emph{delta de Dirac} se define como el límite, cuando $\alpha\to\infty$, de una sucesión de funciones $f_\alpha(x)$ (por
ejemplo, un pulso rectangular de anchura $1/\alpha$ y altura $\alpha$, un triángulo de base $2/\alpha$ y altura $\alpha$ o una
gaussiana estrecha de área unidad) (\cref{fig:delta}):
\[
  \delta(x)=\lim_{\alpha\to\infty}f_\alpha(x),\qquad
  \delta(x)=\begin{cases}0,&x\neq0,\\\infty,&x=0,\end{cases}\qquad
  \int_{-\infty}^{\infty}\delta(x)\,\dif x=1 .
\]
Análogamente, $\delta(x-a)$ es nula salvo en $x=a$, donde diverge. Su propiedad esencial es la de \emph{muestreo}:
\begin{equation}\label{eq:delta-muestreo}
  \int_{-\infty}^{\infty}f(x)\,\delta(x-a)\,\dif x=f(a)
\end{equation}
para toda $f$ continua en $a$.
\end{definicion}

% Figura: sucesiones que definen la delta de Dirac.
\begin{figure}[ht]
\centering
\pgfplotsset{deltaax/.style={tema, width=0.335\textwidth, height=4.3cm, xmin=-1.3, xmax=1.3, ymin=0, ymax=4.4,
   xtick={-1,0,1}, ytick={0,1,2,4}, xlabel={$x$}, grid=major}}
\begin{tikzpicture}
  \begin{axis}[deltaax, name=A, ylabel={$f_\alpha(x)$}, title={\small\scshape (a)\; pulsos}]
    \addplot[azulclaro!80!azulacero, line width=1.1pt, const plot, mark=none] coordinates {(-1.3,0) (-0.5,0) (-0.5,1) (0.5,1) (0.5,0) (1.3,0)};
    \addplot[azulacero, line width=1.1pt, const plot, mark=none] coordinates {(-1.3,0) (-0.25,0) (-0.25,2) (0.25,2) (0.25,0) (1.3,0)};
    \addplot[azulprof, line width=1.1pt, const plot, mark=none] coordinates {(-1.3,0) (-0.125,0) (-0.125,4) (0.125,4) (0.125,0) (1.3,0)};
  \end{axis}
  \begin{axis}[deltaax, at={(A.east)}, anchor=west, xshift=0.6cm, yticklabels={}, title={\small\scshape (b)\; triángulos}]
    \addplot[azulclaro!80!azulacero, line width=1.1pt] coordinates {(-1.3,0) (-1,0) (0,1) (1,0) (1.3,0)};
    \addplot[azulacero, line width=1.1pt] coordinates {(-1.3,0) (-0.5,0) (0,2) (0.5,0) (1.3,0)};
    \addplot[azulprof, line width=1.1pt] coordinates {(-1.3,0) (-0.25,0) (0,4) (0.25,0) (1.3,0)};
  \end{axis}
  \begin{axis}[deltaax, at={(A.east)}, anchor=west, xshift=5.85cm, yticklabels={}, samples=200, domain=-1.3:1.3,
               title={\small\scshape (c)\; gaussianas}]
    \addplot[azulclaro!80!azulacero, line width=1.1pt] {1/(0.6*sqrt(2*pi))*exp(-x^2/(2*0.6^2))};
    \addplot[azulacero, line width=1.1pt] {1/(0.3*sqrt(2*pi))*exp(-x^2/(2*0.3^2))};
    \addplot[azulprof, line width=1.1pt] {1/(0.1*sqrt(2*pi))*exp(-x^2/(2*0.1^2))};
  \end{axis}
\end{tikzpicture}
\caption{La delta de Dirac como límite ($\alpha\to\infty$) de picos cada vez más estrechos y altos, todos de
\emph{área unidad}. (a)~Pulsos rectangulares de anchura $1/\alpha$ y altura $\alpha$ ($\alpha=1,2,4$).
(b)~Triángulos de base $2/\alpha$ y altura $\alpha$. (c)~Gaussianas $\frac{1}{\sigma\sqrt{2\pi}}\eul^{-x^{2}/2\sigma^{2}}$
con $\sigma=0{,}6;\,0{,}3;\,0{,}1$. Cuando el pico se afila, el área no cambia y el valor en $x=0$ diverge.}
\label{fig:delta}
\end{figure}


Estrictamente, la delta no es una función sino una \emph{función generalizada} o \emph{distribución}, y solo adquiere sentido
bajo el signo integral. La propiedad de muestreo se entiende así: como $\delta(x-a)=0$ para $x\neq a$, en la integral solo cuenta el
valor de $f$ en $x=a$, que sale fuera de ella: $\int f(x)\delta(x-a)\dif x=f(a)\int\delta(x-a)\dif x=f(a)$.

\begin{ejemplo}[Cálculos con la delta]\label{ej:delta}
Con \cref{eq:delta-muestreo} (la delta cae en el interior del intervalo de integración):
\[
  \int_{-\infty}^{\infty}(3x-1)\,\delta(x+1)\,\dif x=3(-1)-1=-4,\qquad
  \int_{0}^{3}x^{3}\,\delta(x-2)\,\dif x=2^{3}=8 .
\]
\end{ejemplo}

En tres dimensiones se define $\delta^{3}(\vc{r})=\delta(x)\,\delta(y)\,\delta(z)$ y, para un punto $\vc{a}$ del espacio,
\begin{equation}\label{eq:delta3}
\begin{gathered}
  \delta^{3}(\vc{r}-\vc{a})=\begin{cases}0,&\vc{r}\neq\vc{a},\\\infty,&\vc{r}=\vc{a},\end{cases}\qquad
  \int_{V}\delta^{3}(\vc{r}-\vc{a})\,\dV=\begin{cases}1,&\vc{a}\in V,\\0,&\vc{a}\notin V,\end{cases}\\[4pt]
  \int_{V}f(\vc{r})\,\delta^{3}(\vc{r}-\vc{a})\,\dV=\begin{cases}f(\vc{a}),&\vc{a}\in V,\\0,&\vc{a}\notin V.\end{cases}
\end{gathered}
\end{equation}

\begin{ejemplo}[Delta tridimensional]\label{ej:delta3}
$\displaystyle\int_{\R^{3}}\bigl(r^{2}+\vc{r}\cdot\vc{a}+a^{2}\bigr)\,\delta^{3}(\vc{r}-\vc{a})\,\dV=a^{2}+a^{2}+a^{2}=3a^{2}$.
\end{ejemplo}

Con ella podemos precisar qué ocurre con el campo de una carga puntual.

\begin{teorema}[Divergencia del campo coulombiano y laplaciano de $1/R$]\label{teo:delta}
Sea $\vc{R}=\vc{r}-\vc{r}\,'$. Entonces, derivando respecto de $\vc{r}$,
\begin{equation}\label{eq:div-delta}
  \nab\cdot\Bigl(\frac{\uR}{R^{2}}\Bigr)=\nab\cdot\Bigl(\frac{\vc{R}}{R^{3}}\Bigr)=4\pi\,\delta^{3}(\vc{R}),
  \qquad
  \Lap\Bigl(\frac{1}{R}\Bigr)=-4\pi\,\delta^{3}(\vc{R}).
\end{equation}
\end{teorema}
\begin{proof}
Tomemos el origen en la fuente ($\vc{r}\,'=\vc{0}$, $\vc{R}=\vc{r}$). Para $R\neq0$ el campo es radial, $A_R=1/R^{2}$, y en coordenadas
esféricas $\nab\cdot\vc{A}=\frac{1}{R^{2}}\pd{}{R}\bigl(R^{2}A_R\bigr)=\frac{1}{R^{2}}\pd{(1)}{R}=0$. Pero en $R=0$ el campo diverge. Para
ver qué ocurre allí, integramos sobre un volumen $V$ cualquiera con el teorema de la divergencia y el \cref{lema:omega-4pi}:
\[
  \int_{V}\nab\cdot\Bigl(\frac{\uR}{R^{2}}\Bigr)\dV=\oiint_{\partial V}\frac{\uR\cdot\dS}{R^{2}}=\oiint_{\partial V}\dOm
  =\begin{cases}0,&\text{si }R=0\text{ no está en }V,\\4\pi,&\text{si }R=0\text{ sí está en }V.\end{cases}
\]
Una cantidad que vale cero fuera del origen y cuya integral sobre cualquier volumen que lo contiene es $4\pi$ es, por definición,
$4\pi\,\delta^{3}(\vc{R})$. Para el laplaciano, como $\nab(1/R)=-\vc{R}/R^{3}$, se tiene
$\Lap(1/R)=\nab\cdot\nab(1/R)=-\nab\cdot(\vc{R}/R^{3})=-4\pi\delta^{3}(\vc{R})$.
\end{proof}

\begin{observacion}[Ni cero ni infinito]\label{obs:delta-div}
En coordenadas esféricas, $\Lap(1/R)=\frac{1}{R^{2}}\pd{}{R}\bigl(R^{2}\pd{}{R}\frac{1}{R}\bigr)$ vale cero para $R\neq0$ y no está definido en
$R=0$, donde diverge, igual que $\delta(R)$. El teorema precisa cuánto: lo suficiente para que la integral sobre un entorno
del origen sea $-4\pi$. Un cálculo ingenuo que olvide el origen diría que $\nab\cdot(\uR/R^{2})=0$ en todas partes y
que el flujo $4\pi$ de este campo sale «de la nada»; la delta es lo que restituye la coherencia entre la divergencia y el flujo.
\end{observacion}

\begin{corolario}[Ley de Gauss diferencial de una carga puntual]\label{cor:gauss-puntual}
Para una carga puntual $\qp$ en $\vc{r}\,'$, $\rho(\vc{r})=\qp\,\delta^{3}(\vc{r}-\vc{r}\,')$ y
\[
  \nab\cdot\vE(\vc{r})=\frac{\qp}{4\pi\eo}\,\nab\cdot\Bigl(\frac{\vc{R}}{R^{3}}\Bigr)=\frac{\qp}{4\pi\eo}\,4\pi\,\delta^{3}(\vc{R})
  =\frac{\qp}{\eo}\,\delta^{3}(\vc{R})=\frac{\rho}{\eo}.
\]
La divergencia del campo de una carga puntual es nula en todas partes salvo en el punto que ocupa la carga, donde diverge.
\end{corolario}

Con la delta podemos dar una segunda demostración, directa, de la ley de Gauss diferencial \emph{sin pasar por la forma integral}.

\begin{proof}[Demostración directa de \cref{eq:gauss-dif}]
Partimos del campo de una distribución volúmica (\cref{eq:campo-vol}) y derivamos respecto del punto campo $\vc{r}$, es decir, respecto de las
coordenadas del punto campo y no de las de la fuente, de modo que $\nab$ actúa solo sobre $\vc{R}/R^{3}$ y se puede introducir bajo
la integral:
\[
  \nab\cdot\vE(\vc{r})=\frac{1}{4\pi\eo}\int_{V'}\rho(\vc{r}\,')\,\nab\cdot\Bigl(\frac{\vc{R}}{R^{3}}\Bigr)\dVp
  =\frac{1}{4\pi\eo}\int_{V'}\rho(\vc{r}\,')\,4\pi\,\delta^{3}(\vc{r}-\vc{r}\,')\,\dVp=\frac{\rho(\vc{r})}{\eo},
\]
por el \cref{teo:delta} y la propiedad de muestreo (\cref{eq:delta3}).
\end{proof}

% -------------------------------------------------------------------------
\subsection{Cálculo del campo con la ley de Gauss: simetrías}\label{sec:gauss-simetria}
% -------------------------------------------------------------------------
Si la distribución de carga posee suficiente simetría, la ley de Gauss permite calcular $\vE$ de manera sencilla. Hay que hacer
hincapié en que la ley es válida \emph{siempre}, pero que despejar $\vE$ de ella solo es posible si existe la simetría adecuada:
hay que elegir una superficie de integración (\emph{superficie gaussiana}, una construcción matemática que no tiene por qué ser real)
en la cual, en todos sus puntos, $\vE$ sea constante en módulo y, además, paralelo o perpendicular a $\dS$. Sólo hay tres
simetrías que funcionan:

\begin{center}\small
\renewcommand{\arraystretch}{1.3}
\begin{tabular}{@{}lll@{}}
\toprule
Simetría & Superficie de Gauss & Campo\\
\midrule
Esférica & esfera de radio $r$ & $\vE=E(r)\,\ur$\\
Cilíndrica & cilindro coaxial de radio $r$ y altura $\ell$ & $\vE=E(r)\,\urho$\\
Plana & cilindro (o paralelepípedo) recto que atraviesa el plano & $\vE=E(z)\,\uz$\\
\bottomrule
\end{tabular}
\end{center}

En esos casos, $\oiint_{S}\vE\cdot\dS=E\,S$ y se despeja $E(\vc{r})=\Qenc(\vc{r})/(\eo\,S(\vc{r}))$. Los seminarios
\cref{sem:esfera,sem:hilo-gauss,sem:plano-gauss,sem:lamina,sem:cascara,sem:cilindro} aplican el método a cada una de ellas.

\begin{ejemplo}[De los campos a las fuentes]\label{ej:inverso}
La ley de Gauss diferencial también se usa en sentido inverso: dado el campo, \emph{deducir la carga}. Si $\vE(\vc{r})=k\,z\,\uz$,
entonces $\nab\cdot\vE=k$ y por \cref{eq:gauss-dif} la densidad es $\rho=k\,\eo$: una distribución volúmica de carga \emph{uniforme} en todo el espacio.
\end{ejemplo}

% =========================================================================
\section{El potencial eléctrico}\label{sec:potencial}
% =========================================================================
Hasta ahora hemos descrito los efectos electrostáticos únicamente mediante el campo vectorial $\vE$. Escribiendo el campo de una
carga puntual de otra forma podremos expresar la misma información mediante un \emph{campo escalar}, lo cual es más
conveniente para muchos propósitos: un número por punto en lugar de tres.

% -------------------------------------------------------------------------
\subsection{Definición del potencial}\label{sec:def-potencial}
% -------------------------------------------------------------------------
Partimos del campo de una carga puntual (\cref{eq:campo-puntual}) y usamos que $\vc{R}/R^{3}=-\nab(1/R)$, que es el caso $n=-1$ del
gradiente de $R^{n}$~\cite{TEMA0} ($\nab R^{n}=nR^{n-1}\uR$):
\begin{equation}\label{eq:campo-grad}
  \vE(\vc{r})=\frac{\qp}{4\pi\eo}\,\frac{\vc{R}}{R^{3}}=-\frac{\qp}{4\pi\eo}\,\nab\Bigl(\frac{1}{R}\Bigr)
  =-\nab\Bigl(\frac{\qp}{4\pi\eo R}\Bigr).
\end{equation}

\begin{definicion}[Potencial electrostático]\label{def:potencial}
Se llama \emph{potencial eléctrico} (o escalar) de una carga puntual $\qp$ al campo escalar
\begin{equation}\label{eq:potencial-puntual}
  \phi(\vc{r})=\frac{1}{4\pi\eo}\,\frac{\qp}{R}+\text{cte},\qquad R=\norm{\vc{r}-\vc{r}\,'}.
\end{equation}
La relación entre el campo vectorial $\vE:\R^{3}\to\R^{3}$ y el campo escalar $\phi:\R^{3}\to\R$ es
\begin{equation}\label{eq:E-grad-phi}
  \boxed{\;\vE(\vc{r})=-\nab\phi(\vc{r})\;}
\end{equation}
\end{definicion}

\begin{observacion}[Propiedades del potencial]\label{obs:potencial}
\begin{enumerate}[label=(\alph*)]
  \item \emph{Significado del gradiente.} $\norm{\nab\phi}$ es la derivada de $\phi$ en la dirección de máxima variación, y $\nab\phi$
        apunta en esa dirección, que coincide con la de $\vE$ en cada punto; pero, por el signo menos, $\vE$ y $\nabla\phi$ tienen
        sentidos opuestos: el potencial \emph{aumenta} al movernos en sentido contrario al campo. Como $\dif\phi=\nab\phi\cdot\dl$,
        resulta $\dif\phi=-\vE\cdot\dl$.
  \item \emph{Unidades.} En el SI, $[\phi]=\mathrm{V}$ (voltio $=\mathrm{N\,m/C}$), y de aquí $[\vE]=\mathrm{V/m}$.
  \item \emph{No está unívocamente definido.} Sumándole cualquier constante se obtiene el mismo campo (el gradiente de una constante
        es nulo). Ocurre como con la altitud de un lugar: siempre se da respecto de una referencia, el nivel del mar. En electrostática se toma
        habitualmente como nivel de referencia (al que se asigna el valor cero) un punto infinitamente alejado de las cargas: si
        $r\to\infty$, $\phi=\lim\bigl(\frac{1}{4\pi\eo}\frac{\qp}{R}+\text{cte}\bigr)=\text{cte}$, luego $\text{cte}=0$. Pero hay
        problemas ideales en los que las cargas llegan al infinito (un hilo o un plano infinitos); en ellos hay que tomar otro punto como
        origen de potencial, e incluso a veces se asigna el cero a la superficie de la Tierra. En cada problema hay que especificarlo, y
        como el potencial no es único conviene hablar de \emph{diferencias de potencial}. A partir de ahora escribimos $\phi$ para el potencial con
        la constante ya elegida.
\end{enumerate}
\end{observacion}

Por el principio de superposición, el potencial de un sistema de cargas es la suma de los potenciales de cada una, y el de una
distribución continua la integral correspondiente:

\begin{proposicion}[Potencial de sistemas y distribuciones de carga]\label{prop:potencial-dist}
\begin{align}
  \phi(\vc{r})&=\frac{1}{4\pi\eo}\sum_{i=1}^{N}\frac{\qp_i}{R_i}+\text{cte}, &&\text{(discreta)}\label{eq:pot-discreto}\\
  \phi(\vc{r})&=\frac{1}{4\pi\eo}\int_{V'}\frac{\rho(\vc{r}\,')\,\dVp}{R}+\text{cte}, &&\text{(volúmica)}\label{eq:pot-vol}\\
  \phi(\vc{r})&=\frac{1}{4\pi\eo}\int_{S'}\frac{\sigma(\vc{r}\,')\,\dSp}{R}+\text{cte}, \quad
  \phi(\vc{r})=\frac{1}{4\pi\eo}\int_{l'}\frac{\lambda(\vc{r}\,')\,\dlp}{R}+\text{cte}. &&\text{(superficial, lineal)}\label{eq:pot-sup-lin}
\end{align}
\end{proposicion}
\begin{proof}
El gradiente es lineal y $\vE=-\nab\phi$ para cada carga puntual, con la misma constante aditiva; la superposición de campos
(\cref{eq:campo-discreto,eq:campo-vol,eq:campo-sup-lin}) es entonces la superposición de potenciales.
\end{proof}

% -------------------------------------------------------------------------
\subsection{Significado físico: el potencial como trabajo por unidad de carga}\label{sec:potencial-signif}
% -------------------------------------------------------------------------
Consideremos una carga puntual $q$ (de prueba, $q\to0$) en equilibrio bajo la acción de la fuerza electrostática $\vF_q=q\vE$ y de
una fuerza \emph{mecánica} externa $\vF$ (\cref{fig:trabajo}a): $\vF_q+\vF=q\vE+\vF=\vc{0}$. Supongamos que la carga se mueve
\emph{infinitamente despacio} (de forma reversible) de un punto $A=\vc{r}_1$ a otro $B=\vc{r}_2$ por una trayectoria dada $\Gamma$. La
velocidad es esencialmente nula en todo momento, también la energía cinética, y la carga está siempre en equilibrio o casi, de modo que la
ecuación anterior es aplicable. Este procedimiento permite calcular el trabajo reversible, en sentido termodinámico, que realiza la fuerza
mecánica contra las del campo, sin efectos disipativos ni de fricción.

% Figura: trabajo de la fuerza externa contra el campo y circulación nula.
\begin{figure}[ht]
\centering
\begin{minipage}[t]{0.52\textwidth}\centering
\begin{tikzpicture}[baseline=(current bounding box.north), x=1cm, y=1cm]
  \coordinate (A) at (0,0); \coordinate (B) at (5.2,2.2);
  \draw[trayecto] (A) .. controls (0.8,2.2) and (2.0,-0.9) .. (3.0,0.9) .. controls (3.6,1.9) and (4.4,1.2) .. (B);
  \draw[vect, cm1, line width=0.9pt] (3.0,0.9) -- ++(-0.15,-0.9) node[below, font=\small, text=cm1] {$\vc{F}_{q}=q\vE$};
  \draw[vectrojo, line width=0.9pt] (3.0,0.9) -- ++(0.15,0.9) node[above, font=\small, text=acento] {$\vc{F}=-q\vE$};
  \draw[-{Stealth[length=4pt]}, azulprof, line width=0.8pt] (3.0,0.9) -- ++(0.55,0.72) node[right, font=\scriptsize] {$\dl$};
  \node[cargapru, minimum size=8pt] at (3.0,0.9) {};
  \fill[azulprof] (A) circle (1.8pt) node[below left, font=\small] {$A$};
  \fill[azulprof] (B) circle (1.8pt) node[above right, font=\small] {$B$};
  \node[font=\small, text=acento, anchor=north west] at (0.7,2.1) {$\Gamma$};
\end{tikzpicture}\\[2pt]{\small\scshape (a)}
\end{minipage}\hfill
\begin{minipage}[t]{0.44\textwidth}\centering
\begin{tikzpicture}[baseline=(current bounding box.north), x=1cm, y=1cm]
  \draw[trayecto, orienta] (0,0) .. controls (0.7,1.8) and (2.3,2.0) .. (3.2,0.2);
  \draw[trayecto, orienta, azulacero] (3.2,0.2) .. controls (2.8,-1.2) and (0.7,-1.5) .. (0,0);
  \fill[azulprof] (0,0) circle (1.8pt) node[left, font=\small] {$1$};
  \fill[azulprof] (3.2,0.2) circle (1.8pt) node[right, font=\small] {$2$};
  \node[font=\small, text=acento] at (1.6,1.75) {$1\to2$};
  \node[font=\small, text=azulacero] at (1.6,-1.2) {$2\to1$};
  \node[font=\small, text=azulprof] at (1.6,0.2) {$C$};
\end{tikzpicture}\\[2pt]{\small\scshape (b)}
\end{minipage}
\caption{(a)~Una fuerza externa $\vF$ lleva, muy despacio, una carga de prueba de $A$ a $B$
contra la fuerza del campo $q\vE$; en cada punto ambas se compensan. (b)~Un circuito cerrado $C$ se
descompone en dos tramos de $1$ a $2$ y de $2$ a $1$; si la integral de línea de $\vE$ no depende del
camino, la circulación a lo largo de $C$ es nula.}
\label{fig:trabajo}
\end{figure}


\begin{proposicion}[Trabajo y diferencia de potencial]\label{prop:trabajo}
El trabajo realizado por la fuerza mecánica $\vF=-q\vE$ para llevar la carga $q$ de $A$ a $B$ por cualquier camino $\Gamma$ es
\begin{equation}\label{eq:trabajo}
  W_{AB}=\int_{A}^{B}\vF\cdot\dl=-q\int_{A}^{B}\vE\cdot\dl=q\int_{A}^{B}\nab\phi\cdot\dl=q\bigl[\phi(B)-\phi(A)\bigr].
\end{equation}
En particular, para mover una carga \emph{positiva unidad} desde el infinito hasta un punto $P$ (con $\phi(\infty)=0$),
\begin{equation}\label{eq:trabajo-unidad}
  W=\int_{\infty}^{P}\vF\cdot\dif\vc{r}=-\int_{\infty}^{P}\vE\cdot\dif\vc{r}=\int_{\infty}^{P}\dif\phi=\phi_P-\phi_\infty=\phi_P.
\end{equation}
\end{proposicion}
\begin{proof}
Se sustituye $\vF=-q\vE$, se usa $\vE=-\nab\phi$ y el teorema fundamental de las integrales de línea (\emph{teorema del gradiente}~\cite{TEMA0}):
$\int_A^B\nab\phi\cdot\dl=\phi(B)-\phi(A)$, que no depende de $\Gamma$.
\end{proof}

El potencial no es, pues, más que el \emph{trabajo por unidad de carga} necesario para desplazar una carga unidad desde el infinito
hasta el punto, \emph{sin importar el camino seguido} (si la fuente es positiva, es trabajo realizado contra el campo). Si el
trabajo suministrado fuese mayor que $q\phi_P$, el exceso de energía se transformaría en energía cinética, el sistema dejaría de estar
en reposo y abandonaríamos la electrostática.

% -------------------------------------------------------------------------
\subsection{El campo electrostático es conservativo e irrotacional}\label{sec:rotacional}
% -------------------------------------------------------------------------
\begin{teorema}[Circulación nula: el campo electrostático es conservativo]\label{teo:conservativo}
Entre dos puntos arbitrarios,
\begin{equation}\label{eq:int-linea}
  -\int_{1}^{2}\vE\cdot\dif\vc{r}=\int_{1}^{2}\nab\phi\cdot\dif\vc{r}=\int_{1}^{2}\dif\phi=\phi(\vc{r}_2)-\phi(\vc{r}_1),
\end{equation}
luego la integral de línea de $\vE$ es independiente del camino seguido y, para cualquier circuito cerrado $C$,
\begin{equation}\label{eq:circulacion}
  \oint_{C}\vE\cdot\dif\vc{r}=\int_{1}^{2}\vE\cdot\dif\vc{r}+\int_{2}^{1}\vE\cdot\dif\vc{r}=0 .
\end{equation}
\end{teorema}
\begin{proof}
La primera cadena de igualdades es el teorema del gradiente. Para la segunda, descomponemos $C$ en dos tramos de $1$ a $2$ y de $2$ a $1$
(\cref{fig:trabajo}b); como $\int_{2}^{1}=-\int_{1}^{2}$ y la integral entre 1 y 2 no depende del camino, la suma es cero.
\end{proof}

\begin{teorema}[El campo electrostático es irrotacional]\label{teo:irrotacional}
En electrostática,
\begin{equation}\label{eq:rotacional}
  \boxed{\;\nab\times\vE=\vc{0}\;}
\end{equation}
\end{teorema}
\begin{proof}
Por el teorema de Stokes~\cite{TEMA0} y \cref{eq:circulacion}, para toda superficie $S$ con borde $C=\partial S$,
\[
  \oint_{C}\vE\cdot\dif\vc{r}=\iint_{S}(\nab\times\vE)\cdot\dS=0 .
\]
El resultado es válido para cualquier trayectoria $C$ y, por tanto, para cualquier superficie $S$; como en la demostración de la ley
de Gauss diferencial, esto obliga a que el integrando $\nab\times\vE$ sea idénticamente nulo.
\end{proof}

\begin{observacion}[Comprobación directa y alcance]\label{obs:rotacional}
La anulación del rotacional se comprueba también directamente: para una distribución volúmica,
$\nab\times\vE=\frac{1}{4\pi\eo}\int_{V'}\rho(\vc{r}\,')\,\nab\times\bigl(\vc{R}/R^{3}\bigr)\dVp=\vc{0}$, pues
$\nab\times(R^{n}\uR)=\vc{0}$ para todo $n$ (caso $n=-2$)~\cite{TEMA0}. \emph{Cuidado}: $\nab\times\vE=\vc{0}$ es un resultado de la
\emph{electrostática}; con campos magnéticos variables en el tiempo (Faraday) dejará de cumplirse. Y a la inversa: si $\vE=-\nab\phi$, entonces
$\nab\times\vE=-\nab\times(\nab\phi)=\vc{0}$ porque el rotacional de un gradiente siempre es nulo (si $\phi$ es suficientemente regular);
un campo con $\nab\times\vE=\vc{0}$ en una región simplemente conexa puede escribirse como el gradiente de un campo escalar~\cite{TEMA0}.
\end{observacion}

Conocidas la divergencia y el rotacional de un campo vectorial en todos los puntos de una región, el campo queda determinado. Es el contenido
del siguiente resultado.

\begin{teorema}[Teorema de Helmholtz~\cite{HELMHOLTZ}]\label{teo:helmholtz}
Sea $\vc{F}$ un campo vectorial con $\nab\cdot\vc{F}=b(\vc{r})$ y $\nab\times\vc{F}=\vc{c}(\vc{r})$ dados en todo el espacio, y tal que $\vc{F}\to\vc{0}$
en el infinito suficientemente deprisa (y las fuentes $b,\vc{c}$ están localizadas). Entonces $\vc{F}$ queda determinado de forma única.
\end{teorema}

En electrostática, $\nab\cdot\vE=\rho/\eo$ y $\nab\times\vE=\vc{0}$ (con las condiciones de regularidad del teorema, que el campo cumple), de modo que el
campo electrostático está \emph{unívocamente determinado} por la distribución de carga. Es la razón por la que la ley de Gauss en forma diferencial y la
irrotacionalidad contienen toda la información del campo, y no solo una parte.

% -------------------------------------------------------------------------
\subsection{Ecuaciones de Poisson y de Laplace}\label{sec:poisson}
% -------------------------------------------------------------------------
¿Qué debemos conocer para calcular $\phi$? Sustituimos $\vE=-\nab\phi$ en las ecuaciones del campo.

\begin{teorema}[Ecuaciones de Poisson y de Laplace]\label{teo:poisson}
El potencial electrostático satisface la \emph{ecuación de Poisson}
\begin{equation}\label{eq:poisson}
  \boxed{\;\Lap\phi=-\frac{\rho}{\eo}\;}
\end{equation}
y, en recintos sin fuentes ($\rho=0$), la \emph{ecuación de Laplace}
\begin{equation}\label{eq:laplace}
  \Lap\phi=0 .
\end{equation}
\end{teorema}
\begin{proof}
Con $\vE=-\nab\phi$ en la ley de Gauss diferencial,
$\nab\cdot\vE=\nab\cdot\bigl(-\nab\phi\bigr)=-\Lap\phi=\rho/\eo$. En una región con $\rho=0$ se obtiene \cref{eq:laplace}.
\end{proof}

Si en lugar de la divergencia sustituimos en el rotacional, $\nab\times\vE=\nab\times(-\nab\phi)=\vc{0}$ se cumple idénticamente: la irrotacionalidad
\emph{no aporta} ninguna ecuación adicional para $\phi$; basta una sola ecuación, la de Poisson. Es una ventaja práctica: para los campos vectoriales
hacen falta dos ecuaciones (divergencia y rotacional), pero para los escalares solo una, y al ser $\phi$ una función escalar resulta más sencillo resolverla y
calcular después $\vE=-\nab\phi$. Para una carga puntual, $\Lap\phi=\frac{\qp}{4\pi\eo}\Lap(1/R)=-\frac{\qp}{\eo}\delta^{3}(\vc{R})$, coherente con
el \cref{teo:delta}.

% -------------------------------------------------------------------------
\subsection{Líneas de campo y superficies equipotenciales}\label{sec:lineas}
% -------------------------------------------------------------------------
El campo eléctrico se suele representar, en cada punto, mediante un vector; globalmente, mediante sus líneas de fuerza.

\begin{definicion}[Líneas de campo y superficies equipotenciales]\label{def:lineas}
Las \emph{líneas de fuerza} (o \emph{de campo}) son las líneas que en cada punto son tangentes al vector $\vE$ en ese punto. Las
\emph{superficies equipotenciales} son aquellas en las que $\phi(\vc{r})$ es constante.
\end{definicion}

\begin{proposicion}[Propiedades de las líneas de campo y de las equipotenciales]\label{prop:lineas}
\begin{enumerate}[label=(\alph*)]
  \item Las líneas de fuerza \emph{no se cortan}: si lo hicieran, $\vE$ tendría dos direcciones en un punto y no estaría unívocamente definido.
  \item Salen de las cargas positivas y terminan en las negativas (son líneas abiertas), y se dibujan de modo que su densidad sea
        proporcional a $\norm{\vE}$: el campo es más intenso donde están más juntas.
  \item Las superficies equipotenciales son perpendiculares en cada punto a las líneas de fuerza.
  \item Las líneas de fuerza apuntan siempre hacia donde el potencial disminuye.
\end{enumerate}
\end{proposicion}
\begin{proof}
(c) De $\phi(\vc{r})=-\int\vE\cdot\dif\vc{r}+\text{cte}$ se tiene $\dif\phi=-\vE\cdot\dif\vc{r}$. Si nos movemos sobre una superficie equipotencial,
$\dif\phi=0$, luego $\vE\perp\dif\vc{r}$ para todo desplazamiento tangente a la superficie. (d) Si nos desplazamos en el sentido del campo,
$\dif\vc{r}\parallel\vE$ y $\dif\phi=-\norm{\vE}\,\dif r<0$.
\end{proof}

La \cref{fig:lineas} muestra las líneas de campo y las equipotenciales de una carga puntual, de un dipolo y de dos cargas iguales.

% Figura: líneas de campo y superficies equipotenciales (corte meridiano).
\begin{figure}[ht]
\centering
\IfFileExists{rapido.flag}{\fbox{\parbox{0.9\textwidth}{\centering\small Figura omitida (modo rápido)}}}{%
\newcommand{\marcoL}{\draw[grispizarra, line width=0.5pt] (-3,-3) rectangle (3,3);}
\tikzset{lcampo/.style={cm1, line width=0.55pt, postaction={decorate}, decoration={markings, mark=at position 0.5 with {\arrow{Stealth[length=3.6pt]}}}},
         lequi/.style={azulacero, line width=0.5pt}}
\begin{minipage}[t]{0.32\textwidth}\centering
\begin{tikzpicture}[x=0.8cm, y=0.8cm, baseline=(current bounding box.north)]
  \begin{scope}
    \clip (-3,-3) rectangle (3,3);
    \foreach \r in {0.5,1.0,1.5,2.0,2.5,3.0}{\draw[lequi] (0,0) circle (\r);}
    \foreach \a in {0,22.5,...,337.5}{%
      \draw[lcampo] (\a:0.2) -- (\a:3.3);}
  \end{scope}
  \marcoL
  \node[cargapos] at (0,0) {$+$};
\end{tikzpicture}\\[2pt]{\small\scshape (a)\; carga puntual}
\end{minipage}\hfill
\begin{minipage}[t]{0.32\textwidth}\centering
\begin{tikzpicture}[x=0.8cm, y=0.8cm, baseline=(current bounding box.north)]
  \def\qs{-1}
  \begin{scope}
    \clip (-3,-3) rectangle (3,3);
    % equipotenciales (positivas: línea continua; negativas, por simetría: discontinua)
    \foreach \xs in {-1.25,-1.4,-1.65,-2.0,-2.6}{
      \equipC{\xs}{0}{0.05}{140}{-1}
      \draw[lequi] \ruta; \begin{scope}[yscale=-1]\draw[lequi] \ruta;\end{scope}
      \begin{scope}[xscale=-1]\draw[lequi, densely dashed] \ruta;\end{scope}
      \begin{scope}[xscale=-1,yscale=-1]\draw[lequi, densely dashed] \ruta;\end{scope}}
    \draw[lequi] (0,-3) -- (0,3);
    % líneas de campo (8 conos de igual flujo desde la carga positiva)
    \foreach \k in {1,...,8}{
      \pgfmathsetmacro{\psi}{acos(1-(2*\k-1)/8)}
      \pgfmathsetmacro{\sx}{-1+0.13*cos(\psi)}\pgfmathsetmacro{\sy}{0.13*sin(\psi)}
      \lineaC{\sx}{\sy}{0.05}{200}{1}
      \draw[lcampo] \ruta;
      \begin{scope}[yscale=-1]\draw[lcampo] \ruta;\end{scope}}
  \end{scope}
  \marcoL
  \node[cargapos] at (-1,0) {$+$};
  \node[carganeg] at (1,0) {$-$};
\end{tikzpicture}\\[2pt]{\small\scshape (b)\; dipolo}
\end{minipage}\hfill
\begin{minipage}[t]{0.32\textwidth}\centering
\begin{tikzpicture}[x=0.8cm, y=0.8cm, baseline=(current bounding box.north)]
  \def\qs{1}
  \begin{scope}
    \clip (-3,-3) rectangle (3,3);
    \foreach \xs in {-1.2,-1.3,-1.45,-1.8,-2.2,-2.7}{
      \equipC{\xs}{0}{0.05}{150}{-1}
      \foreach \sxx/\syy in {1/1,1/-1,-1/1,-1/-1}{\begin{scope}[xscale=\sxx,yscale=\syy]\draw[lequi] \ruta;\end{scope}}}
    \foreach \k in {1,...,8}{
      \pgfmathsetmacro{\psi}{acos(1-(2*\k-1)/8)}
      \pgfmathsetmacro{\sx}{-1+0.13*cos(\psi)}\pgfmathsetmacro{\sy}{0.13*sin(\psi)}
      \lineaC{\sx}{\sy}{0.05}{200}{1}
      \foreach \sxx/\syy in {1/1,1/-1,-1/1,-1/-1}{\begin{scope}[xscale=\sxx,yscale=\syy]\draw[lcampo] \ruta;\end{scope}}}
  \end{scope}
  \marcoL
  \node[cargapos] at (-1,0) {$+$};
  \node[cargapos] at (1,0) {$+$};
\end{tikzpicture}\\[2pt]{\small\scshape (c)\; dos cargas iguales}
\end{minipage}
}%
\caption{Líneas de campo (púrpura, con el sentido del campo) y superficies equipotenciales (azul, vistas en
un plano meridiano) de (a)~una carga puntual, (b)~un dipolo y (c)~dos cargas iguales. Las líneas se trazan con
igual flujo entre dos consecutivas, de modo que su densidad es proporcional a $\norm{\vE}$. Salen de las cargas
positivas y llegan a las negativas, nunca se cortan y son perpendiculares a las equipotenciales, hacia donde
el potencial disminuye. En (b) las equipotenciales de potencial negativo se trazan a trazos y la mediatriz es la
de potencial nulo; en (c) hay un punto de campo nulo en el punto medio.}
\label{fig:lineas}
\end{figure}


% =========================================================================
\section{Condiciones de contorno}\label{sec:contorno}
% =========================================================================
Al estudiar distintas situaciones se observa que el campo eléctrico sufre una discontinuidad al atravesar una superficie cargada con densidad $\sigma$
(un plano, una cáscara esférica). Interesa saber qué ocurre exactamente en esa frontera. Es de esperar que, en un caso real, ninguna propiedad
electromagnética cambie de forma abrupta, sino continuamente a lo largo de una \emph{capa de transición} de anchura $h$: los campos físicos son continuos y con
derivadas continuas. Pero no podemos esperar conocer los detalles de esa transición, y por eso se sustituye la situación real por una idealizada, en la que
$h\to0$ (\cref{fig:capa}). Como consecuencia, los campos pueden mostrar una discontinuidad al pasar la frontera.

% Figura: capa de transición real frente a la discontinuidad idealizada.
\begin{figure}[ht]
\centering
\pgfplotsset{capaax/.style={tema, width=0.44\textwidth, height=3.8cm, xmin=-1.6, xmax=1.6, ymin=-0.1, ymax=1.2,
   xtick=\empty, ytick=\empty, grid=none, xlabel={$x$}, ylabel={$\vc{E}$}, axis lines=left, clip=false}}
\begin{tikzpicture}
  \begin{axis}[capaax, name=R, title={\small\scshape (a)\; real}]
    \fill[azulfondo] (axis cs:-0.4,-0.1) rectangle (axis cs:0.4,1.2);
    \draw[grispizarra, densely dashed, line width=0.4pt] (axis cs:-0.4,-0.1) -- (axis cs:-0.4,1.2);
    \draw[grispizarra, densely dashed, line width=0.4pt] (axis cs:0.4,-0.1) -- (axis cs:0.4,1.2);
    \addplot[azulprof, line width=1.2pt, domain=-1.6:1.6, samples=100] {0.15+0.85/(1+exp(-9*x))};
    \draw[<->, grispizarra, line width=0.4pt] (axis cs:-0.4,-0.02) -- (axis cs:0.4,-0.02);
    \node[font=\scriptsize, text=grispizarra, anchor=north] at (axis cs:0,-0.04) {$h$};
    \node[font=\scriptsize, text=azulprof, align=center] at (axis cs:0,0.55) {capa de\\transición};
  \end{axis}
  \begin{axis}[capaax, at={(R.east)}, anchor=west, xshift=0.8cm, title={\small\scshape (b)\; ideal ($h\to0$)}]
    \draw[grispizarra, densely dashed, line width=0.4pt] (axis cs:0,-0.1) -- (axis cs:0,1.2);
    \addplot[azulprof, line width=1.2pt, const plot, mark=none] coordinates {(-1.6,0.15) (0,0.15) (0,1.0) (1.6,1.0)};
    \node[font=\scriptsize, text=azulprof, anchor=west, align=left] at (axis cs:0.05,0.55) {superficie de\\[-2pt]discontinuidad};
  \end{axis}
\end{tikzpicture}
\caption{Un campo físico varía de forma continua (y con derivadas continuas) a través de una capa
de transición de anchura $h$ (a); como no se conocen los detalles de esa transición, se sustituye por una
superficie de discontinuidad, el límite $h\to0$ (b).}
\label{fig:capa}
\end{figure}


Lo analizamos primero para un campo vectorial genérico $\vc{F}$ y luego particularizamos a $\vE$ y a $\phi$.

% -------------------------------------------------------------------------
\subsection{Hipótesis de regularidad}\label{sec:hipotesis}
% -------------------------------------------------------------------------
\begin{definicion}[Campo con discontinuidad en una superficie]\label{def:salto}
Sea $S$ una superficie que separa dos regiones $V_1$ y $V_2$, y sea $\un$ la normal a $S$ que apunta del medio~1 al~2. Decimos que el campo
$\vc{F}:\R^{3}\to\R^{3}$ \emph{presenta una discontinuidad de salto en $S$} si $\vc{F}\in C^{1}(V_1)\cap C^{1}(V_2)$ y admite límites laterales finitos al
aproximarse a $S$,
\[
  \vc{F}_1(\vc{r})=\lim_{\vc{r}\to S^{-}}\vc{F},\qquad\vc{F}_2(\vc{r})=\lim_{\vc{r}\to S^{+}}\vc{F},\qquad
  \norm{\vc{F}_1}<\infty,\quad\norm{\vc{F}_2}<\infty ,
\]
para los puntos de $S$.
\end{definicion}
En términos intuitivos: $\vc{F}$ puede ser discontinuo al atravesar $S$, pero no puede divergir en $S$; debe permanecer acotado en un entorno de la superficie
y ser regular a cada lado.

\begin{observacion}[Qué permite esta hipótesis]\label{obs:hipotesis}
En la condición que obtendremos con el teorema de Gauss, la superficie lateral del cilindro tiene área $O(h)$ y, como $\vc{F}=O(1)$, su flujo es $O(h)\to0$. En la que se
obtiene con Stokes, los dos lados verticales del circuito tienen longitud $O(h)$ y por tanto $\int\vc{F}\cdot\dl=O(h)\to0$. La misma hipótesis justifica, entonces, que despreciemos
estas partes de la superficie y del contorno.
\end{observacion}

% -------------------------------------------------------------------------
\subsection{Condiciones de contorno de un campo vectorial}\label{sec:contorno-general}
% -------------------------------------------------------------------------
Sea $\vc{F}$ un campo del que conocemos sus fuentes,
\begin{equation}\label{eq:fuentes}
  \nab\cdot\vc{F}(\vc{r})=b(\vc{r}),\qquad\nab\times\vc{F}(\vc{r})=\vc{c}(\vc{r}).
\end{equation}
Estas ecuaciones no se pueden aplicar directamente a la discontinuidad. Las usaremos en la capa de transición, donde todo es continuo y con derivada continua, y después
tomaremos límites cuando $h\to0$.

\begin{teorema}[Condiciones de contorno generales]\label{teo:contorno}
Bajo las hipótesis de la \cref{def:salto}, con $\nab\cdot\vc{F}=b$ y $\nab\times\vc{F}=\vc{c}$ en la capa de transición, los saltos de las componentes normal y
tangencial de $\vc{F}$ al cruzar $S$ son
\begin{equation}\label{eq:contorno-normal}
  \boxed{\;F_{2n}-F_{1n}=(\vc{F}_2-\vc{F}_1)\cdot\un=\lim_{h\to0}\bigl(h\,\nab\cdot\vc{F}\bigr)=\lim_{h\to0}(h\,b)\;}
\end{equation}
\begin{equation}\label{eq:contorno-tangencial}
  \boxed{\;\un\times(\vc{F}_2-\vc{F}_1)=\lim_{h\to0}h\,\vc{c}\quad\Longleftrightarrow\quad
  \vc{F}_{2t}-\vc{F}_{1t}=\lim_{h\to0}h\,\bigl(\nab\times\vc{F}\bigr)\times\un\;}
\end{equation}
\end{teorema}
\begin{proof}
\emph{Componente normal.} Escogemos un cilindro de altura $h$ construido en la capa de transición, con bases de área $\Delta S$ paralelas a la superficie
(\figpanel{fig:contorno}{a}) y normales $\un_2=\un$ y $\un_1=-\un$. Integrando $\nab\cdot\vc{F}$ en el volumen $V$ del cilindro y usando el teorema de la divergencia,
\[
  \iiint_{V}\nab\cdot\vc{F}\,\dV=\oiint_{S_V}\vc{F}\cdot\dS=\iiint_{V}b\,\dV .
\]
Si el cilindro es suficientemente pequeño como para que $\vc{F}$ sea constante en sus bases y $b$ constante en todo el volumen,
\[
  \oiint_{S_V}\vc{F}\cdot\dS=\vc{F}_2\cdot\un\,\Delta S+\vc{F}_1\cdot(-\un)\,\Delta S+\iint_{S_{\mathrm{lat}}}\vc{F}\cdot\dS=b\,\Delta S\,h .
\]
Tomando límites con $h\to0$ la integral lateral se anula (\cref{obs:hipotesis}); sin embargo, como $b$ puede tener un valor muy alto en la discontinuidad
(una singularidad concentrada en la superficie), el producto $b\,\Delta S\,h$ puede ser distinto de cero. Dividiendo por $\Delta S$ se obtiene \cref{eq:contorno-normal}.

\emph{Componente tangencial.} Escogemos un pequeño circuito rectangular $C$ en un plano perpendicular a la superficie (\figpanel{fig:contorno}{b}), con dos lados de longitud $\Delta l$
paralelos a $S$, uno a cada lado de ella, y otros dos de longitud $h$ en la capa de transición. Sea $\un'$ la normal al rectángulo (que sale del papel) y $\vc{t}=\un'\times\un$ el
versor tangente a la superficie en el plano del circuito, de modo que $\Delta\vc{l}_2=\Delta l\,\vc{t}$ y $\Delta\vc{l}_1=-\Delta l\,\vc{t}$. Por el teorema de Stokes,
\[
  \iint_{S_C}(\nab\times\vc{F})\cdot\dS=\oint_{C}\vc{F}\cdot\dl=\iint_{S_C}\vc{c}\cdot\dS,
\]
y suponiendo el rectángulo suficientemente pequeño para tomar $\vc{F}$ y $\vc{c}$ constantes,
\[
  \oint_{C}\vc{F}\cdot\dl=\vc{F}_2\cdot\Delta\vc{l}_2+\vc{F}_1\cdot\Delta\vc{l}_1+\int_{\text{lados verticales}}\vc{F}\cdot\dl=\vc{c}\cdot\un'\,h\,\Delta l .
\]
Los lados verticales miden $O(h)$ y contribuyen $O(h)\to0$. Tomando el límite $h\to0$ queda
$(\vc{F}_2-\vc{F}_1)\cdot\vc{t}\,\Delta l=\lim\vc{c}\cdot\un'\,h\,\Delta l$, es decir,
$(\vc{F}_2-\vc{F}_1)\cdot(\un'\times\un)=\lim\vc{c}\cdot\un'\,h$. Con la propiedad del producto mixto
$\vc{A}\cdot(\vc{B}\times\vc{C})=\vc{B}\cdot(\vc{C}\times\vc{A})$, el primer miembro es
$\un'\cdot\bigl[\un\times(\vc{F}_2-\vc{F}_1)\bigr]$, y sacando factor común a $\un'$,
\[
  \un'\cdot\Bigl[\un\times(\vc{F}_2-\vc{F}_1)-\lim_{h\to0}h\,\vc{c}\Bigr]=0 .
\]
Como el circuito puede orientarse con cualquier $\un'$ tangente a la superficie, se obtiene \cref{eq:contorno-tangencial}. Multiplicando vectorialmente por $\un$ (con
$\un\times[\un\times\vc{u}]=-\vc{u}_t$) se obtiene la forma equivalente para las componentes tangenciales.
\end{proof}

% -------------------------------------------------------------------------
\subsection{Condiciones de contorno del campo y del potencial electrostáticos}\label{sec:contorno-electrico}
% -------------------------------------------------------------------------
\begin{corolario}[Condiciones de contorno de $\vE$]\label{cor:contorno-E}
Sea $\sigma$ la densidad superficial de carga en la frontera, entendida como el límite $\sigma=\lim_{h\to0}(\rho\,h)$ de una densidad volúmica
muy grande en la capa de transición. Entonces
\begin{equation}\label{eq:contorno-E}
  \boxed{\;\un\cdot(\vE_2-\vE_1)=\lim_{h\to0}\Bigl(h\,\frac{\rho}{\eo}\Bigr)=\frac{\sigma}{\eo},\qquad
  \un\times(\vE_2-\vE_1)=\vc{0}\;}
\end{equation}
Es decir: las componentes normales a la superficie de separación son \emph{discontinuas}, con un salto $\sigma/\eo$, si existe densidad superficial de carga en la
frontera, mientras que las componentes tangenciales \emph{siempre se conservan}: $\vE_{2t}=\vE_{1t}$.
\end{corolario}
\begin{proof}
Basta tomar $b=\rho/\eo$ y $\vc{c}=\vc{0}$ (porque $\nab\times\vE=\vc{0}$) en el \cref{teo:contorno}; el límite del segundo miembro de la ecuación del rotacional es $\lim h\cdot0=0$.
\end{proof}

\begin{corolario}[Continuidad del potencial]\label{cor:contorno-phi}
El potencial electrostático es continuo al atravesar cualquier superficie de separación: $\phi_2=\phi_1$.
\end{corolario}
\begin{proof}
Sean $\phi_1$ y $\phi_2$ los valores del potencial a cada lado de la capa de transición. La integral puede hacerse sobre cualquier trayectoria a través de la capa
(\cref{eq:int-linea}); tomamos una perpendicular a la superficie (\figpanel{fig:contorno}{c}):
\[
  \phi_2-\phi_1=-\int_{1}^{2}\vE\cdot\dl=-\int_{1}^{2}E_n\,\dif l=-\langle E_n\rangle\,h,
\]
donde $\langle E_n\rangle$ es el valor medio de la componente normal de $\vE$ en la capa de anchura $h$. Tomando el límite,
$\phi_2-\phi_1=-\lim_{h\to0}\langle E_n\rangle\,h=0$, puesto que el campo eléctrico nunca es infinito (a lo sumo $E_n$ tiene una discontinuidad, pero sigue siendo finito).
\end{proof}

% Figura: construcciones para las condiciones de contorno (cilindro, circuito, camino del potencial).
\begin{figure}[ht]
\centering
\tikzset{medio/.style={font=\small, circle, draw=grispizarra, line width=0.4pt, inner sep=1.2pt, fill=white},
         lamina/.style={fill=grisclaro!60, draw=grispizarra!70, line width=0.4pt}}
% (a) cilindro de altura h
\begin{minipage}[t]{0.98\textwidth}\centering
\begin{tikzpicture}[baseline=(current bounding box.north), x=1cm, y=1cm]
  % vista oblicua
  \fill[lamina] (0,0) -- (3.4,0) -- (4.3,0.9) -- (0.9,0.9) -- cycle;
  \draw[vect, azulprof] (0.55,0.45) -- ++(0,1.5) node[above, font=\small] {$\un$};
  \node[medio] at (0.3,0.12) {2}; \node[medio] at (0.3,-0.35) {1};
  \draw[rojoteo, line width=0.7pt] (2.0,0.2) -- (2.0,0.95); \draw[rojoteo, line width=0.7pt] (2.9,0.2) -- (2.9,0.95);
  \draw[rojoteo, line width=0.7pt, fill=white, fill opacity=0.6] (2.45,0.95) ellipse (0.45 and 0.15);
  \draw[rojoteo, line width=0.7pt] (2.0,0.2) arc (180:360:0.45 and 0.15);
  \draw[rojoteo, line width=0.4pt, densely dashed] (2.0,0.2) arc (180:0:0.45 and 0.15);
  \node[font=\scriptsize, text=rojoteo, anchor=west] at (3.0,0.6) {$h$};
  % corte transversal
  \begin{scope}[shift={(6.3,0.2)}]
    \fill[lamina] (-0.2,-0.04) rectangle (3.6,0.04);
    \draw[grispizarra, line width=0.5pt] (-0.2,0) -- (3.6,0);
    \draw[azulprof, ->, line width=0.8pt] (-0.1,-0.5) -- (-0.1,1.0) node[above, font=\small] {$\un$};
    \node[medio] at (-0.45,0.35) {2}; \node[medio] at (-0.45,-0.35) {1};
    \draw[rojoteo, line width=0.9pt] (1.2,-0.6) rectangle (2.4,0.6);
    \draw[->, line width=0.9pt, azulacero] (1.8,0.6) -- ++(0,0.55) node[right, font=\scriptsize] {$\un_2=\un$};
    \draw[->, line width=0.9pt, azulacero] (1.8,-0.6) -- ++(0,-0.55) node[right, font=\scriptsize] {$\un_1=-\un$};
    \draw[<->, rojoteo, line width=0.5pt] (2.7,-0.6) -- (2.7,0.6) node[midway, right, font=\scriptsize] {$h$};
    \draw[<->, rojoteo, line width=0.5pt] (1.2,-0.85) -- (2.4,-0.85) node[midway, below, font=\scriptsize] {$\Delta S$};
  \end{scope}
\end{tikzpicture}\\[2pt]{\small\scshape (a)\; componente normal: cilindro de altura $h$ en la capa de transición}
\end{minipage}\\[10pt]
% (b) circuito
\begin{minipage}[t]{0.98\textwidth}\centering
\begin{tikzpicture}[baseline=(current bounding box.north), x=1cm, y=1cm]
  \fill[lamina] (0,0) -- (3.4,0) -- (4.3,0.9) -- (0.9,0.9) -- cycle;
  \draw[azulprof, ->, line width=0.8pt] (0.55,0.45) -- ++(0,1.5) node[above, font=\small] {$\un$};
  \node[medio] at (0.3,0.12) {2}; \node[medio] at (0.3,-0.35) {1};
  \draw[rojoteo, line width=0.8pt] (1.8,-0.3) rectangle (3.3,1.1);
  \begin{scope}[shift={(6.3,0.2)}]
    \fill[lamina] (-0.2,-0.04) rectangle (3.6,0.04);
    \draw[azulprof, ->, line width=0.8pt] (-0.1,-0.5) -- (-0.1,1.0) node[above, font=\small] {$\un$};
    \node[medio] at (-0.45,0.35) {2}; \node[medio] at (-0.45,-0.35) {1};
    \draw[rojoteo, line width=0.9pt, orienta] (1.0,-0.6) -- (1.0,0.6) -- (2.6,0.6) -- (2.6,-0.6) -- cycle;
    \draw[->, azulacero, line width=0.9pt] (1.8,0.6) -- ++(0.7,0) node[above, font=\scriptsize] {$\vc{t}$};
    \draw[->, azulacero, line width=0.9pt] (1.8,-0.6) -- ++(-0.7,0);
    \node[font=\scriptsize, text=azulacero, anchor=north] at (1.1,-0.64) {$\Delta\vc{l}_1=-\Delta l\,\vc{t}$};
    \node[font=\scriptsize, text=azulacero, anchor=south west] at (1.4,0.66) {$\Delta\vc{l}_2=\Delta l\,\vc{t}$};
    \node[font=\scriptsize, text=rojoteo, anchor=west] at (2.65,0) {$h$};
    \draw[azulprof, line width=0.7pt] (1.8,0) circle (0.12);
    \fill[azulprof] (1.8,0) circle (0.9pt);
    \node[font=\scriptsize, text=azulprof, anchor=north west] at (1.85,-0.08) {$\un'$};
  \end{scope}
\end{tikzpicture}\\[2pt]{\small\scshape (b)\; componente tangencial: circuito rectangular $C$ ($\un'$ sale de la figura)}
\end{minipage}\\[10pt]
% (c) camino del potencial
\begin{minipage}[t]{0.98\textwidth}\centering
\begin{tikzpicture}[baseline=(current bounding box.north), x=1cm, y=1cm]
  \fill[lamina] (0,0) -- (3.4,0) -- (4.3,0.9) -- (0.9,0.9) -- cycle;
  \draw[azulprof, ->, line width=0.8pt] (0.55,0.45) -- ++(0,1.5) node[above, font=\small] {$\un$};
  \node[medio] at (0.3,0.12) {2}; \node[medio] at (0.3,-0.35) {1};
  \draw[rojoteo, line width=0.9pt] (2.2,-0.55) -- (2.2,0.55);
  \fill[rojoteo] (2.2,0.55) circle (1.6pt) node[above right, font=\small] {$B$};
  \fill[rojoteo] (2.2,-0.55) circle (1.6pt) node[below right, font=\small] {$A$};
  \begin{scope}[shift={(6.3,0.2)}]
    \fill[lamina] (-0.2,-0.04) rectangle (3.6,0.04);
    \draw[azulprof, ->, line width=0.8pt] (-0.1,-0.5) -- (-0.1,1.0) node[above, font=\small] {$\un$};
    \node[medio] at (-0.45,0.35) {2}; \node[medio] at (-0.45,-0.35) {1};
    \draw[rojoteo, line width=0.9pt, orienta] (1.8,-0.6) -- (1.8,0.6);
    \fill[rojoteo] (1.8,0.6) circle (1.6pt) node[above right, font=\small] {$B$};
    \fill[rojoteo] (1.8,-0.6) circle (1.6pt) node[below right, font=\small] {$A$};
    \draw[<->, rojoteo, line width=0.5pt] (2.5,-0.6) -- (2.5,0.6) node[midway, right, font=\scriptsize] {$h$};
  \end{scope}
\end{tikzpicture}\\[2pt]{\small\scshape (c)\; potencial: camino $A\to B$ perpendicular a la superficie}
\end{minipage}
\caption{Construcciones geométricas para obtener las condiciones de contorno (la región 1 está
debajo de la superficie y la 2 encima; $\un$ apunta de 1 a 2). (a)~Un cilindro de bases de área $\Delta S$ paralelas a la
superficie y altura $h$, contenido en la capa de transición, para integrar la divergencia. (b)~Un
rectángulo, con dos lados de longitud $\Delta l$ paralelos a la superficie y dos de longitud $h$ (los
dos lados verticales), para integrar el rotacional. (c)~Un camino perpendicular a la
superficie, de $A$ a $B$, para el potencial.}
\label{fig:contorno}
\end{figure}


\begin{observacion}[Consecuencias y comprobaciones]\label{obs:contorno-ejemplos}
Las condiciones se pueden reformular con el potencial: como $\vE=-\nab\phi$, la continuidad de $\vE_t$ equivale a la de $\phi$ (que ya hemos demostrado) y el salto de la
derivada normal es $\partial_n\phi|_2-\partial_n\phi|_1=-\sigma/\eo$. Algunos ejemplos que se comprueban en los seminarios:
\begin{itemize}
  \item \emph{Plano infinito} (\cref{sem:plano}): $\vE=\frac{\sigma}{2\eo}\,\mathrm{sgn}(z)\,\uz$, de modo que el salto en $z=0$ es $\frac{\sigma}{2\eo}-\bigl(-\frac{\sigma}{2\eo}\bigr)=\frac{\sigma}{\eo}$.
  \item \emph{Cáscara esférica} (\cref{sem:cascara}): $E=0$ para $r<a$ y $E=\frac{Q}{4\pi\eo a^{2}}=\frac{\sigma}{\eo}$ para $r=a^{+}$, mientras que $\phi$ es continuo.
  \item \emph{Lámina con densidad volúmica} (\cref{sem:lamina}): no hay carga superficial ($\sigma=0$) y por eso $\vE$ es \emph{continuo} en las caras de la lámina.
\end{itemize}
\end{observacion}

% =========================================================================
\section{Seminarios prácticos}\label{sec:seminarios}
% =========================================================================
Esta sección resuelve, con todo detalle, los doce problemas de la relación del tema~\cite{PROBLEMAS-T1}. Cada uno sigue la misma estructura: el
\emph{enunciado}, la \emph{idea} que guía la resolución (qué estamos buscando y por qué el método funciona), la \emph{resolución} paso a paso
y, cuando merece la pena, un breve comentario sobre la \emph{interpretación} física del resultado. Los problemas son de dos familias:
\begin{itemize}
  \item \emph{Integración directa} (problemas 1.1, 1.3, 1.5, 1.7, 1.12): se aplica la definición del campo (\cref{eq:campo-vol,eq:campo-sup-lin}) o del potencial
        (\cref{eq:pot-vol,eq:pot-sup-lin}). Funciona siempre, pero exige elegir bien el sistema de coordenadas y aprovechar la simetría para anular componentes.
  \item \emph{Ley de Gauss} (problemas 1.2, 1.4, 1.6, 1.8, 1.9, 1.10): se elige una superficie gaussiana adaptada a la simetría (\cref{sec:gauss-simetria}) y se despeja $E$ de
        $\oiint\vE\cdot\dS=\Qenc/\eo$. Es mucho más corta, pero solo posible cuando la simetría es esférica, cilíndrica o plana. El potencial se obtiene después
        integrando el campo, $\phi(\vc{r})-\phi(\vc{r}_0)=-\int_{\vc{r}_0}^{\vc{r}}\vE\cdot\dl$.
\end{itemize}
El problema 1.11 se resuelve por superposición de los campos de tres planos, calculados antes. La \cref{tab:resultados} recoge, al final de la sección, todos los resultados.

% =========================================================================
\subsection{Problema 1.1. Esfera uniformemente cargada: campo exterior por integración directa}\label{sem:esfera-directa}
% =========================================================================
\begin{ejercicio}[Campo exterior de una esfera, por integración directa]\label{ejer:esfera-directa}
Calcular, por integración directa, el campo eléctrico creado por una carga $Q$ uniformemente distribuida en una esfera de radio $a$ en los puntos exteriores a dicha esfera.
\end{ejercicio}

\begin{resolucion}
\textbf{Idea.} La esfera es una distribución volúmica: cada elemento $\rho\,\dVp$ crea un campo coulombiano y hay que sumarlos todos (\cref{eq:campo-vol}). La simetría nos deja elegir el
punto campo donde más nos convenga: como el problema es invariante frente a rotaciones, basta calcular el campo \emph{en el eje $z$} (\cref{fig:sem-esfera}); el resultado en un punto
cualquiera se obtiene luego sustituyendo $z$ por la distancia $r$ y $\uz$ por $\ur$.

\textbf{Paso 1: geometría.} El punto campo es $\vc{r}=z\,\uz$ con $z>a$, y el punto fuente, en esféricas,
$\vc{r}\,'=r'\bigl[\cos\theta'\,\uz+\sen\theta'(\cos\varphi'\,\ux+\sen\varphi'\,\uy)\bigr]$. Entonces
\[
  \vc{R}=\vc{r}-\vc{r}\,'=(z-r'\cos\theta')\,\uz-r'\sen\theta'(\cos\varphi'\,\ux+\sen\varphi'\,\uy),\qquad
  R=\bigl(z^{2}+r'^{2}-2zr'\cos\theta'\bigr)^{1/2}.
\]
La densidad es constante, $\rho=Q/(\tfrac43\pi a^{3})$, y $\dVp=r'^{2}\sen\theta'\,\dif r'\,\dif\theta'\,\dif\varphi'$.

\textbf{Paso 2: la simetría anula dos componentes.} En la integral de $\vE$, las componentes $x$ e $y$ llevan $\cos\varphi'$ y $\sen\varphi'$, cuya integral sobre $[0,2\pi]$ es cero. Solo
sobrevive la componente $z$:
\[
  E_z=\frac{\rho}{4\pi\eo}\int_{0}^{a}\!\!\int_{0}^{\pi}\!\!\int_{0}^{2\pi}\frac{(z-r'\cos\theta')\,r'^{2}\sen\theta'}{\bigl(z^{2}+r'^{2}-2zr'\cos\theta'\bigr)^{3/2}}\,\dif\varphi'\,\dif\theta'\,\dif r'
  =\frac{\rho}{2\eo}\int_{0}^{a}r'^{2}\,I(r')\,\dif r',
\]
donde hemos integrado en $\varphi'$ (que da $2\pi$) y llamado
\[
  I(r')=\int_{0}^{\pi}\frac{(z-r'\cos\theta')\,\sen\theta'}{\bigl(z^{2}+r'^{2}-2zr'\cos\theta'\bigr)^{3/2}}\,\dif\theta' .
\]

\textbf{Paso 3: la integral angular.} Con el cambio $\mu=\cos\theta'$ ($\dif\mu=-\sen\theta'\,\dif\theta'$) la integral se hace racional. Es más cómodo todavía tomar como variable la propia distancia
$R$: de $R^{2}=z^{2}+r'^{2}-2zr'\mu$ se tiene $\dif\mu=-R\,\dif R/(zr')$ y $z-r'\mu=(z^{2}-r'^{2}+R^{2})/(2z)$, de modo que
\[
  I(r')=\int_{|z-r'|}^{z+r'}\frac{z^{2}-r'^{2}+R^{2}}{2z\,R^{3}}\,\frac{R\,\dif R}{zr'}
  =\frac{1}{2z^{2}r'}\Bigl[-\frac{z^{2}-r'^{2}}{R}+R\Bigr]_{|z-r'|}^{z+r'} .
\]
(Es la integral $\int\frac{ax-b}{(a^{2}+b^{2}-2abx)^{3/2}}\dif x$ de la hoja de problemas.) Como el punto campo es exterior, $z>a\ge r'$, y por tanto $|z-r'|=z-r'$:
\[
  I(r')=\frac{1}{2z^{2}r'}\Bigl[\bigl(-(z-r')+(z+r')\bigr)-\bigl(-(z+r')+(z-r')\bigr)\Bigr]=\frac{1}{2z^{2}r'}\,(2r'+2r')=\frac{2}{z^{2}} .
\]

\textbf{Paso 4: la integral radial.} Queda
\[
  E_z=\frac{\rho}{2\eo}\int_{0}^{a}r'^{2}\,\frac{2}{z^{2}}\,\dif r'=\frac{\rho}{\eo z^{2}}\,\frac{a^{3}}{3}=\frac{\rho\,\frac43\pi a^{3}}{4\pi\eo z^{2}}=\frac{Q}{4\pi\eo z^{2}},
\]
y, restituyendo el vector,
\[
  \boxed{\;\vE(\vc{r})=\frac{Q}{4\pi\eo\,r^{2}}\,\ur\qquad(r>a)\;}
\]
\end{resolucion}

\begin{observacion}[Interpretación]
El campo exterior de una esfera uniforme es \emph{exactamente} el de una carga puntual $Q$ situada en su centro, un resultado análogo al teorema de Newton para la gravedad. Lo notable es
que la cancelación se produce a pesar de que los elementos más cercanos al punto campo pesan más: el factor $1/R^{2}$ hace que la contribución de los cercanos y de los lejanos se compense exactamente,
y esto solo ocurre para el exponente 2 (\cref{prop:cavendish}). Con la ley de Gauss (\cref{sem:esfera}) se llega al mismo resultado en tres líneas, pero aquí hemos comprobado que ambos métodos coinciden.
\end{observacion}

% Figura (Problema 1.1): esfera uniforme, punto campo en el eje z.
\begin{figure}[ht]
\centering
\tdplotsetmaincoords{64}{125}
\begin{tikzpicture}[tdplot_main_coords, scale=1.25]
  \def\RR{1.7}
  \draw[eje] (0,0,0) -- (0,3.2,0) node[right] {\small $y$};
  \draw[eje] (0,0,0) -- (3.1,0,0) node[below left] {\small $x$};
  \tdEsfera{\RR}{azulprof}
  \draw[eje] (0,0,0) -- (0,0,4.8) node[above] {\small $z$};
  \coordinate (P) at (0,0,3.5);
  \coordinate (S) at (0.2,1.35,0.6);
  % vectores
  \draw[->, azulclaro!75!azulacero, line width=0.9pt] (0,0,0) -- (S);
  \node[font=\small, text=azulprof, anchor=north] at ($(0,0,0)!0.55!(S)+(0,0,-0.08)$) {$\vc{r}\,'$};
  \draw[vect, azulprof, line width=1.2pt] (S) -- (P);
  \node[font=\small, text=azulprof, anchor=west] at ($(S)!0.55!(P)+(0.0,0.18,0)$) {$\vc{R}$};
  \draw[->, azulacero, line width=1.1pt] (0,0,0) -- (P);
  \node[font=\small, text=azulprof, anchor=east] at (0,0,2.9) {$\vc{r}=z\,\uz$};
  \draw[vectrojo, line width=1.5pt] (P) -- (0,0,4.3);
  \node[font=\small, text=acento, anchor=east] at (0,-0.08,4.2) {$\vE$};
  \fill[azulprof] (P) circle (1.5pt) node[right=2pt, font=\small] {$P$};
  % elemento de volumen
  \tdParche[acento, fill=acento!35, line width=0.6pt]{0.2,1.35,0.6}{0.13,-0.03,0}{0.03,0.12,0.1}
  \fill[acento] (S) circle (1.2pt);
  \node[font=\small, text=acento, anchor=north west] at ($(S)+(0.12,0.1,-0.05)$) {$\dV'$};
  % ángulo theta' entre el eje z y r' (arco en el plano que contiene a ambos)
  \draw[verdeteo, line width=0.7pt, -{Stealth[length=3.5pt]}, domain=0:62, samples=24, variable=\tth]
    plot ({0.85*sin(\tth)*cos(81.6)},{0.85*sin(\tth)*sin(81.6)},{0.85*cos(\tth)});
  \node[font=\small, text=verdeteo, fill=white, fill opacity=0.8, text opacity=1, inner sep=1pt]
    at ({1.2*sin(30)*cos(81.6)},{1.2*sin(30)*sin(81.6)},{1.2*cos(30)}) {$\theta'$};
  % radio a (en el plano de la pantalla)
  \begin{scope}[tdplot_screen_coords]
    \draw[grispizarra, arrows={Stealth[length=3.5pt]-Stealth[length=3.5pt]}, line width=0.5pt] (0,-0.05) -- (-\RR,-0.05);
    \node[font=\small, text=azulprof, anchor=north] at (-0.85,-0.08) {$a$};
    \node[font=\small, text=azulprof] at (0.55,-1.15) {$\rho$};
  \end{scope}
\end{tikzpicture}
\caption{Problema 1.1. Esfera de radio $a$ con densidad volúmica uniforme $\rho$. El punto campo $P$ está en el eje $z$ a distancia $z>a$ del centro,
$\vc{r}=z\,\uz$; el elemento $\dV'$ está en $\vc{r}\,'=r'\,\uh{r'}$ y $\vc{R}=\vc{r}-\vc{r}\,'$ lo une con $P$. Por simetría
solo sobrevive la componente $z$ del campo.}
\label{fig:sem-esfera}
\end{figure}


% =========================================================================
\subsection{Problema 1.2. Esfera uniformemente cargada: campo y potencial en todo el espacio}\label{sem:esfera}
% =========================================================================
\begin{ejercicio}[Campo y potencial de una esfera maciza]\label{ejer:esfera}
Calcular mediante el teorema de Gauss el campo y el potencial eléctrico para una distribución uniforme de carga $Q$ en una esfera de radio $a$ en todos los puntos del espacio.
Representar gráficamente cómo varían el potencial y el campo en función de la distancia al origen. Comprobar que se cumplen las ecuaciones de Maxwell estudiadas.
\end{ejercicio}

\begin{resolucion}
\textbf{Idea.} La distribución es esféricamente simétrica, así que $\vE=E(r)\,\ur$ (puramente radial y dependiente solo de la distancia al centro). Elegimos como superficie gaussiana
una esfera $S_{\mathrm{G}}$ concéntrica de radio $r$ (\figpanel{fig:sem-gauss-esfera}{a}). En ella $E(r)$ es constante y $\vE\parallel\dS=\dsc\,\ur$, de modo que
\[
  \oiint_{S_{\mathrm{G}}}\vE\cdot\dS=E(r)\oiint_{S_{\mathrm{G}}}\dsc=4\pi r^{2}\,E(r)=\frac{\Qenc}{\eo}.
\]
Distinguimos dos regiones, según la superficie gaussiana encierre solo una parte de la carga o toda ella. La densidad es $\rho=Q/(\frac43\pi a^{3})=3Q/(4\pi a^{3})$.

\textbf{Campo, caso I: $r<a$.} La esfera de radio $r$ encierra solo una parte, $\Qenc=\rho\,\tfrac43\pi r^{3}$. Entonces $4\pi r^{2}E=\frac{1}{\eo}\rho\frac43\pi r^{3}$ y
\[
  E(r)=\frac{\rho\,r}{3\eo}=\frac{Q\,r}{4\pi\eo\,a^{3}}\qquad(r<a).
\]
\textbf{Campo, caso II: $r\ge a$.} Ahora $\Qenc=Q$, y $E(r)=Q/(4\pi\eo r^{2})$. En conjunto,
\[
  \boxed{\;\vE(\vc{r})=\begin{cases}\dfrac{Q\,r}{4\pi\eo a^{3}}\,\ur,&0\le r<a,\\[8pt]\dfrac{Q}{4\pi\eo r^{2}}\,\ur,&r\ge a.\end{cases}\;}
\]
Ambas expresiones coinciden en $r=a$: $E(a)=Q/(4\pi\eo a^{2})$. Es lo esperable, porque no hay carga superficial (\cref{cor:contorno-E}).

\textbf{Potencial.} Integramos el campo a lo largo de una trayectoria radial, $\dl=\dif r\,\ur$, de modo que $\vE\cdot\dl=E(r)\,\dif r$. Con la referencia en el infinito, $\phi(\infty)=0$,
para $r\ge a$:
\[
  \phi(r)=-\int_{\infty}^{r}\frac{Q}{4\pi\eo r'^{2}}\,\dif r'=\frac{Q}{4\pi\eo}\Bigl[\frac{1}{r'}\Bigr]_{\infty}^{r}=\frac{Q}{4\pi\eo\,r},\qquad\text{en particular}\quad\phi(a)=\frac{Q}{4\pi\eo a}.
\]
Para $r<a$ \emph{no podemos} volver a tomar el infinito como origen de la integral, porque la expresión del campo cambia en $r=a$. Partimos del valor ya conocido $\phi(a)$:
\[
  \phi(r)-\phi(a)=-\int_{a}^{r}\frac{\rho\,r'}{3\eo}\,\dif r'=-\frac{\rho}{6\eo}\,(r^{2}-a^{2})
  \;\Longrightarrow\;
  \phi(r)=\frac{Q}{4\pi\eo a}+\frac{Q}{8\pi\eo a^{3}}(a^{2}-r^{2}),
\]
usando $\rho/6\eo=Q/(8\pi\eo a^{3})$. Agrupando términos,
\[
  \boxed{\;\phi(r)=\begin{cases}\dfrac{Q}{8\pi\eo a}\Bigl(3-\dfrac{r^{2}}{a^{2}}\Bigr),&0\le r\le a,\\[8pt]\dfrac{Q}{4\pi\eo\,r},&r\ge a.\end{cases}\;}
\]
Comprobación de continuidad: en $r=a$, $\frac{Q}{8\pi\eo a}(3-1)=\frac{Q}{4\pi\eo a}$, igual que la expresión exterior. Además, $\phi(0)=3Q/(8\pi\eo a)$, así que el potencial es máximo en el centro (con $Q>0$) mientras que
el campo allí es nulo, $\vE(0)=\vc{0}$.

\textbf{Gráficas} (\figpanel{fig:sem-gauss-esfera}{b) y~(c}). El campo crece linealmente hasta un máximo en la superficie y luego cae como $1/r^{2}$; el potencial decrece parabólicamente en el interior y como $1/r$ fuera. No hay que confundir la dependencia
$\phi\propto1/r$ del exterior con la $E\propto1/r^{2}$ del campo.
\end{resolucion}

\textbf{Tarea: comprobación de las ecuaciones de Maxwell.} Verifiquemos que los resultados anteriores son compatibles con $\nab\cdot\vE=\rho/\eo$, $\nab\times\vE=\vc{0}$ y $\Lap\phi=-\rho/\eo$.

\begin{resolucion}[Comprobación]
\textbf{1) Divergencia} (primera ecuación de Maxwell). Para un campo radial $\vE=E_r(r)\,\ur$, en esféricas $\nab\cdot\vE=\frac{1}{r^{2}}\frac{\dif}{\dif r}(r^{2}E_r)$.
\begin{itemize}
  \item \emph{Interior} ($r<a$): $r^{2}E_r=\frac{Q}{4\pi\eo a^{3}}\,r^{3}$, luego $\nab\cdot\vE=\frac{1}{r^{2}}\frac{3Q\,r^{2}}{4\pi\eo a^{3}}=\frac{3Q}{4\pi\eo a^{3}}=\frac{\rho}{\eo}>0$ (hay fuentes, el campo diverge).
  \item \emph{Exterior} ($r>a$): $r^{2}E_r=Q/4\pi\eo=\text{cte}$, luego $\nab\cdot\vE=0=\rho/\eo$, ya que fuera de la esfera no hay carga.
\end{itemize}
En $r=a$ debemos tratar aparte la posible contribución superficial; aquí no existe carga superficial ($\sigma\to0$), por lo que tanto $\vE$ como $\phi$ son continuos.

\textbf{2) Rotacional.} Un campo de la forma $\vE=E_r(r)\,\ur$ tiene $E_\theta=E_\varphi=0$ y $E_r$ solo depende de $r$ ($\partial_\theta E_r=\partial_\varphi E_r=0$); todas las componentes del rotacional en esféricas se anulan,
$\nab\times\vE=\vc{0}$ para todo $r$. Igualmente, $\vE=-\nab\phi$ implica $\nab\times\vE=-\nab\times(\nab\phi)=\vc{0}$.

\textbf{3) Ecuación de Poisson.} Como $\phi=\phi(r)$, $\Lap\phi=\frac{1}{r^{2}}\frac{\dif}{\dif r}\bigl(r^{2}\frac{\dif\phi}{\dif r}\bigr)$.
\begin{itemize}
  \item \emph{Interior}: $\phi'=-\frac{Q\,r}{4\pi\eo a^{3}}$, $r^{2}\phi'=-\frac{Q\,r^{3}}{4\pi\eo a^{3}}$, y $\Lap\phi=\frac{1}{r^{2}}\bigl(-\frac{3Q\,r^{2}}{4\pi\eo a^{3}}\bigr)=-\frac{3Q}{4\pi\eo a^{3}}=-\frac{\rho}{\eo}$.
  \item \emph{Exterior}: $\phi'=-\frac{Q}{4\pi\eo r^{2}}$, $r^{2}\phi'=\text{cte}$, luego $\Lap\phi=0$: se cumple la ecuación de Laplace, como corresponde a una región sin carga.
\end{itemize}
\end{resolucion}

% Figura (Problemas 1.2 y 1.9): superficies gaussianas y gráficas E(r), phi(r) para esfera maciza y cáscara.
\begin{figure}[ht]
\centering
\begin{minipage}[t]{0.30\textwidth}\centering
\begin{tikzpicture}[baseline=(current bounding box.north), scale=0.95]
  \shade[ball color=azulclaro!50, opacity=0.30] (0,0) circle (1.2);
  \draw[azulprof, line width=0.8pt] (0,0) circle (1.2);
  \draw[verdeteo, densely dashed, line width=0.7pt] (0,0) circle (0.7);
  \draw[verdeteo, densely dashed, line width=0.7pt] (0,0) circle (2.0);
  \foreach \a in {0,45,...,315}{\draw[-{Stealth[length=4pt]}, cm1, line width=0.6pt] (\a:2.0) -- (\a:2.45);}
  \foreach \a in {0,45,...,315}{\draw[-{Stealth[length=3pt]}, cm1, line width=0.6pt] (\a:0.7) -- (\a:0.95);}
  \foreach \p in {(-0.4,0.9),(0.9,0.3),(-0.95,-0.35),(0.3,-0.95)}{\node[font=\tiny, text=azulacero] at \p {$+$};}
  \draw[grispizarra, ->, line width=0.4pt] (0,0) -- (-125:1.2) node[pos=0.5, left, font=\scriptsize, text=azulprof] {$a$};
  \node[font=\scriptsize, text=verdeteo, anchor=south west] at (40:0.62) {$r$};
  \node[font=\scriptsize, text=verdeteo, anchor=north west] at (-45:2.05) {$r$};
  \node[font=\scriptsize, text=verdeteo] at (0,2.7) {$S_{\mathrm{G}}$};
\end{tikzpicture}\\[2pt]{\small\scshape (a)}
\end{minipage}\hfill
\begin{minipage}[t]{0.33\textwidth}\centering
\begin{tikzpicture}[baseline=(current bounding box.north)]
  \begin{axis}[tema, width=\textwidth, height=4.2cm, xlabel={$r/a$}, ylabel={$E/E(a)$}, xmin=0, xmax=3, ymin=0, ymax=1.15,
     xtick={0,1,2,3}, ytick={0,0.5,1}, domain=0:3, samples=120]
    \addplot[azulprof, line width=1.1pt, domain=0:1] {x};
    \addplot[azulprof, line width=1.1pt, domain=1:3] {1/x^2};
    \addplot[grispizarra, densely dashed] coordinates {(1,0) (1,1)};
    \node[font=\scriptsize, text=azulprof, anchor=south west] at (axis cs:0.15,0.35) {$\propto r$};
    \node[font=\scriptsize, text=azulprof, anchor=south west] at (axis cs:1.3,0.55) {$\propto 1/r^{2}$};
  \end{axis}
\end{tikzpicture}\\[2pt]{\small\scshape (b)}
\end{minipage}\hfill
\begin{minipage}[t]{0.33\textwidth}\centering
\begin{tikzpicture}[baseline=(current bounding box.north)]
  \begin{axis}[tema, width=\textwidth, height=4.2cm, xlabel={$r/a$}, ylabel={$\phi/\phi(a)$}, xmin=0, xmax=3, ymin=0, ymax=1.6,
     xtick={0,1,2,3}, ytick={0,0.5,1,1.5}, domain=0:3, samples=120]
    \addplot[azulacero, line width=1.1pt, domain=0:1] {(3-x^2)/2};
    \addplot[azulacero, line width=1.1pt, domain=1:3] {1/x};
    \addplot[grispizarra, densely dashed] coordinates {(1,0) (1,1)};
    \node[font=\scriptsize, text=azulacero, anchor=south west] at (axis cs:1.35,0.5) {$\propto 1/r$};
    \node[font=\scriptsize, text=azulacero, anchor=south west] at (axis cs:0.05,0.75) {$\tfrac{3}{2}-\tfrac{r^{2}}{2a^{2}}$};
  \end{axis}
\end{tikzpicture}\\[2pt]{\small\scshape (c)}
\end{minipage}
\caption{Problema 1.2. (a)~Corte de la esfera uniformemente cargada y de las dos superficies gaussianas esféricas (a trazos), una
interior ($r<a$) y otra exterior ($r>a$). (b)~El campo crece linealmente dentro y decae como $1/r^{2}$ fuera, con
valor máximo $E(a)$ en la superficie. (c)~El potencial, máximo en el centro, es parabólico dentro y decae como $1/r$ fuera; $E$ y $\phi$ son continuos en $r=a$.}
\label{fig:sem-gauss-esfera}
\end{figure}


% =========================================================================
\subsection{Problema 1.3. Hilo infinito: campo y potencial por integración directa}\label{sem:hilo}
% =========================================================================
\begin{ejercicio}[Hilo infinito, integración directa]\label{ejer:hilo}
Calcular por integración directa el campo y el potencial creados por un hilo infinito con una densidad lineal de carga uniforme $\lambda$.
\end{ejercicio}

\begin{resolucion}
\textbf{Idea.} Colocamos el hilo sobre el eje $z$ y usamos coordenadas cilíndricas. Por la simetría de traslación a lo largo del hilo y de rotación en torno a él, el campo solo puede depender de la distancia $r$ al hilo y
no puede tener componente en $z$: para cada elemento en $z'$ hay otro simétrico en $-z'$ cuya componente axial lo cancela (\figpanel{fig:sem-hilo}{a}).

\textbf{Campo.} El punto campo es $\vc{r}=r\,\urho$ y el fuente $\vc{r}\,'=z'\,\uz$, con $\dif q'=\lambda\,\dif z'$ ($\dlp=\dif z'$). Entonces
$\vc{R}=r\,\urho-z'\,\uz$, $R=(r^{2}+z'^{2})^{1/2}$ y, con la \cref{eq:campo-sup-lin},
\[
  \vE(\vc{r})=\frac{\lambda}{4\pi\eo}\int_{-\infty}^{\infty}\frac{r\,\urho-z'\,\uz}{(r^{2}+z'^{2})^{3/2}}\,\dif z' .
\]
La componente $z$ contiene el integrando impar $z'/(r^{2}+z'^{2})^{3/2}$ sobre un intervalo simétrico: es nula (\emph{la componente $z$ se anula por la simetría del problema}). Para la radial usamos la primitiva
$\int\frac{\dif z'}{(r^{2}+z'^{2})^{3/2}}=\frac{z'}{r^{2}\sqrt{r^{2}+z'^{2}}}$ (se comprueba derivando; es el caso $\int\frac{\dif x}{(x^{2}+c^{2})^{3/2}}=\frac{x}{c^{2}\sqrt{x^{2}+c^{2}}}$):
\[
  E_r=\frac{\lambda\,r}{4\pi\eo}\Bigl[\frac{z'}{r^{2}\sqrt{r^{2}+z'^{2}}}\Bigr]_{-\infty}^{\infty}=\frac{\lambda}{4\pi\eo\,r}\bigl[1-(-1)\bigr]=\frac{\lambda}{2\pi\eo\,r},
\]
pues $z'/\sqrt{r^{2}+z'^{2}}\to\pm1$ cuando $z'\to\pm\infty$. Por tanto
\[
  \boxed{\;\vE(\vc{r})=\frac{\lambda}{2\pi\eo\,r}\,\urho\;}
\]

\textbf{Potencial.} Como $\vE=E_r(r)\,\urho$ y $\vE=-\nab\phi$, basta integrar la componente radial. Observemos que \emph{no podemos} tomar $\phi(\infty)=0$: una carga que se extiende hasta el infinito hace divergir la integral
($\int^{\infty}E_r\dif r\propto\ln r\to\infty$). Tomamos entonces una referencia a distancia $r_0$, $\phi(r_0)=0$:
\[
  \phi(r)=-\int_{r_0}^{r}\frac{\lambda}{2\pi\eo\,r'}\,\dif r'=-\frac{\lambda}{2\pi\eo}\ln\frac{r}{r_0}
  \;\Longrightarrow\;\boxed{\;\phi(r)=\frac{\lambda}{2\pi\eo}\,\ln\frac{r_0}{r}\;}
\]
El potencial decrece logarítmicamente con la distancia al hilo (positivo para $r<r_0$ si $\lambda>0$).
\end{resolucion}

\begin{observacion}[Hilo finito y el límite de hilo infinito]\label{obs:hilo-finito}
Para un hilo \emph{finito} de longitud $L$, sobre $-L/2\le z'\le L/2$, la misma integral, con $\vc{R}=r\,\urho+(z-z')\,\uz$, da (se sustituye $u=z-z'$ y se usa la primitiva anterior y $\int\frac{u\,\dif u}{(r^{2}+u^{2})^{3/2}}=-\frac{1}{\sqrt{r^{2}+u^{2}}}$)
\begin{align*}
  \vE_{\mathrm{ex}}&=\frac{\lambda}{4\pi\eo\,r}\left[\frac{z+L/2}{\sqrt{r^{2}+(z+L/2)^{2}}}-\frac{z-L/2}{\sqrt{r^{2}+(z-L/2)^{2}}}\right]\urho\\
  &\quad-\frac{\lambda}{4\pi\eo}\left[\frac{1}{\sqrt{r^{2}+(z+L/2)^{2}}}-\frac{1}{\sqrt{r^{2}+(z-L/2)^{2}}}\right]\uz .
\end{align*}
Si $L\to\infty$, el corchete radial tiende a $2$ y el axial a $0$, y se recupera la aproximación de hilo infinito, $\vE_{\mathrm{ap}}=\frac{\lambda}{2\pi\eo r}\,\urho$ (en el esquema de la hoja del curso aparece $4\pi$ en el denominador;
el factor correcto es $2\pi$, porque el corchete vale $2$). La \figpanel{fig:sem-hilo}{b) y~(c} muestra qué tan buena es esa aproximación para un hilo finito: cerca del hilo y lejos de sus extremos el error es pequeño, y crece a medida que nos
alejamos radialmente de él, y también al acercarnos a la altura de los extremos.
\end{observacion}

% Figura (Problema 1.3): hilo infinito (geometría) e hilo finito frente a la aproximación de hilo infinito.
\begin{figure}[ht]
\centering
\pgfmathdeclarefunction{Erfin}{2}{%
  \begingroup
  \pgfmathparse{(1/(#1))*(((#2)+0.5)/sqrt((#1)*(#1)+((#2)+0.5)*((#2)+0.5))-((#2)-0.5)/sqrt((#1)*(#1)+((#2)-0.5)*((#2)-0.5)))}%
  \pgfmathsmuggle\pgfmathresult\endgroup}
\pgfmathdeclarefunction{Ezfin}{2}{%
  \begingroup
  \pgfmathparse{-(1/sqrt((#1)*(#1)+((#2)+0.5)*((#2)+0.5))-1/sqrt((#1)*(#1)+((#2)-0.5)*((#2)-0.5)))}%
  \pgfmathsmuggle\pgfmathresult\endgroup}
\pgfmathdeclarefunction{errfin}{2}{%
  \begingroup
  \pgfmathparse{100*abs(2/(#1)-sqrt(Erfin(#1,#2)*Erfin(#1,#2)+Ezfin(#1,#2)*Ezfin(#1,#2)))/sqrt(Erfin(#1,#2)*Erfin(#1,#2)+Ezfin(#1,#2)*Ezfin(#1,#2))}%
  \pgfmathsmuggle\pgfmathresult\endgroup}
\begin{minipage}[t]{0.36\textwidth}\centering
\begin{tikzpicture}[baseline=(current bounding box.north), x=1cm, y=1cm]
  \draw[eje] (-0.2,0) -- (3.4,0) node[below] {\small $x$};
  \draw[line width=2.2pt, azulprof] (0,-2.0) -- (0,2.3);
  \draw[eje] (0,2.3) -- (0,2.8) node[above] {\small $z$};
  \fill[azulprof] (0,0) circle (1.4pt);
  \coordinate (P) at (2.3,0);
  \coordinate (A) at (0,1.2); \coordinate (B) at (0,-1.2);
  \fill[azulacero] (A) circle (2pt) node[left, font=\small] {$\dif q'$};
  \fill[azulacero] (B) circle (2pt) node[left, font=\small] {$\dif q'$};
  \draw[->, azulclaro!70!azulacero, line width=0.7pt] (A) -- (P) node[pos=0.45, above right, font=\small, text=azulprof] {$\vc{R}$};
  \draw[->, azulclaro!70!azulacero, line width=0.7pt] (B) -- (P);
  \draw[vectrojo] (P) -- ++(0.55,0.28) node[above right, font=\scriptsize, text=acento] {$\dif\vE_{1}$};
  \draw[vectrojo] (P) -- ++(0.55,-0.28) node[below right, font=\scriptsize, text=acento] {$\dif\vE_{2}$};
  \draw[vect, azulprof, line width=1.2pt] (P) -- ++(1.1,0) node[right, font=\small] {$\vE$};
  \draw[grispizarra, densely dotted] (0,0) -- (P) node[midway, below, font=\small, text=azulprof] {$r$};
  \node[font=\scriptsize, text=grispizarra, anchor=west] at (-1.5,1.2) {$z'$};
  \node[font=\scriptsize, text=grispizarra, anchor=west] at (-1.5,-1.2) {$-z'$};
\end{tikzpicture}\\[2pt]{\small\scshape (a)}
\end{minipage}\hfill
\begin{minipage}[t]{0.62\textwidth}\centering
\begin{tikzpicture}[baseline=(current bounding box.north), x=7.6cm, y=2.4cm]
  \draw[grispizarra, line width=0.5pt] (-0.55,-1.05) rectangle (0.55,1.05);
  \foreach \xx in {-0.45,-0.35,...,0.46}{\foreach \zz in {-0.9,-0.8,...,0.91}{%
    \pgfmathsetmacro{\ra}{abs(\xx)}%
    \pgfmathsetmacro{\er}{Erfin(\ra,\zz)}\pgfmathsetmacro{\ez}{Ezfin(\ra,\zz)}%
    \pgfmathsetmacro{\em}{sqrt(\er*\er+\ez*\ez)}%
    \pgfmathsetmacro{\ft}{min(1,(\em^0.4)/3.2)}%
    \mapcolor{\ft}%
    \pgfmathsetmacro{\ux}{sign(\xx)*\er/\em}\pgfmathsetmacro{\uz}{\ez/\em}%
    \pgfmathsetmacro{\ll}{0.03*(0.55+\ft)}%
    \draw[-{Stealth[length=3pt,width=3pt]}, line width=0.5pt, mapc] (\xx-\ux*\ll*0.9,\zz-\uz*\ll*0.9/3.42) -- ++({\ux*\ll*1.8},{\uz*\ll*1.8/3.42});}}
  \draw[line width=2.2pt, azulprof] (0,-0.5) -- (0,0.5);
  \foreach \zz in {0,0.1,0.2,0.3}{\draw[azulprof, densely dashed, line width=0.5pt] (-0.5,\zz) -- (0.5,\zz);
     \node[font=\tiny, text=azulprof, fill=white, inner sep=0.6pt, anchor=west] at (0.505,\zz) {\pgfmathprintnumber[fixed,precision=1,use comma]{\zz}};}
  \foreach \t in {-0.4,-0.2,0,0.2,0.4}{\node[font=\tiny, anchor=north] at (\t,-1.05) {\pgfmathprintnumber[fixed,precision=1,use comma]{\t}};}
  \node[font=\scriptsize, anchor=north] at (0,-1.25) {$x/L$};
  \foreach \t in {-1,-0.5,0,0.5,1}{\node[font=\tiny, anchor=east] at (-0.555,\t) {\pgfmathprintnumber[fixed,precision=1,use comma]{\t}};}
  \node[font=\scriptsize, rotate=90, anchor=south] at (-0.66,0) {$z/L$};
\end{tikzpicture}\\[2pt]{\small\scshape (b)}
\end{minipage}\\[6pt]
\pgfplotsset{errax/.style={tema, width=0.265\textwidth, height=3.3cm, xmin=0, xmax=0.5, ymin=0, ymax=100,
   xtick={0,0.25,0.5}, ytick={0,50,100}, xlabel={$r/L$}, domain=0.02:0.5, samples=60, grid=major}}
\begin{tikzpicture}
  \begin{axis}[errax, name=a, ylabel={error (\%)}, title={\scriptsize $z/L=0{,}0$}]
    \addplot[azulprof, line width=1.0pt] {errfin(x,0.0)};
  \end{axis}
  \begin{axis}[errax, at={(a.east)}, anchor=west, xshift=0.35cm, name=b, yticklabels={}, title={\scriptsize $z/L=0{,}1$}]
    \addplot[azulprof, line width=1.0pt] {errfin(x,0.1)};
  \end{axis}
  \begin{axis}[errax, at={(b.east)}, anchor=west, xshift=0.35cm, name=c, yticklabels={}, title={\scriptsize $z/L=0{,}2$}]
    \addplot[azulprof, line width=1.0pt] {errfin(x,0.2)};
  \end{axis}
  \begin{axis}[errax, at={(c.east)}, anchor=west, xshift=0.35cm, yticklabels={}, title={\scriptsize $z/L=0{,}3$}]
    \addplot[azulprof, line width=1.0pt] {errfin(x,0.3)};
  \end{axis}
\end{tikzpicture}\\[2pt]{\small\scshape (c)}
\caption{Problema 1.3. (a)~Hilo infinito: los elementos $\dif q'$ en $z'$ y $-z'$ crean campos cuyas componentes $z$ se cancelan y
solo sobrevive la radial. (b)~Campo de un hilo \emph{finito} de longitud $L$ (trazo grueso) en el plano $(x,z)$, en unidades
de $\lambda/4\pi\eo L$ (el color y la longitud de las flechas codifican $\norm{\vE}$, de forma comprimida). (c)~Error relativo
entre el módulo del campo exacto del hilo finito y el de la aproximación de hilo infinito, $\lambda/2\pi\eo r$, a lo largo de las rectas horizontales
$z/L=0;\,0{,}1;\,0{,}2;\,0{,}3$ marcadas a trazos en (b): crece al alejarnos, pero la aproximación es muy buena cerca del hilo y lejos de sus extremos.}
\label{fig:sem-hilo}
\end{figure}


% =========================================================================
\subsection{Problema 1.4. Hilo infinito: campo por la ley de Gauss}\label{sem:hilo-gauss}
% =========================================================================
\begin{ejercicio}[Hilo infinito, ley de Gauss]\label{ejer:hilo-gauss}
Calcular mediante el teorema de Gauss el campo creado por un hilo infinito con una densidad lineal de carga uniforme $\lambda$.
\end{ejercicio}

\begin{resolucion}
\textbf{Idea.} Simetría cilíndrica: $\vE=E(r)\,\urho$ (la dirección radial cilíndrica, que en la hoja de problemas se escribe $\uh{r}$). La superficie gaussiana es un cilindro coaxial de radio $r$ y longitud $\ell$
(\cref{fig:sem-gauss-cil}), formado por la superficie lateral $S_L$ y dos tapas.

\textbf{Flujo.} En las tapas $\vE\perp\dS$ (el campo es radial y la normal de la tapa es axial), luego su flujo es cero. En la superficie lateral $\vE\parallel\dS$ y $E$ es constante:
\[
  \oiint_{S}\vE\cdot\dS=\iint_{S_L}\vE\cdot\dS+\iint_{\text{tapas}}\vE\cdot\dS=E\iint_{S_L}\dsc+0=E\cdot2\pi r\ell .
\]
La carga encerrada es $\Qenc=\lambda\,\ell$. Por la ley de Gauss,
\[
  2\pi r\ell\,E=\frac{\lambda\,\ell}{\eo}\quad\Longrightarrow\quad E(r)=\frac{\lambda}{2\pi\eo\,r}.
\]
Como sabemos que el campo tiene dirección radial, podemos escribirlo en forma vectorial:
\[
  \boxed{\;\vE(\vc{r})=\frac{\lambda}{2\pi\eo\,r}\,\urho\;}
\]
\end{resolucion}

\begin{observacion}[Comparación]
Coincide con el resultado por integración directa (\cref{sem:hilo}), donde costó varios pasos. La longitud $\ell$ de la superficie gaussiana se cancela, como debe: ninguna propiedad del campo puede depender de una elección arbitraria.
\end{observacion}

% Figura (Problema 1.4): superficie gaussiana cilíndrica alrededor de un hilo infinito.
\begin{figure}[ht]
\centering
\begin{tikzpicture}[x=1cm, y=1cm]
  \def\rx{1.4}\def\ry{0.42}\def\hc{3.2}
  \draw[line width=2.2pt, azulprof] (0,-0.9) -- (0,\hc+0.9);
  \node[font=\small, text=azulprof, anchor=west] at (0.08,\hc+0.7) {$\lambda$};
  % cilindro de Gauss
  \fill[verdeteo, opacity=0.06] (-\rx,0) -- (-\rx,\hc) arc (180:360:{\rx} and \ry) -- (\rx,0) arc (0:-180:{\rx} and \ry) -- cycle;
  \draw[verdeteo, line width=0.8pt] (-\rx,0) -- (-\rx,\hc) (\rx,0) -- (\rx,\hc);
  \draw[verdeteo, line width=0.8pt] (0,\hc) ellipse ({\rx} and \ry);
  \draw[verdeteo, line width=0.8pt] (-\rx,0) arc (180:360:{\rx} and \ry);
  \draw[verdeteo, line width=0.5pt, densely dashed] (-\rx,0) arc (180:0:{\rx} and \ry);
  \node[font=\small, text=verdeteo, anchor=east] at (-\rx-0.05,\hc-0.2) {$S_{\mathrm{G}}$};
  % campo radial sobre la superficie lateral
  \foreach \z in {0.6,1.6,2.6}{
    \draw[-{Stealth[length=4pt]}, cm1, line width=0.7pt] (\rx,\z) -- (\rx+0.8,\z);
    \draw[-{Stealth[length=4pt]}, cm1, line width=0.7pt] (-\rx,\z) -- (-\rx-0.8,\z);
    \draw[-{Stealth[length=4pt]}, cm1, line width=0.7pt] (0.5,\z-\ry-0.0) ++(0,-0.0) -- ++(0.28,-0.38);
    \draw[-{Stealth[length=4pt]}, cm1, line width=0.7pt] (-0.5,\z-\ry-0.0) -- ++(-0.28,-0.38);}
  \draw[vectverde, line width=0.9pt] (\rx,1.1) -- ++(0.5,0) node[right, font=\scriptsize, text=verdeteo] {$\dS$};
  \draw[grispizarra, <->, line width=0.4pt] (0,2.15) -- (\rx,2.15) node[midway, above, font=\scriptsize, text=azulprof] {$r$};
  \draw[grispizarra, <->, line width=0.4pt] (\rx+1.25,0) -- (\rx+1.25,\hc) node[midway, right, font=\small] {$\ell$};
\end{tikzpicture}
\caption{Problema 1.4. Superficie gaussiana cilíndrica, coaxial con el hilo, de radio $r$ y altura $\ell$. El campo es radial, de módulo constante sobre la
superficie lateral y paralelo a $\dS$ ($\vE\parallel\dS$); en las tapas $\vE\perp\dS$ y no hay flujo.}
\label{fig:sem-gauss-cil}
\end{figure}


% =========================================================================
\subsection{Problema 1.5. Plano infinito: campo por integración directa}\label{sem:plano}
% =========================================================================
\begin{ejercicio}[Plano infinito, integración directa]\label{ejer:plano}
Calcular por integración directa el campo $\vE$ creado por una superficie plana infinitamente grande con una densidad superficial de carga uniforme $\sigma$.
\end{ejercicio}

\begin{resolucion}
\textbf{Idea.} Situamos el plano en $z'=0$. Como la distribución es invariante frente a traslaciones en $x$ e $y$, el campo solo puede depender de $z$, y podemos tomar el punto campo sobre el eje $z$: $\vc{r}=z\,\uz$
(\figpanel{fig:sem-plano}{a}). Por la simetría especular respecto del plano, el campo en $-z$ es el reflejado del campo en $+z$.

\textbf{Geometría.} El elemento de carga está en $\vc{r}\,'=x'\,\ux+y'\,\uy$, con $\dif q'=\sigma\,\dSp$. Para aprovechar la simetría pasamos a coordenadas polares en el plano (usamos $s'$ y no $\rho$, para no
confundirlo con la densidad volúmica de carga): $x'=s'\cos\varphi'$, $y'=s'\sen\varphi'$, $\dSp=s'\,\dif s'\,\dif\varphi'$, $s'\in[0,\infty)$, $\varphi'\in[0,2\pi)$. Entonces
\[
  \vc{R}=\vc{r}-\vc{r}\,'=-s'\cos\varphi'\,\ux-s'\sen\varphi'\,\uy+z\,\uz,\qquad R=\sqrt{s'^{2}+z^{2}} .
\]
\textbf{Integral.} Con la \cref{eq:campo-sup-lin},
\[
  \vE(z)=\frac{\sigma}{4\pi\eo}\int_{0}^{2\pi}\!\!\int_{0}^{\infty}\frac{-s'\cos\varphi'\,\ux-s'\sen\varphi'\,\uy+z\,\uz}{(s'^{2}+z^{2})^{3/2}}\;s'\,\dif s'\,\dif\varphi' .
\]
Las componentes $x$ e $y$ contienen $\cos\varphi'$ y $\sen\varphi'$ y su integral en $\varphi'$ sobre un período completo es nula: $E_x=E_y=0$ (traslación simétrica en $x$ e $y$). Solo queda
\[
  \vE(z)=\frac{\sigma z}{4\pi\eo}\,2\pi\int_{0}^{\infty}\frac{s'\,\dif s'}{(s'^{2}+z^{2})^{3/2}}\,\uz=\frac{\sigma z}{2\eo}\Bigl[-\frac{1}{\sqrt{s'^{2}+z^{2}}}\Bigr]_{0}^{\infty}\uz=\frac{\sigma z}{2\eo}\,\frac{1}{\sqrt{z^{2}}}\,\uz,
\]
y como $\sqrt{z^{2}}=\lvert z\rvert$,
\[
  \boxed{\;\vE(\vc{r})=\frac{\sigma}{2\eo}\,\frac{z}{\lvert z\rvert}\,\uz=\frac{\sigma}{2\eo}\,\sgn(z)\,\uz\;}
\]
\end{resolucion}

\begin{observacion}[Un campo uniforme a cada lado, y un salto]
El campo no depende de la distancia al plano: a cada lado es uniforme, de módulo $\sigma/2\eo$, saliente si $\sigma>0$. (Al alejarnos, cada elemento del plano crea menos campo, pero son más los elementos que contribuyen con una componente axial apreciable; la compensación es exacta.) La densidad superficial de carga implica una discontinuidad de $\sigma/2\eo-(-\sigma/2\eo)=\sigma/\eo$ en $z=0$, exactamente la que predicen las condiciones de contorno (\cref{cor:contorno-E}).
\end{observacion}

% =========================================================================
\subsection{Problema 1.6. Plano infinito: campo por la ley de Gauss}\label{sem:plano-gauss}
% =========================================================================
\begin{ejercicio}[Plano infinito, ley de Gauss]\label{ejer:plano-gauss}
Calcular mediante el teorema de Gauss el campo creado por una superficie plana infinitamente grande con una densidad superficial de carga uniforme $\sigma$.
\end{ejercicio}

\begin{resolucion}
\textbf{Idea.} Simetría plana (especular): $\vE=E_z(z)\,\uz$, y $E_z(-z)=-E_z(z)$. La superficie gaussiana es un cilindro recto de radio $R_c$ con eje $z$, cuyas dos tapas, de área $S_{\mathrm{tapa}}=\pi R_c^{2}$, están a ambos lados del plano y a la misma
distancia $\lvert z\rvert$ (\figpanel{fig:sem-plano}{b}).

\textbf{Flujo.} Sobre la superficie lateral, $\vE\perp\dS$ y no hay flujo; en las tapas, $\vE\parallel\dS$ y $E$ es el mismo:
\[
  \flujo=\iint_{S_L}\vE\cdot\dS+\iint_{\text{tapas}}\vE\cdot\dS=0+E\,(2S_{\mathrm{tapa}}).
\]
La carga encerrada es la del disco del plano interior al cilindro: $\Qenc=\sigma\,\pi R_c^{2}=\sigma S_{\mathrm{tapa}}$. Entonces $2E\,S_{\mathrm{tapa}}=\sigma S_{\mathrm{tapa}}/\eo$, y
\[
  \boxed{\;\vE(\vc{r})=\frac{\sigma}{2\eo}\,\frac{z}{\lvert z\rvert}\,\uz\;}
\]
igual que en el problema anterior. El área de las tapas se simplifica, y $E$ no depende de $z$.
\end{resolucion}

% Figura (Problemas 1.5 y 1.6): plano infinito cargado, integración directa y superficie gaussiana.
\begin{figure}[ht]
\centering
\tdplotsetmaincoords{66}{125}
\begin{minipage}[t]{0.50\textwidth}\centering
\begin{tikzpicture}[tdplot_main_coords, scale=0.95, baseline=(current bounding box.north)]
  \fill[grisclaro!70, opacity=0.8] (-2.6,-2.6,0) -- (2.6,-2.6,0) -- (2.6,2.6,0) -- (-2.6,2.6,0) -- cycle;
  \draw[grispizarra!70, line width=0.4pt] (-2.6,-2.6,0) -- (2.6,-2.6,0) -- (2.6,2.6,0) -- (-2.6,2.6,0) -- cycle;
  \draw[eje] (0,0,0) -- (3.0,0,0) node[below left] {\small $x$};
  \draw[eje] (0,0,0) -- (0,3.0,0) node[right] {\small $y$};
  \draw[eje] (0,0,0) -- (0,0,3.6) node[above] {\small $z$};
  \coordinate (P) at (0,0,2.4); \coordinate (S) at (1.4,0.9,0);
  \draw[azulprof, densely dotted, line width=0.5pt] (0,0,0) -- (1.4,0,0) -- (S);
  \draw[verdeteo, line width=0.5pt, domain=0:32, samples=20] plot ({0.8*cos(\x)},{0.8*sin(\x)},0);
  \node[font=\scriptsize, text=verdeteo] at (1.0,0.25,0) {$\varphi'$};
  \draw[->, azulclaro!80!azulacero, line width=0.8pt] (0,0,0) -- (S) node[pos=0.55, below, font=\small, text=azulprof] {$\vc{r}\,'$};
  \draw[vect, azulprof, line width=1.1pt] (S) -- (P) node[pos=0.5, right, font=\small] {$\vc{R}$};
  \draw[vectrojo, line width=1.2pt] (P) -- ++(0,0,0.9) node[right, font=\small, text=acento] {$\vE$};
  \fill[acento] (S) circle (1.5pt) node[above right, font=\small, text=acento] {$\dSp$};
  \node[font=\small, text=azulprof] at (-1.8,1.6,0) {$\sigma$};
  \fill[azulprof] (P) circle (1.5pt) node[left, font=\small] {$P$};
\end{tikzpicture}\\[2pt]{\small\scshape (a)\; integración directa}
\end{minipage}\hfill
\begin{minipage}[t]{0.46\textwidth}\centering
\begin{tikzpicture}[tdplot_main_coords, scale=0.95, baseline=(current bounding box.north)]
  \def\Rc{0.9}\def\Hc{1.4}
  \fill[grisclaro!70, opacity=0.8] (-2.6,-2.6,0) -- (2.6,-2.6,0) -- (2.6,2.6,0) -- (-2.6,2.6,0) -- cycle;
  \draw[grispizarra!70, line width=0.4pt] (-2.6,-2.6,0) -- (2.6,-2.6,0) -- (2.6,2.6,0) -- (-2.6,2.6,0) -- cycle;
  \draw[eje] (0,0,-2.0) -- (0,0,3.2) node[above] {\small $z$};
  \foreach \zz in {\Hc,-\Hc}{
    \draw[verdeteo, line width=0.8pt, domain=0:360, samples=60, fill=verdeteo, fill opacity=0.08] plot ({\Rc*cos(\x)},{\Rc*sin(\x)},\zz);}
  \draw[verdeteo, line width=0.8pt] ({\Rc*cos(35)},{\Rc*sin(35)},-\Hc) -- ({\Rc*cos(35)},{\Rc*sin(35)},\Hc);
  \draw[verdeteo, line width=0.8pt] ({\Rc*cos(215)},{\Rc*sin(215)},-\Hc) -- ({\Rc*cos(215)},{\Rc*sin(215)},\Hc);
  \draw[vectrojo, line width=1.1pt] (0,0,\Hc) -- (0,0,\Hc+0.9) node[right, font=\small, text=acento] {$\vE$};
  \draw[vectrojo, line width=1.1pt] (0,0,-\Hc) -- (0,0,-\Hc-0.9) node[right, font=\small, text=acento] {$\vE$};
  \draw[vectverde, line width=0.9pt] (0.5,0.3,\Hc) -- ++(0,0,0.5) node[right, font=\scriptsize, text=verdeteo] {$\dS$};
  \node[font=\small, text=verdeteo] at (1.3,1.3,\Hc-0.4) {$S_{\mathrm{G}}$};
  \node[font=\small, text=azulprof] at (-1.8,1.6,0) {$\sigma$};
\end{tikzpicture}\\[2pt]{\small\scshape (b)\; superficie gaussiana}
\end{minipage}
\caption{Problemas 1.5 y 1.6. (a)~Plano $z'=0$ con densidad superficial uniforme $\sigma$; el punto campo está sobre el eje $z$, el elemento $\dSp$ en
coordenadas polares $(s',\varphi')$ y $\vc{R}=\vc{r}-\vc{r}\,'$. (b)~Superficie gaussiana: cilindro recto que atraviesa el plano, con las dos tapas a la misma distancia.
Por la simetría especular, $\vE$ es perpendicular al plano y de módulo igual a ambos lados, con sentidos opuestos.}
\label{fig:sem-plano}
\end{figure}


% =========================================================================
\subsection{Problema 1.7. Anillo, disco y plano: campo y potencial sobre el eje}\label{sem:anillo}
% =========================================================================
\begin{ejercicio}[Del anillo al plano]\label{ejer:anillo}
Calcular el campo y el potencial en el eje de simetría perpendicular a un anillo de radio $a$ con una carga total $Q$ uniformemente distribuida. Utilizar dicho resultado para calcular el campo y el potencial para un disco de radio $a$
y un plano infinito uniformemente cargados.
\end{ejercicio}

\begin{resolucion}
\textbf{Idea.} El anillo es una distribución lineal; el disco es una colección de anillos concéntricos y el plano es un disco con $a\to\infty$. Así, solo hay que \emph{integrar} una vez el resultado del anillo.

\textbf{Anillo.} Con el anillo en el plano $z'=0$ y el punto campo en $\vc{r}=z\,\uz$ (\figpanel{fig:sem-anillo}{a}), todos los puntos del anillo están a la misma distancia $R=\sqrt{a^{2}+z^{2}}$ del punto campo. Con $\lambda=Q/2\pi a$ y $\dlp=a\,\dif\varphi'$,
\[
  \phi(z)=\frac{1}{4\pi\eo}\int\frac{\lambda\,\dlp}{R}=\frac{1}{4\pi\eo}\,\frac{Q}{\sqrt{a^{2}+z^{2}}},\qquad\boxed{\;\phi_{\mathrm{anillo}}(z)=\frac{Q}{4\pi\eo\,(a^{2}+z^{2})^{1/2}}\;}
\]
(el potencial es mucho más sencillo que el campo, pues $R$ es constante y no hay que descomponer vectores). Para el campo, el elemento $\dif q'$ crea un $\dif\vE$ con una componente axial $\dif E_z=\frac{\dif q'}{4\pi\eo R^{2}}\cos\alpha=\frac{\dif q'}{4\pi\eo R^{2}}\frac{z}{R}$
y otra perpendicular al eje que se anula al sumar sobre elementos diametralmente opuestos. Por tanto
\[
  \boxed{\;\vE_{\mathrm{anillo}}(z)=\frac{Q\,z}{4\pi\eo\,(a^{2}+z^{2})^{3/2}}\,\uz\;}
\]
Coincide con $-\partial_z\phi$ \emph{sobre el eje}: $-\frac{\dif}{\dif z}\frac{Q}{4\pi\eo}(a^{2}+z^{2})^{-1/2}=\frac{Q}{4\pi\eo}\frac{z}{(a^{2}+z^{2})^{3/2}}$ (derivar $\phi$ respecto de $z$ solo da la componente axial, y es lo único que hay en el eje).

\textbf{Disco.} Un disco de radio $a$ con densidad $\sigma$ se descompone en anillos de radio $r'\in[0,a]$ y anchura $\dif r'$, con carga $\dif q'=\sigma\,2\pi r'\,\dif r'$ (\figpanel{fig:sem-anillo}{b}). Sumando los potenciales y los campos de los anillos:
\[
  \phi_{\mathrm{disco}}(z)=\int_{0}^{a}\frac{\sigma\,2\pi r'\,\dif r'}{4\pi\eo\sqrt{r'^{2}+z^{2}}}=\frac{\sigma}{2\eo}\Bigl[\sqrt{r'^{2}+z^{2}}\Bigr]_{0}^{a}=\frac{\sigma}{2\eo}\Bigl(\sqrt{a^{2}+z^{2}}-\lvert z\rvert\Bigr),
\]
\[
  \vE_{\mathrm{disco}}(z)=\int_{0}^{a}\frac{\sigma\,2\pi r'\,z\,\dif r'}{4\pi\eo\,(r'^{2}+z^{2})^{3/2}}\,\uz=\frac{\sigma z}{2\eo}\Bigl[-\frac{1}{\sqrt{r'^{2}+z^{2}}}\Bigr]_{0}^{a}\uz
  =\boxed{\;\frac{\sigma}{2\eo}\Bigl(\frac{z}{\lvert z\rvert}-\frac{z}{\sqrt{a^{2}+z^{2}}}\Bigr)\uz\;}
\]
(en el potencial, $\sqrt{z^{2}}=\lvert z\rvert$ en el extremo inferior).

\textbf{Plano infinito.} Es el límite $a\to\infty$. El campo del disco tiende a
\[
  \boxed{\;\vE_{\mathrm{plano}}(z)=\frac{\sigma}{2\eo}\,\frac{z}{\lvert z\rvert}\,\uz\;}
\]
(el término $z/\sqrt{a^{2}+z^{2}}\to0$), que es el resultado de los problemas 1.5 y 1.6. En cambio, el potencial \emph{diverge}: $\phi_{\mathrm{disco}}=\frac{\sigma}{2\eo}(a-\lvert z\rvert+O(1/a))$, con una constante $\sigma a/2\eo\to\infty$ independiente de $z$. Como el potencial
está definido salvo una constante, la restamos (equivale a tomar $\phi(0)=0$ como referencia) y queda
\[
  \boxed{\;\phi_{\mathrm{plano}}(z)=-\frac{\sigma}{2\eo}\,\lvert z\rvert\;}\qquad\Bigl(\text{y se comprueba que }-\frac{\dif\phi}{\dif z}=\frac{\sigma}{2\eo}\,\frac{z}{\lvert z\rvert}\Bigr).
\]
\end{resolucion}

\begin{observacion}[Interpretación y casos límite]
\begin{enumerate}[label=(\alph*)]
  \item \emph{Lejos del anillo o del disco}, $\lvert z\rvert\gg a$: $\vE_{\mathrm{anillo}}\to\frac{Q}{4\pi\eo z^{2}}\,\uz$, y para el disco (con $Q=\sigma\pi a^{2}$) el desarrollo $\frac{z}{\lvert z\rvert}-\frac{z}{\sqrt{a^{2}+z^{2}}}\simeq\frac{a^{2}}{2z^{2}}$ da
        también $\frac{Q}{4\pi\eo z^{2}}$: visto desde lejos, todo objeto cargado es una carga puntual.
  \item \emph{El máximo del anillo.} $E_{\mathrm{anillo}}$ se anula en $z=0$ (en el centro las contribuciones opuestas se cancelan) y en el infinito; tiene un máximo en $\frac{\dif}{\dif z}\bigl[z(a^{2}+z^{2})^{-3/2}\bigr]=0\Leftrightarrow z=a/\sqrt2$,
        de valor $\frac{Q}{4\pi\eo a^{2}}\frac{2}{3\sqrt3}$ (\cref{fig:sem-anillo}).
  \item \emph{El salto del disco.} $\vE_{\mathrm{disco}}$ es discontinuo en $z=0$ para cualquier radio $a$, con salto $\frac{\sigma}{2\eo}-\bigl(-\frac{\sigma}{2\eo}\bigr)=\frac{\sigma}{\eo}$ (la ampliación de la \figpanel{fig:sem-anillo}{c}): es la condición de contorno
        de \cref{eq:contorno-E}, que solo depende de la densidad local de carga y no del tamaño del disco. En cambio, el potencial es continuo.
\end{enumerate}
\end{observacion}

% Figura (Problema 1.7): anillo y disco; campos y potenciales sobre el eje.
\begin{figure}[ht]
\centering
\tdplotsetmaincoords{66}{125}
\begin{minipage}[t]{0.48\textwidth}\centering
\begin{tikzpicture}[tdplot_main_coords, scale=0.9, baseline=(current bounding box.north)]
  \draw[eje] (0,0,0) -- (3.0,0,0) node[below left] {\small $x$};
  \draw[eje] (0,0,0) -- (0,3.0,0) node[right] {\small $y$};
  \draw[eje] (0,0,0) -- (0,0,3.4) node[above] {\small $z$};
  \draw[azulprof, line width=2.0pt, domain=0:360, samples=90] plot ({1.9*cos(\x)},{1.9*sin(\x)},0);
  \coordinate (P) at (0,0,2.3); \coordinate (S) at ({1.9*cos(35)},{1.9*sin(35)},0);
  \draw[vect, azulclaro!70!azulacero, line width=0.8pt] (S) -- (P) node[pos=0.45, right, font=\small, text=azulprof] {$\vc{R}$};
  \draw[vectrojo, line width=1.0pt] (P) -- ++({-0.55*cos(35)*0.62},{-0.55*sin(35)*0.62},-0.55*0.5)  node[below left, font=\scriptsize, text=acento] {$\dif\vE$};
  \draw[vectrojo, line width=1.4pt] (P) -- ++(0,0,0.9) node[right, font=\small, text=acento] {$\vE_{\mathrm{anillo}}$};
  \fill[acento] (S) circle (1.6pt) node[above right, font=\small, text=acento] {$\dif q'$};
  \draw[grispizarra, densely dotted] (0,0,0) -- (S);
  \node[font=\scriptsize, text=azulprof, anchor=south west] at (0.3,-0.2,0) {$a$};
  \node[font=\small, text=azulprof] at (-2.2,0.4,0) {$Q$};
  \fill[azulprof] (P) circle (1.4pt) node[left, font=\small] {$z$};
\end{tikzpicture}\\[2pt]{\small\scshape (a)\; anillo}
\end{minipage}\hfill
\begin{minipage}[t]{0.48\textwidth}\centering
\begin{tikzpicture}[tdplot_main_coords, scale=0.9, baseline=(current bounding box.north)]
  \draw[eje] (0,0,0) -- (3.0,0,0) node[below left] {\small $x$};
  \draw[eje] (0,0,0) -- (0,3.0,0) node[right] {\small $y$};
  \draw[eje] (0,0,0) -- (0,0,3.4) node[above] {\small $z$};
  \draw[azulprof, line width=0.8pt, fill=grisclaro!70, opacity=0.9, domain=0:360, samples=90] plot ({2.1*cos(\x)},{2.1*sin(\x)},0) -- cycle;
  \draw[azulprof, line width=2.4pt, domain=0:360, samples=90] plot ({1.3*cos(\x)},{1.3*sin(\x)},0);
  \draw[azulclaro!80!azulacero, line width=2.4pt, domain=0:360, samples=90, opacity=0.0] plot ({1.35*cos(\x)},{1.35*sin(\x)},0);
  \coordinate (P) at (0,0,2.3);
  \draw[vectrojo, line width=1.4pt] (P) -- ++(0,0,0.9) node[right, font=\small, text=acento] {$\vE_{\mathrm{disco}}$};
  \draw[grispizarra, densely dotted] (0,0,0) -- ({1.3*cos(200)},{1.3*sin(200)},0) node[midway, below, font=\scriptsize, text=azulprof] {$r'$};
  \node[font=\scriptsize, text=azulprof, anchor=north west] at ({1.5*cos(-40)},{1.5*sin(-40)},0) {$\dif r'$};
  \node[font=\small, text=azulprof] at (-2.3,1.2,0) {$\sigma$};
  \draw[grispizarra, densely dotted] (0,0,0) -- ({2.1*cos(300)},{2.1*sin(300)},0) node[midway, below left, font=\scriptsize, text=azulprof] {$a$};
  \fill[azulprof] (P) circle (1.4pt) node[left, font=\small] {$z$};
\end{tikzpicture}\\[2pt]{\small\scshape (b)\; disco}
\end{minipage}\\[8pt]
\pgfplotsset{anax/.style={tema, width=0.40\textwidth, height=3.6cm, xlabel={$z/a$}, domain=-6:6, samples=140}}
\begin{tikzpicture}
  \begin{axis}[anax, name=A, xmin=-6, xmax=6, ymin=-1.1, ymax=1.1, ytick={-1,-0.5,0,0.5,1}, title={\scriptsize anillo},
      legend style={at={(0.98,0.04)}, anchor=south east, font=\tiny}]
    \addplot[azulprof, line width=1.0pt] {x/(1+x^2)^1.5/0.3849};
    \addplot[acento, line width=1.0pt, densely dashed] {1/sqrt(1+x^2)};
    \legend{$E_z/E_{\max}$, $\phi/\phi(0)$}
  \end{axis}
  \begin{axis}[anax, at={(A.east)}, anchor=west, xshift=1.3cm, name=B, xmin=-6, xmax=6, ymin=-1.3, ymax=1.3, ytick={-1,0,1}, title={\scriptsize disco y plano}, legend style={at={(0.98,0.04)}, anchor=south east, font=\tiny}]
    \addplot[azulacero, line width=1.0pt, domain=-6:-0.003] {sign(x)-x/sqrt(1+x^2)};
    \addplot[azulacero, line width=1.0pt, domain=0.003:6] {sign(x)-x/sqrt(1+x^2)};
    \addplot[acento, line width=0.9pt, densely dashed, domain=-6:-0.003] {-1};
    \addplot[acento, line width=0.9pt, densely dashed, domain=0.003:6] {1};
    \legend{$E_{\mathrm{disco}}$,,$E_{\mathrm{plano}}$}
  \end{axis}
  \begin{axis}[anax, at={(A.south)}, anchor=north, yshift=-2.1cm, name=C, xmin=-0.1, xmax=0.1, ymin=-1.3, ymax=1.3, ytick={-1,0,1}, xtick={-0.1,0,0.1}, title={\scriptsize ampliación en $z\approx0$}, domain=-0.1:0.1]
    \addplot[azulacero, line width=1.0pt, domain=-0.1:-0.0005, samples=60] {sign(x)-x/sqrt(1+x^2)};
    \addplot[azulacero, line width=1.0pt, domain=0.0005:0.1, samples=60] {sign(x)-x/sqrt(1+x^2)};
    \addplot[acento, line width=0.9pt, densely dashed, domain=-0.1:-0.0005, samples=2] {-1};
    \addplot[acento, line width=0.9pt, densely dashed, domain=0.0005:0.1, samples=2] {1};
  \end{axis}
  \begin{axis}[anax, at={(B.south)}, anchor=north, yshift=-2.1cm, xmin=-6, xmax=6, ymin=-6.5, ymax=0.2, ytick={-6,-4,-2,0}, title={\scriptsize potencial}, legend style={at={(0.5,0.04)}, anchor=south, font=\tiny}]
    \addplot[azulacero, line width=1.0pt] {sqrt(1+x^2)-1-abs(x)};
    \addplot[acento, line width=0.9pt, densely dashed] {-abs(x)};
    \legend{$\phi_{\mathrm{disco}}$,$\phi_{\mathrm{plano}}$}
  \end{axis}
\end{tikzpicture}\\[2pt]{\small\scshape (c)}
\caption{Problema 1.7. (a)~Anillo de radio $a$ y carga $Q$: las componentes perpendiculares al eje se anulan por simetría y solo queda $E_z$. (b)~El disco
se descompone en anillos de radio $r'$ y anchura $\dif r'$. (c)~Sobre el eje: campo (normalizado a su máximo $Q/(4\pi\eo a^{2})\cdot2/(3\sqrt3)$ en $z=a/\sqrt2$) y
potencial del anillo; campo del disco, que tiende al del plano ($\sigma/2\eo$) y presenta un salto $\sigma/\eo$ en $z=0$ (ampliación); potenciales
(en unidades de $\sigma a/2\eo$ y con $\phi(0)=0$ para el plano).}
\label{fig:sem-anillo}
\end{figure}


% =========================================================================
\subsection{Problema 1.8. Lámina con densidad volúmica de carga}\label{sem:lamina}
% =========================================================================
\begin{ejercicio}[Lámina cargada en volumen]\label{ejer:lamina}
Calcular el campo asociado a una distribución volúmica de carga uniformemente distribuida (densidad $\rho$) entre los dos planos infinitos $z=+a/2$ y $z=-a/2$.
\end{ejercicio}

\begin{resolucion}
\textbf{Idea.} Es una distribución con simetría plana, aunque no infinitamente delgada: $\vE=E_z(z)\,\uz$ con $E_z(-z)=-E_z(z)$. La integración directa sería muy trabajosa, pero la ley de Gauss es inmediata. Tomamos un cilindro recto, de bases de área $A$ paralelas a la lámina y
simétricas respecto de ella, de modo que sus tapas están en $\pm z$ (\cref{fig:sem-lamina}). Sobre la superficie lateral $\vE\perp\dS$, y el flujo es $2E\,A$ en ambos casos.

\textbf{Fuera, $\lvert z\rvert>a/2$.} El cilindro atraviesa toda la lámina y encierra $\Qenc=\rho\,A\,a$:
\[
  2E\,A=\frac{1}{\eo}\rho A a\quad\Longrightarrow\quad\boxed{\;\vE(\vc{r})=\frac{\rho\,a}{2\eo}\,\frac{z}{\lvert z\rvert}\,\uz\qquad\bigl(\lvert z\rvert>a/2\bigr)\;}
\]
con $\rho a$ la carga por unidad de superficie: igual que un plano con $\sigma=\rho a$.

\textbf{Dentro, $\lvert z\rvert<a/2$.} Ahora el cilindro solo encierra un espesor $2\lvert z\rvert$ de lámina: $\Qenc=\rho\,A\cdot2z$ (tomando $z>0$; por simetría se obtiene lo mismo para $z<0$ con el signo adecuado). Entonces $2E\,A=\frac{1}{\eo}\,2z\,\rho A$ y
\[
  \boxed{\;\vE(\vc{r})=\frac{\rho\,z}{\eo}\,\uz\qquad\bigl(\lvert z\rvert<a/2\bigr)\;}
\]
\end{resolucion}

\begin{observacion}[Interpretación]
El campo \emph{crece linealmente} con $z$ dentro de la lámina (dependencia lineal, $\vE\propto z$) hasta alcanzar $\rho a/2\eo$ en las caras ($z=\pm a/2$), y es constante fuera (\figpanel{fig:sem-lamina}{c}). Con una densidad volúmica de carga \emph{el campo es continuo}: en $z=\pm a/2$ las dos expresiones dan $\rho a/2\eo$,
porque no hay carga superficial, de acuerdo con las condiciones de contorno. Si adelgazamos la lámina con $\rho\to\infty$ manteniendo $\sigma=\rho a$ fijo, la rampa interior se hace infinitamente empinada y se recupera el plano de los problemas 1.5--1.6, con su salto $\sigma/\eo$.
\end{observacion}

% Figura (Problema 1.8): lámina con densidad volúmica; superficies gaussianas y campo.
\begin{figure}[ht]
\centering
\tdplotsetmaincoords{66}{125}
\newcommand{\laminaBox}[1]{%
  % caja de espesor 1 (z de -0.5 a 0.5): caras visibles +x, +y y superior
  \fill[azulclaro!55, opacity=#1] (2.2,-2.2,-0.5) -- (2.2,2.2,-0.5) -- (2.2,2.2,0.5) -- (2.2,-2.2,0.5) -- cycle;
  \fill[azulclaro!75, opacity=#1] (-2.2,2.2,-0.5) -- (2.2,2.2,-0.5) -- (2.2,2.2,0.5) -- (-2.2,2.2,0.5) -- cycle;
  \fill[azulclaro!35, opacity=#1] (-2.2,-2.2,0.5) -- (2.2,-2.2,0.5) -- (2.2,2.2,0.5) -- (-2.2,2.2,0.5) -- cycle;
  \draw[azulprof, line width=0.6pt] (2.2,-2.2,-0.5) -- (2.2,2.2,-0.5) -- (-2.2,2.2,-0.5) (2.2,2.2,-0.5) -- (2.2,2.2,0.5);
  \draw[azulprof, line width=0.6pt] (-2.2,-2.2,0.5) -- (2.2,-2.2,0.5) -- (2.2,2.2,0.5) -- (-2.2,2.2,0.5) -- cycle;
  \draw[azulprof, line width=0.6pt] (2.2,-2.2,-0.5) -- (2.2,-2.2,0.5) (-2.2,2.2,-0.5) -- (-2.2,2.2,0.5);}
\begin{minipage}[t]{0.31\textwidth}\centering
\begin{tikzpicture}[tdplot_main_coords, scale=0.78, baseline=(current bounding box.north)]
  \laminaBox{0.85}
  \tdCilindro[verdeteo!10]{0.8}{-1.7}{1.7}{verdeteo}{0.8pt}
  \draw[vectrojo, line width=1.2pt] (0,0,1.7) -- (0,0,2.5); \draw[vectrojo, line width=1.2pt] (0,0,-1.7) -- (0,0,-2.5);
  \node[font=\small, text=azulprof, anchor=west] at (2.25,0.0,0.0) {$\rho$};
  \node[font=\scriptsize, text=verdeteo, anchor=west] at ({0.95*cos(125)},{0.95*sin(125)},1.55) {$S_{\mathrm{G}}$};
  \draw[grispizarra, <->, line width=0.4pt] (2.45,-2.45,-0.5) -- (2.45,-2.45,0.5) node[midway, left, font=\scriptsize] {$a$};
\end{tikzpicture}\\[2pt]{\small\scshape (a)\; $\lvert z\rvert>a/2$}
\end{minipage}\hfill
\begin{minipage}[t]{0.31\textwidth}\centering
\begin{tikzpicture}[tdplot_main_coords, scale=0.78, baseline=(current bounding box.north)]
  \laminaBox{0.55}
  \tdCilindro[verdeteo!14]{0.8}{-0.25}{0.25}{verdeteo}{0.8pt}
  \draw[vectrojo, line width=1.2pt] (0,0,0.25) -- (0,0,1.3); \draw[vectrojo, line width=1.2pt] (0,0,-0.25) -- (0,0,-1.3);
  \node[font=\small, text=azulprof, anchor=west] at (2.25,0.0,0.0) {$\rho$};
  \node[font=\scriptsize, text=verdeteo, anchor=west] at ({1.0*cos(125)},{1.0*sin(125)},0.45) {$S_{\mathrm{G}}$};
  \draw[grispizarra, <->, line width=0.4pt] ({1.05*cos(305)},{1.05*sin(305)},-0.25) -- ({1.05*cos(305)},{1.05*sin(305)},0.25) node[midway, left, font=\scriptsize] {$2z$};
\end{tikzpicture}\\[2pt]{\small\scshape (b)\; $\lvert z\rvert<a/2$}
\end{minipage}\hfill
\begin{minipage}[t]{0.35\textwidth}\centering
\begin{tikzpicture}[baseline=(current bounding box.north)]
  \begin{axis}[tema, width=\textwidth, height=4.4cm, xlabel={$z/a$}, ylabel={$E_z/(\rho a/2\eo)$}, xmin=-1.5, xmax=1.5, ymin=-1.3, ymax=1.3,
     xtick={-1,-0.5,0,0.5,1}, xticklabels={$-1$,$-\tfrac12$,0,$\tfrac12$,1}, ytick={-1,0,1}, domain=-1.5:1.5, samples=100]
    \addplot[azulprof, line width=1.2pt, domain=-0.5:0.5] {2*x};
    \addplot[azulprof, line width=1.2pt, domain=0.5:1.5] {1};
    \addplot[azulprof, line width=1.2pt, domain=-1.5:-0.5] {-1};
    \addplot[grispizarra, densely dashed] coordinates {(-0.5,-1.3) (-0.5,1.3)};
    \addplot[grispizarra, densely dashed] coordinates {(0.5,-1.3) (0.5,1.3)};
  \end{axis}
\end{tikzpicture}\\[2pt]{\small\scshape (c)}
\end{minipage}
\caption{Problema 1.8. Lámina de espesor $a$ con densidad volúmica uniforme $\rho$ y superficies gaussianas (cilindros de área de base $A$ con tapas simétricas
respecto del plano medio). (a)~Fuera de la lámina el cilindro la atraviesa y encierra $\rho A a$. (b)~Dentro, encierra solo la parte de espesor $2z$. (c)~El campo es lineal en el interior
y constante, de módulo $\rho a/2\eo$, fuera; es continuo en las caras, porque no hay carga superficial.}
\label{fig:sem-lamina}
\end{figure}


% =========================================================================
\subsection{Problema 1.9. Superficie esférica cargada}\label{sem:cascara}
% =========================================================================
\begin{ejercicio}[Cáscara esférica]\label{ejer:cascara}
Calcular el campo y el potencial creados por una superficie esférica de radio $a$ con una carga total $Q$ uniformemente distribuida. Representar gráficamente cómo varían el campo y el potencial en función de la distancia al origen.
\end{ejercicio}

\begin{resolucion}
\textbf{Idea.} Es el mismo problema que el 1.2, pero con la carga concentrada en la superficie; la simetría esférica sigue permitiendo $\vE=E(r)\,\ur$ y una gaussiana esférica de radio $r$: $4\pi r^{2}E=\Qenc/\eo$.

\textbf{Campo.} Si $r<a$, la gaussiana no encierra carga ($\Qenc=0$): $\vE=\vc{0}$. Si $r>a$, encierra toda la carga: $E=Q/4\pi\eo r^{2}$:
\[
  \boxed{\;\vE(\vc{r})=\begin{cases}\vc{0},&r<a,\\[4pt]\dfrac{Q}{4\pi\eo r^{2}}\,\ur,&r>a.\end{cases}\;}
\]
\textbf{Potencial.} Con $\phi(\infty)=0$, fuera es $\phi=Q/4\pi\eo r$ y en $r=a$ vale $Q/4\pi\eo a$. Dentro el campo es nulo, luego $\phi(r)-\phi(a)=-\int_a^r0\,\dif r'=0$ y el potencial es \emph{constante}:
\[
  \boxed{\;\phi(r)=\begin{cases}\dfrac{Q}{4\pi\eo\,a},&r\le a,\\[6pt]\dfrac{Q}{4\pi\eo\,r},&r\ge a.\end{cases}\;}
\]
\end{resolucion}

\begin{observacion}[Interpretación y condiciones de contorno]
Dentro de la cáscara no hay campo: es la propiedad en que se basa el experimento de Cavendish (\cref{prop:cavendish}) y el análogo del teorema de Newton de las cáscaras. Aunque el interior esté a un potencial no nulo, no hay fuerza (el potencial es constante, su gradiente es nulo).
Al cruzar la superficie, el campo \emph{salta} de $0$ a $Q/4\pi\eo a^{2}=\sigma/\eo$, con $\sigma=Q/4\pi a^{2}$, exactamente lo que dice el \cref{cor:contorno-E}; el potencial, en cambio, es continuo (\cref{cor:contorno-phi}), pero su derivada no: se «dobla» en $r=a$ (\cref{fig:sem-cascara}).
Observemos la diferencia con la esfera maciza (\cref{sem:esfera}), donde $E$ y $\phi$ eran continuos.
\end{observacion}

% Figura (Problema 1.9): cáscara esférica; campo discontinuo y potencial continuo.
\begin{figure}[ht]
\centering
\begin{minipage}[t]{0.28\textwidth}\centering
\begin{tikzpicture}[baseline=(current bounding box.north), scale=0.9]
  \draw[azulprof, line width=1.6pt] (0,0) circle (1.1);
  \foreach \a in {20,70,...,340}{\node[font=\tiny, text=azulacero] at (\a:1.1) {$+$};}
  \draw[verdeteo, densely dashed, line width=0.7pt] (0,0) circle (0.55);
  \draw[verdeteo, densely dashed, line width=0.7pt] (0,0) circle (1.9);
  \foreach \a in {0,45,...,315}{\draw[-{Stealth[length=4pt]}, cm1, line width=0.6pt] (\a:1.9) -- (\a:2.35);}
  \node[font=\scriptsize, text=azulprof] at (0,0) {$E=0$};
  \draw[grispizarra, ->, line width=0.4pt] (0,0) -- (-120:1.1) node[midway, left, font=\scriptsize, text=azulprof] {$a$};
\end{tikzpicture}\\[2pt]{\small\scshape (a)}
\end{minipage}\hfill
\begin{minipage}[t]{0.33\textwidth}\centering
\begin{tikzpicture}[baseline=(current bounding box.north)]
  \begin{axis}[tema, width=\textwidth, height=4.2cm, xlabel={$r/a$}, ylabel={$E/E(a^{+})$}, xmin=0, xmax=3, ymin=0, ymax=1.15,
     xtick={0,1,2,3}, ytick={0,0.5,1}, domain=0:3, samples=120]
    \addplot[azulprof, line width=1.1pt, domain=0:1, samples=2] {0};
    \addplot[azulprof, line width=1.1pt, domain=1:3] {1/x^2};
    \addplot[azulprof, line width=0.8pt, densely dotted] coordinates {(1,0) (1,1)};
    \node[font=\scriptsize, text=azulprof, anchor=south west] at (axis cs:1.3,0.5) {$\propto 1/r^{2}$};
    \node[font=\scriptsize, text=grispizarra, anchor=west] at (axis cs:1.05,0.12) {salto $\sigma/\eo$};
  \end{axis}
\end{tikzpicture}\\[2pt]{\small\scshape (b)}
\end{minipage}\hfill
\begin{minipage}[t]{0.33\textwidth}\centering
\begin{tikzpicture}[baseline=(current bounding box.north)]
  \begin{axis}[tema, width=\textwidth, height=4.2cm, xlabel={$r/a$}, ylabel={$\phi/\phi(a)$}, xmin=0, xmax=3, ymin=0, ymax=1.3,
     xtick={0,1,2,3}, ytick={0,0.5,1}, domain=0:3, samples=120]
    \addplot[azulacero, line width=1.1pt, domain=0:1, samples=2] {1};
    \addplot[azulacero, line width=1.1pt, domain=1:3] {1/x};
    \node[font=\scriptsize, text=azulacero, anchor=south west] at (axis cs:1.3,0.55) {$\propto 1/r$};
  \end{axis}
\end{tikzpicture}\\[2pt]{\small\scshape (c)}
\end{minipage}
\caption{Problema 1.9. (a)~Cáscara esférica de radio $a$ con carga $Q$ y las superficies gaussianas interior y exterior. (b)~El campo es nulo dentro y coulombiano fuera; salta en $r=a$ en
$\sigma/\eo=Q/(4\pi\eo a^{2})$. (c)~El potencial es constante dentro ($Q/4\pi\eo a$), decae como $1/r$ fuera y es continuo.}
\label{fig:sem-cascara}
\end{figure}


% =========================================================================
\subsection{Problema 1.10. Cilindro infinito con densidad volúmica de carga}\label{sem:cilindro}
% =========================================================================
\begin{ejercicio}[Cilindro cargado]\label{ejer:cilindro}
Hallar el campo y el potencial asociados a una distribución uniforme de carga con densidad $\rho$ en un cilindro infinitamente largo de radio $a$. Calcular la diferencia de potencial entre el eje ($r=0$) y el borde ($r=a$) del cilindro.
\end{ejercicio}

\begin{resolucion}
\textbf{Idea.} Simetría cilíndrica, como en el hilo, pero con la carga repartida por el volumen. Gaussiana: cilindro coaxial de radio $r$ y altura $h$ (\figpanel{fig:sem-cilindro}{a}). Flujo: solo la superficie lateral, $2\pi r h\,E$.

\textbf{Campo.} Exterior ($r>a$): $\Qenc=\rho\,\pi a^{2}h$, luego $2\pi rh\,E=\rho\pi a^{2}h/\eo$. Interior ($r<a$): $\Qenc=\rho\,\pi r^{2}h$ y $2\pi rh\,E=\rho\pi r^{2}h/\eo$:
\[
  \boxed{\;\vE(\vc{r})=\begin{cases}\dfrac{\rho\,r}{2\eo}\,\urho,&r<a,\\[8pt]\dfrac{\rho\,a^{2}}{2\eo\,r}\,\urho,&r>a.\end{cases}\;}
\]
Coinciden en $r=a$ ($E=\rho a/2\eo$). Fuera, es el del hilo infinito con $\lambda=\rho\pi a^{2}$ (el cilindro se ve como una línea).

\textbf{Potencial.} Al igual que en el hilo infinito, la carga llega al infinito y no se puede tomar $\phi(\infty)=0$; tomamos $\phi(r_0)=0$ con $r_0>a$. Exterior ($r\ge a$):
\[
  \phi(r)=-\int_{r_0}^{r}\frac{\rho a^{2}}{2\eo r'}\,\dif r'=\frac{\rho a^{2}}{2\eo}\,\ln\frac{r_0}{r}.
\]
En el borde, $\phi(a)=\frac{\rho a^{2}}{2\eo}\ln\frac{r_0}{a}$. Interior ($r<a$), partiendo de $\phi(a)$ (no podemos integrar desde $r_0$ de una vez porque la expresión del campo cambia en $r=a$):
\[
  \phi(r)=\phi(a)-\int_{a}^{r}\frac{\rho\,r'}{2\eo}\,\dif r'=\frac{\rho a^{2}}{2\eo}\ln\frac{r_0}{a}-\frac{\rho}{4\eo}\,(r^{2}-a^{2}),
\]
es decir,
\[
  \boxed{\;\phi(r)=\begin{cases}-\dfrac{\rho}{4\eo}\Bigl[\,r^{2}-a^{2}-2a^{2}\ln\dfrac{r_0}{a}\Bigr],&r\le a,\\[8pt]\dfrac{\rho\,a^{2}}{2\eo}\,\ln\dfrac{r_0}{r},&r\ge a.\end{cases}\;}
\]
\textbf{Diferencia de potencial entre el eje y el borde.} Del resultado interior, $\phi(0)=\frac{\rho}{4\eo}\bigl(a^{2}+2a^{2}\ln\frac{r_0}{a}\bigr)$ y $\phi(a)=\frac{\rho}{4\eo}\,2a^{2}\ln\frac{r_0}{a}$, y por tanto
\[
  \boxed{\;\phi(0)-\phi(a)=\frac{\rho\,a^{2}}{4\eo}\;}
\]
independiente de $r_0$ (la referencia se cancela en toda diferencia de potencial; directamente, $\phi(0)-\phi(a)=\int_0^a E\,\dif r=\int_0^a\frac{\rho r}{2\eo}\dif r$).
\end{resolucion}

\begin{observacion}[La referencia del potencial]
Como la carga se extiende hasta el infinito, el potencial depende de una elección arbitraria de $r_0$. La \figpanel{fig:sem-cilindro}{c} lo muestra para dos referencias, $\phi(r_0)=0$ y $\phi(0)=0$: las curvas están desplazadas una respecto de la otra por una constante, pero tienen la misma forma, el mismo campo y la misma diferencia de potencial. Eso es lo físicamente relevante.
\end{observacion}

% Figura (Problema 1.10): cilindro infinito cargado y superficies gaussianas coaxiales (r>a y r<a).
\begin{figure}[ht]
\centering
\tikzset{
  ejeC/.style={-{Stealth[length=5pt]}, line width=0.6pt, black!80},
  cotaC/.style={arrows={Stealth[length=4pt]-Stealth[length=4pt]}, line width=0.5pt, black!85},
  fondoB/.style={fill=white, fill opacity=0.85, text opacity=1, inner sep=1pt, rounded corners=1pt}}
% \cilCarga{rx}{ry}{y0}{y1}: cilindro cargado (sombreado)
\newcommand{\cilCarga}[4]{%
  \shade[left color=azulacero!55, right color=azulclaro!55, middle color=azulclaro!30]
     (-#1,#3) -- (-#1,#4) arc (180:360:{#1} and #2) -- (#1,#3) arc (0:-180:{#1} and #2) -- cycle;
  \draw[azulprof, line width=0.6pt] (-#1,#3) -- (-#1,#4) (#1,#3) -- (#1,#4);
  \draw[azulprof, line width=0.6pt] (-#1,#3) arc (180:360:{#1} and #2);
  \fill[azulclaro!45] (0,#4) ellipse ({#1} and #2);
  \draw[azulprof, line width=0.6pt] (0,#4) ellipse ({#1} and #2);}
% \cilGaussAtras / \cilGaussDelante{rx}{ry}{y0}{y1}
\newcommand{\cilGaussAtras}[4]{%
  \draw[verdeteo, line width=0.5pt, densely dashed] (-#1,#3) arc (180:0:{#1} and #2);
  \draw[verdeteo, line width=0.5pt, densely dashed] (-#1,#4) arc (180:0:{#1} and #2);}
\newcommand{\cilGaussDelante}[4]{%
  \fill[verdeteo, opacity=0.08] (-#1,#3) -- (-#1,#4) arc (180:360:{#1} and #2) -- (#1,#3) arc (0:-180:{#1} and #2) -- cycle;
  \draw[verdeteo, line width=0.8pt] (-#1,#3) -- (-#1,#4) (#1,#3) -- (#1,#4);
  \draw[verdeteo, line width=0.8pt] (-#1,#3) arc (180:360:{#1} and #2);
  \draw[verdeteo, line width=0.8pt] (-#1,#4) arc (180:360:{#1} and #2);}
% ---------------------------------------------------------------- (a) r > a
\begin{minipage}[b]{0.5\textwidth}\centering
\begin{tikzpicture}[x=1cm, y=1cm]
  \def\ra{0.62}\def\rya{0.17}\def\rg{1.55}\def\ryg{0.42}
  \cilGaussAtras{\rg}{\ryg}{-1.25}{1.25}
  \cilCarga{\ra}{\rya}{-2.3}{2.1}
  % ejes (el eje z sale de la tapa superior)
  \draw[ejeC] (0,0) -- (0,2.75) node[above, font=\small] {$z$};
  \draw[ejeC] (0,0) -- (2.15,-0.23) node[below right=-1pt, font=\small, fill=white, fill opacity=0.85, text opacity=1, inner sep=1pt] {$y$};
  \draw[ejeC] (0,0) -- (-1.15,-0.85) node[below left=-1pt, font=\small, fill=white, fill opacity=0.85, text opacity=1, inner sep=1pt] {$x$};
  \cilGaussDelante{\rg}{\ryg}{-1.25}{1.25}
  % radios a y r, altura h
  \draw[-{Stealth[length=4.5pt]}, line width=0.8pt, azulprof] (0,2.1) -- (\ra,2.1);
  \node[font=\small, text=azulprof, anchor=south] at (0.31,2.12) {$a$};
  \draw[-{Stealth[length=4.5pt]}, line width=0.8pt, verdeteo] (0,1.25) -- (\rg,1.25);
  \node[font=\small, text=verdeteo, fondoB, anchor=south] at (1.15,1.27) {$r$};
  \draw[cotaC] (\rg+0.85,-1.25) -- (\rg+0.85,1.25) node[midway, right, font=\small] {$h$};
  \node[font=\small, text=azulprof, fondoB] at (-0.32,-1.75) {$\rho$};
  \node[font=\small, text=verdeteo, anchor=east] at (-\rg-0.05,1.35) {$S_{\mathrm{G}}$};
\end{tikzpicture}\\[2pt]{\small\scshape (a)\; $r>a$}
\end{minipage}\hfill
% ---------------------------------------------------------------- (b) r < a
\begin{minipage}[b]{0.46\textwidth}\centering
\begin{tikzpicture}[x=1cm, y=1cm]
  \def\ra{1.15}\def\rya{0.3}\def\rg{0.62}\def\ryg{0.17}
  \cilCarga{\ra}{\rya}{-2.3}{2.1}
  \cilGaussAtras{\rg}{\ryg}{-1.1}{1.1}
  \draw[ejeC] (0,0) -- (0,2.75) node[above, font=\small] {$z$};
  \draw[ejeC] (0,0) -- (2.3,-0.25) node[right, font=\small] {$y$};
  \draw[ejeC] (0,0) -- (-1.15,-0.85) node[below left=-1pt, font=\small, fill=white, fill opacity=0.85, text opacity=1, inner sep=1pt] {$x$};
  \cilGaussDelante{\rg}{\ryg}{-1.1}{1.1}
  \draw[-{Stealth[length=4.5pt]}, line width=0.8pt, azulprof] (0,2.1) -- (\ra,2.1);
  \node[font=\small, text=azulprof, anchor=south] at (0.6,2.12) {$a$};
  \draw[-{Stealth[length=4.5pt]}, line width=0.8pt, verdeteo] (0,1.1) -- (\rg,1.1);
  \node[font=\small, text=verdeteo, fondoB, anchor=south] at (0.36,1.12) {$r$};
  \draw[cotaC] (\ra+0.35,-1.1) -- (\ra+0.35,1.1) node[midway, right, font=\small] {$h$};
  \draw[grispizarra, line width=0.3pt] (\rg+0.05,0.55) -- (\ra+0.35,0.55);
  \draw[grispizarra, line width=0.3pt] (\rg+0.05,-0.55) -- (\ra+0.35,-0.55);
  \node[font=\small, text=azulprof, fondoB] at (-0.85,-1.75) {$\rho$};
  \node[font=\small, text=verdeteo, fondoB, anchor=east] at (-\rg-0.08,0.6) {$S_{\mathrm{G}}$};
\end{tikzpicture}\\[2pt]{\small\scshape (b)\; $r<a$}
\end{minipage}
\caption{Problema 1.10. Cilindro infinito de radio $a$ con densidad volúmica uniforme $\rho$ (en azul) y superficie gaussiana coaxial $S_{\mathrm{G}}$ (en verde) de radio $r$ y
altura $h$: (a)~exterior ($r>a$), que encierra $\rho\pi a^{2}h$; (b)~interior ($r<a$), que encierra $\rho\pi r^{2}h$. Solo la superficie lateral contribuye al flujo.}
\label{fig:sem-cilindro}
\end{figure}

\begin{figure}[ht]
\centering
\begin{tikzpicture}
  \begin{axis}[tema, name=EA, width=0.47\textwidth, height=5.0cm, xlabel={$r/a$}, ylabel={$E\,\big/\,(\rho a/2\eo)$}, xmin=0, xmax=6, ymin=0, ymax=1.15,
     xtick={0,1,2,3,4,5,6}, ytick={0,0.2,0.4,0.6,0.8,1}, domain=0:6, samples=150]
    \addplot[azulacero, line width=1.4pt, domain=0:1] {x};
    \addplot[azulacero, line width=1.4pt, domain=1:6] {1/x};
    \addplot[grispizarra, densely dashed] coordinates {(1,0) (1,1)};
    \node[font=\small, text=azulacero, anchor=south east] at (axis cs:0.95,0.08) {$\propto r$};
    \node[font=\small, text=azulacero, anchor=south west] at (axis cs:2.2,0.5) {$\propto 1/r$};
  \end{axis}
  \begin{axis}[tema, at={(EA.east)}, anchor=west, xshift=1.6cm, width=0.47\textwidth, height=5.0cm, xlabel={$r/a$}, ylabel={$\phi\,\big/\,(\rho a^{2}/2\eo)$},
     xmin=0, xmax=6, ymin=-2.5, ymax=2.0, xtick={0,1,2,3,4,5,6}, ytick={-2,-1,0,1}, domain=0:6, samples=150]
    \addplot[grispizarra, densely dashed] coordinates {(1,-2.5) (1,2.0)};
    % phi(r0)=0 con r0=3a
    \addplot[azulacero, line width=1.4pt, domain=0.001:1] {0.5*(1-x^2)+ln(3)};
    \addplot[azulacero, line width=1.4pt, domain=1:6] {ln(3/x)};
    % phi(0)=0
    \addplot[acento, line width=1.4pt, domain=0.001:1] {-x^2/2};
    \addplot[acento, line width=1.4pt, domain=1:6] {-ln(x)-0.5};
    \node[font=\small, text=azulacero, anchor=south west] at (axis cs:2.3,0.25) {$\phi(r_0)=0$};
    \node[font=\small, text=acento, anchor=south east] at (axis cs:5.9,-1.95) {$\phi(0)=0$};
  \end{axis}
\end{tikzpicture}
\caption{Problema 1.10. (Izquierda) Módulo del campo: crece linealmente en el interior del cilindro y decae como $1/r$ fuera. (Derecha) Potencial con dos elecciones de
referencia, $\phi(r_0)=0$ con $r_0=3a$ (azul) y $\phi(0)=0$ (ámbar). Al extenderse la carga hasta el infinito no puede tomarse $\phi(\infty)=0$; las dos curvas difieren en una
constante, y la diferencia $\phi(0)-\phi(a)=\rho a^{2}/4\eo$ no depende de la referencia.}
\label{fig:sem-cilindro-graf}
\end{figure}


% =========================================================================
\subsection{Problema 1.11. Tres planos ortogonales cargados}\label{sem:planos}
% =========================================================================
\begin{ejercicio}[Tres planos]\label{ejer:planos}
Calcular el campo creado por tres planos ortogonales ($x=0$, $y=0$ y $z=0$) infinitamente grandes y cargados uniformemente con densidades superficiales $\sigma_a$, $\sigma_b$ y $\sigma_c$.
\end{ejercicio}

\begin{resolucion}
\textbf{Idea.} Es un ejercicio de \emph{superposición}: la fuerza (y el campo) es lineal en las cargas, así que el campo total es la suma de los campos de cada plano, tomados como si los demás no existieran (\cref{def:superposicion}). Cada uno es el
del plano infinito (\cref{sem:plano}) con la normal adecuada:
\[
  \vE_a=\frac{\sigma_a}{2\eo}\,\frac{x}{\lvert x\rvert}\,\ux,\qquad\vE_b=\frac{\sigma_b}{2\eo}\,\frac{y}{\lvert y\rvert}\,\uy,\qquad\vE_c=\frac{\sigma_c}{2\eo}\,\frac{z}{\lvert z\rvert}\,\uz .
\]
Sumando,
\[
  \boxed{\;\vE(\vc{r})=\frac{1}{2\eo}\Bigl(\sigma_a\,\frac{x}{\lvert x\rvert}\,\ux+\sigma_b\,\frac{y}{\lvert y\rvert}\,\uy+\sigma_c\,\frac{z}{\lvert z\rvert}\,\uz\Bigr)\;}
\]
\end{resolucion}

\begin{observacion}[Interpretación]
El espacio queda dividido en ocho octantes, y en cada uno el campo es \emph{uniforme} y cambia de signo en las componentes que atraviesan un plano. Si las tres densidades son iguales y positivas, el campo es siempre saliente: de módulo $\frac{\sqrt3\,\sigma}{2\eo}$ y dirección
$(\pm\ux\pm\uy\pm\uz)/\sqrt3$ en cada octante (\figpanel{fig:sem-planos}{b}). Cada componente tiene su propio salto $\sigma_i/\eo$ al cruzar su plano, como exigen las condiciones de contorno: las componentes tangenciales a un plano (las de los otros dos) son continuas a su paso.
\end{observacion}

% Figura (Problema 1.11): tres planos ortogonales cargados; campo por superposición.
\begin{figure}[ht]
\centering
\tdplotsetmaincoords{62}{128}
\begin{minipage}[t]{0.48\textwidth}\centering
\begin{tikzpicture}[tdplot_main_coords, scale=0.88, baseline=(current bounding box.north)]
  \def\L{2.3}
  \fill[azulclaro, opacity=0.30] (0,-\L,-\L) -- (0,\L,-\L) -- (0,\L,\L) -- (0,-\L,\L) -- cycle;   % x=0
  \fill[verdeteo, opacity=0.18] (-\L,0,-\L) -- (\L,0,-\L) -- (\L,0,\L) -- (-\L,0,\L) -- cycle;   % y=0
  \fill[acento, opacity=0.20] (-\L,-\L,0) -- (\L,-\L,0) -- (\L,\L,0) -- (-\L,\L,0) -- cycle;      % z=0
  \draw[azulacero, line width=0.4pt] (0,-\L,-\L) -- (0,\L,-\L) -- (0,\L,\L) -- (0,-\L,\L) -- cycle;
  \draw[verdeteo, line width=0.4pt] (-\L,0,-\L) -- (\L,0,-\L) -- (\L,0,\L) -- (-\L,0,\L) -- cycle;
  \draw[acento, line width=0.4pt] (-\L,-\L,0) -- (\L,-\L,0) -- (\L,\L,0) -- (-\L,\L,0) -- cycle;
  \draw[eje] (0,0,0) -- (3.0,0,0) node[below left] {\small $x$};
  \draw[eje] (0,0,0) -- (0,3.0,0) node[right] {\small $y$};
  \draw[eje] (0,0,0) -- (0,0,3.0) node[above] {\small $z$};
  \node[font=\scriptsize, text=azulprof] at (0,-1.6,1.9) {$\sigma_a$};
  \node[font=\scriptsize, text=verdeteo] at (-1.9,0,1.9) {$\sigma_b$};
  \node[font=\scriptsize, text=acento] at (1.9,1.9,0) {$\sigma_c$};
\end{tikzpicture}\\[2pt]{\small\scshape (a)}
\end{minipage}\hfill
\begin{minipage}[t]{0.48\textwidth}\centering
\begin{tikzpicture}[tdplot_main_coords, scale=0.88, baseline=(current bounding box.north)]
  \draw[eje] (-2.6,0,0) -- (2.8,0,0) node[below left] {\small $x$};
  \draw[eje] (0,-2.6,0) -- (0,2.8,0) node[right] {\small $y$};
  \draw[eje] (0,0,-2.6) -- (0,0,2.8) node[above] {\small $z$};
  % un vector por octante: E = (sgn x, sgn y, sgn z)/2e0 con sigma=1 (módulo constante)
  \foreach \sx in {-1,1}{\foreach \sy in {-1,1}{\foreach \sz in {-1,1}{%
    \pgfmathsetmacro{\ft}{0.55+0.25*(\sx*\sy*\sz)}%
    \mapcolor{\ft}%
    \draw[-{Stealth[length=4.5pt,width=4.5pt]}, line width=1.3pt, mapc] ({1.0*\sx},{1.0*\sy},{1.0*\sz}) -- ++({1.0*\sx},{1.0*\sy},{1.0*\sz});}}}
\end{tikzpicture}\\[2pt]{\small\scshape (b)}
\end{minipage}
\caption{Problema 1.11. (a)~Tres planos ortogonales $x=0$, $y=0$ y $z=0$ con densidades $\sigma_a$, $\sigma_b$ y $\sigma_c$. (b)~Para $\sigma_a=\sigma_b=\sigma_c>0$ el campo, suma
de tres campos uniformes de módulo $\sigma/2\eo$ que cambian de sentido al cruzar su plano, es en cada octante $\frac{\sigma}{2\eo}(\pm\ux\pm\uy\pm\uz)$: constante a trozos y
siempre saliente.}
\label{fig:sem-planos}
\end{figure}


% =========================================================================
\subsection{Problema 1.12. Semicáscara esférica: campo en el centro}\label{sem:semiesfera}
% =========================================================================
\begin{ejercicio}[Semicáscara esférica]\label{ejer:semiesfera}
Una carga $Q$ está uniformemente distribuida sobre una superficie semiesférica de radio $a$. Hallar el campo eléctrico en su centro.
\end{ejercicio}

\begin{resolucion}
\textbf{Idea.} Es una integral directa (\cref{eq:campo-sup-lin}), sobre una superficie esférica. Colocamos el centro en el origen, el punto campo $\vc{r}=\vc{0}$, y la semicáscara en la mitad $z<0$, es decir, $\theta'\in[\pi/2,\pi]$ (\cref{fig:sem-semiesfera}).
La densidad superficial es $\sigma=Q/(2\pi a^{2})$ (media esfera: $2\pi a^{2}$).

\textbf{Geometría.} Un punto de la semicáscara: $\vc{r}\,'=a\,\ur'=a\bigl[\cos\theta'\,\uz+\sen\theta'(\cos\varphi'\,\ux+\sen\varphi'\,\uy)\bigr]$, y $\dSp=a^{2}\sen\theta'\,\dif\theta'\,\dif\varphi'$. Entonces
\[
  \vc{R}=\vc{0}-\vc{r}\,'=-a\bigl[\cos\theta'\,\uz+\sen\theta'(\cos\varphi'\,\ux+\sen\varphi'\,\uy)\bigr],\qquad R=a .
\]
\textbf{Integral.} Las componentes $x$, $y$ contienen $\cos\varphi'$, $\sen\varphi'$ y se anulan al integrar en $\varphi'$; solo sobrevive la componente $z$:
\[
  \vE_0=\frac{\sigma}{4\pi\eo}\int\frac{\vc{R}}{R^{3}}\,\dSp=\frac{\sigma}{4\pi\eo}\int_{\pi/2}^{\pi}\!\!\int_{0}^{2\pi}\frac{-a\cos\theta'}{a^{3}}\,a^{2}\sen\theta'\,\dif\varphi'\,\dif\theta'\;\uz
  =-\frac{\sigma}{2\eo}\int_{\pi/2}^{\pi}\cos\theta'\sen\theta'\,\dif\theta'\;\uz .
\]
Como $\int\cos\theta'\sen\theta'\,\dif\theta'=\tfrac12\sen^{2}\theta'$,
\[
  \vE_0=-\frac{\sigma}{2\eo}\Bigl[\tfrac12\sen^{2}\theta'\Bigr]_{\pi/2}^{\pi}\uz=-\frac{\sigma}{2\eo}\Bigl(0-\frac12\Bigr)\uz=\frac{\sigma}{4\eo}\,\uz,
\]
y con $\sigma=Q/2\pi a^{2}$,
\[
  \boxed{\;\vE_0=\frac{\sigma}{4\eo}\,\uz=\frac{Q}{8\pi\eo\,a^{2}}\,\uz\;}
\]
\end{resolucion}

\begin{observacion}[Comprobaciones]
El campo apunta en $+\uz$, alejándose de la carga, como debe: la semicáscara está en $z<0$ y empuja una carga positiva de prueba hacia arriba. Un chequeo independiente: una cáscara esférica \emph{completa} no crea campo en su interior (\cref{sem:cascara}), de modo que la semicáscara
superior debe crear en el centro el campo opuesto, $-\frac{\sigma}{4\eo}\uz$: las dos mitades se cancelan. Además, el valor es justo la mitad del que crearía $Q$ concentrada a distancia $a$ \emph{si} las contribuciones se sumasen sin proyectar, $Q/(4\pi\eo a^{2})$, multiplicada por el
factor geométrico $\langle\cos\theta\rangle=1/2$ del promedio de la proyección sobre el eje: $\frac{Q}{4\pi\eo a^{2}}\cdot\frac12$.
\end{observacion}

% Figura (Problema 1.12): semiesfera cargada; campo en el centro.
\begin{figure}[ht]
\centering
\tdplotsetmaincoords{66}{125}
\begin{tikzpicture}[tdplot_main_coords, scale=1.25]
  \def\RR{1.8}
  % semicáscara inferior: casquete sombreado, borde ecuatorial y meridianos
  \begin{scope}
    \clip plot[domain=\VA-180:\VA, samples=40, variable=\tdt] ({\RR*cos(\tdt)},{\RR*sin(\tdt)},0) -- plot[domain=\VA:\VA-180, samples=2, variable=\tdt] ({\RR*cos(\tdt)},{\RR*sin(\tdt)},0) -- (0,0,-\RR*1.01) -- cycle;
  \end{scope}
  \shade[left color=azulclaro!30, right color=azulclaro!80, opacity=0.85]
    plot[domain=0:180, samples=60, variable=\tdt] ({\RR*cos(\VA)*cos(\tdt)},{\RR*sin(\VA)*cos(\tdt)},{-\RR*sin(\tdt)})
    -- plot[domain=\VA-180:\VA, samples=40, variable=\tdt] ({\RR*cos(\tdt)},{\RR*sin(\tdt)},0) -- cycle;
  \draw[azulprof, line width=0.8pt] plot[domain=0:180, samples=60, variable=\tdt] ({\RR*cos(\VA)*cos(\tdt)},{\RR*sin(\VA)*cos(\tdt)},{-\RR*sin(\tdt)});
  \tdCirc[azulprof, line width=0.8pt]{\RR}{0}
  \foreach \mm in {35,80,170,215,260,305}{
    \draw[azulprof!45, line width=0.4pt, domain=0:90, samples=30, variable=\tdt] plot ({\RR*cos(\mm)*sin(\tdt)},{\RR*sin(\mm)*sin(\tdt)},{-\RR*cos(\tdt)});}
  \draw[eje] (0,0,-2.3) -- (0,0,0); \draw[eje] (0,0,0) -- (0,0,2.6) node[above] {\small $z$};
  \draw[eje] (0,0,0) -- (2.9,0,0) node[below left] {\small $x$};
  \draw[eje] (0,0,0) -- (0,2.9,0) node[right] {\small $y$};
  \def\ph{60}\def\th{140}
  \coordinate (S) at ({\RR*sin(\th)*cos(\ph)},{\RR*sin(\th)*sin(\ph)},{\RR*cos(\th)});
  \draw[vect, azulprof, line width=1.0pt] (S) -- (0,0,0);
  \node[font=\small, text=azulprof, anchor=west] at ($(S)!0.5!(0,0,0)+(0.06,0,0)$) {$\vc{R}=-\vc{r}\,'$};
  \draw[vectrojo, line width=1.0pt] (0,0,0) -- ({0.75*sin(\th)*cos(\ph)*(-1)},{0.75*sin(\th)*sin(\ph)*(-1)},{0.75*cos(\th)*(-1)}) node[left=1pt, font=\scriptsize, text=acento] {$\dif\vE$};
  \draw[vectrojo, line width=1.7pt] (0,0,0) -- (0,0,1.7) node[left, font=\small, text=acento] {$\vE_{0}$};
  \tdParche[acento, fill=acento!45, line width=0.6pt]{\RR*sin(\th)*cos(\ph)*0.985,\RR*sin(\th)*sin(\ph)*0.985,\RR*cos(\th)*0.985}{0.2,-0.1,0}{0.0,0.0,0.2}
  \fill[acento] (S) circle (1.2pt);
  \node[font=\small, text=acento, anchor=north west] at ($(S)+(0.05,0.05,-0.12)$) {$\dif S'$};
  \node[font=\small, text=azulprof] at (1.1,-0.5,-1.3) {$\sigma$};
  \fill[azulprof] (0,0,0) circle (1.4pt) node[above right=1pt, font=\small] {$O$};
\end{tikzpicture}
\caption{Problema 1.12. Semicáscara esférica de radio $a$ (la mitad $z<0$) con carga $Q$ uniforme. El campo creado en el centro $O$ por un elemento $\dif S'$ apunta
según $\vc{R}=-\vc{r}\,'$; las componentes horizontales se cancelan por simetría y el campo total $\vE_0$ es axial (positivo en $z$).}
\label{fig:sem-semiesfera}
\end{figure}


% -------------------------------------------------------------------------
\needspace{22\baselineskip}
\subsection{Resumen de resultados}\label{sec:resumen-seminarios}
% -------------------------------------------------------------------------
\begin{center}
\captionof{table}{Resultados de los problemas del tema ($\vE$ en V/m; $\phi$ en V).}\label{tab:resultados}
\footnotesize
\renewcommand{\arraystretch}{2.3}
\begin{tabular}{@{}p{1.0cm}p{2.4cm}p{6.2cm}p{4.6cm}@{}}
\toprule
Prob. & Distribución & Campo $\vE$ & Potencial $\phi$\\
\midrule
1.1, 1.2 & Esfera maciza, $Q$, $a$ &
  $\dfrac{Q\,r}{4\pi\eo a^{3}}\,\ur$\ \ ($r<a$)\newline $\dfrac{Q}{4\pi\eo r^{2}}\,\ur$\ \ ($r>a$) &
  $\dfrac{Q}{8\pi\eo a}\Bigl(3-\dfrac{r^{2}}{a^{2}}\Bigr)$\ \ ($r<a$)\newline $\dfrac{Q}{4\pi\eo r}$\ \ ($r>a$)\\
1.3, 1.4 & Hilo infinito, $\lambda$ & $\dfrac{\lambda}{2\pi\eo r}\,\urho$ & $\dfrac{\lambda}{2\pi\eo}\ln\dfrac{r_0}{r}$\\
1.5, 1.6 & Plano infinito, $\sigma$ & $\dfrac{\sigma}{2\eo}\,\sgn(z)\,\uz$ & $-\dfrac{\sigma}{2\eo}\lvert z\rvert$\\
1.7 & Anillo, $Q$, $a$ (eje) & $\dfrac{Qz}{4\pi\eo(a^{2}+z^{2})^{3/2}}\,\uz$ & $\dfrac{Q}{4\pi\eo\sqrt{a^{2}+z^{2}}}$\\
1.7 & Disco, $\sigma$, $a$ (eje) & $\dfrac{\sigma}{2\eo}\Bigl(\sgn z-\dfrac{z}{\sqrt{a^{2}+z^{2}}}\Bigr)\uz$ & $\dfrac{\sigma}{2\eo}\bigl(\sqrt{a^{2}+z^{2}}-\lvert z\rvert\bigr)$\\
1.8 & Lámina, $\rho$, espesor $a$ & $\dfrac{\rho a}{2\eo}\sgn(z)\,\uz$\ \ ($\lvert z\rvert>\tfrac a2$)\newline $\dfrac{\rho z}{\eo}\,\uz$\ \ ($\lvert z\rvert<\tfrac a2$) & ---\\
1.9 & Cáscara esférica, $Q$, $a$ & $\vc{0}$\ \ ($r<a$)\newline $\dfrac{Q}{4\pi\eo r^{2}}\,\ur$\ \ ($r>a$) & $\dfrac{Q}{4\pi\eo a}$\ \ ($r\le a$)\newline $\dfrac{Q}{4\pi\eo r}$\ \ ($r\ge a$)\\
1.10 & Cilindro, $\rho$, $a$ & $\dfrac{\rho r}{2\eo}\,\urho$\ \ ($r<a$)\newline $\dfrac{\rho a^{2}}{2\eo r}\,\urho$\ \ ($r>a$) & $\phi(0)-\phi(a)=\dfrac{\rho a^{2}}{4\eo}$\\
1.11 & Tres planos & $\dfrac{1}{2\eo}\bigl(\sigma_a\sgn x\,\ux+\sigma_b\sgn y\,\uy+\sigma_c\sgn z\,\uz\bigr)$ & ---\\
1.12 & Semicáscara, $Q$, $a$ (centro) & $\dfrac{Q}{8\pi\eo a^{2}}\,\uz$ & ---\\
\bottomrule
\end{tabular}
\end{center}

% =========================================================================
\clearpage
\appendix
\section{Cálculo diferencial e integral: significado físico}\label{app:calculo}
% =========================================================================
A lo largo del tema hemos usado el gradiente, la divergencia y el rotacional, definidos y estudiados en el Tema~0~\cite{TEMA0}. Aquí nos detenemos en su \emph{significado físico}, que es el que
da sentido a las ecuaciones del tema, y en dos identidades que se usan continuamente.

% -------------------------------------------------------------------------
\subsection{El gradiente}\label{app:gradiente}
% -------------------------------------------------------------------------
El potencial eléctrico es un campo escalar: a cada punto del espacio le corresponde un valor. Su variación a lo largo de un camino es, en cartesianas,
\[
  \dif\phi=\pd{\phi}{x}\,\dif x+\pd{\phi}{y}\,\dif y+\pd{\phi}{z}\,\dif z=\Bigl(\pd{\phi}{x}\,\ux+\pd{\phi}{y}\,\uy+\pd{\phi}{z}\,\uz\Bigr)\cdot\bigl(\dif x\,\ux+\dif y\,\uy+\dif z\,\uz\bigr)=\nab\phi\cdot\dif\vc{r},
\]
producto escalar del gradiente por el diferencial del vector de posición. Por la definición del producto escalar, $\dif\phi=\norm{\nab\phi}\,\dif r\cos\theta$, con $\theta$ el ángulo entre $\nab\phi$ y $\dif\vc{r}$.

\begin{proposicion}[Significado del gradiente]\label{prop:grad-maximo}
Fijado el módulo $\dif r$ del desplazamiento, la variación $\dif\phi$ es máxima cuando $\theta=0$, es decir, cuando $\dif\vc{r}\parallel\nab\phi$. Por tanto, $\nab\phi$ es el vector que, en cada punto, apunta en la dirección y sentido de máximo crecimiento de
$\phi$, y cuyo módulo es esa tasa máxima de variación, $\norm{\nab\phi}$.
\end{proposicion}
\begin{proof}
$\dif\phi/\dif r=\norm{\nab\phi}\cos\theta\le\norm{\nab\phi}$, con igualdad si y solo si $\cos\theta=1$.
\end{proof}
De aquí que $\vE=-\nab\phi$ apunte hacia donde el potencial decrece más deprisa, y que las líneas de campo sean perpendiculares a las equipotenciales ($\cos\theta=0$ para $\dif\phi=0$).

% -------------------------------------------------------------------------
\subsection{La divergencia}\label{app:divergencia}
% -------------------------------------------------------------------------
La divergencia de un campo vectorial es una magnitud escalar; en cartesianas,
$\nab\cdot\vE=\partial_xE_x+\partial_yE_y+\partial_zE_z$. El nombre está bien escogido: mide cómo el campo \emph{diverge} (se expande) o \emph{converge} (se encoge) en un punto (\cref{fig:div3}).
Dice, por tanto, si en ese punto hay \emph{fuentes} (divergencia positiva) o \emph{sumideros} (negativa).

% Figura: campos con divergencia negativa, nula y positiva.
\begin{figure}[ht]
\centering
\begin{minipage}[t]{0.325\textwidth}\centering
\begin{tikzpicture}[x=0.68cm, y=0.68cm, baseline=(current bounding box.north)]
  \marcoD{-2.7}{2.7}{-2.7}{2.7}{-2,-1,0,1,2}{-2,-1,0,1,2}
  \ejesD{-2.7}{2.7}{-2.7}{2.7}
  \clipD{-2.7}{2.7}{-2.7}{2.7}{\flechas{-\x}{-\y}{-2.4,-1.8,...,2.4}{-2.4,-1.8,...,2.4}{0.20}{3.4}}
\end{tikzpicture}\\[3pt]{\small (a)\; $\nab\cdot\vc{A}<0$}
\end{minipage}\hfill
\begin{minipage}[t]{0.325\textwidth}\centering
\begin{tikzpicture}[x=0.68cm, y=0.68cm, baseline=(current bounding box.north)]
  \marcoD{-2.7}{2.7}{-2.7}{2.7}{-2,-1,0,1,2}{-2,-1,0,1,2}
  \ejesD{-2.7}{2.7}{-2.7}{2.7}
  \clipD{-2.7}{2.7}{-2.7}{2.7}{\flechas{1}{0}{-2.4,-1.8,...,2.4}{-2.4,-1.8,...,2.4}{0.34}{1.2}}
\end{tikzpicture}\\[3pt]{\small (b)\; $\nab\cdot\vc{A}=0$}
\end{minipage}\hfill
\begin{minipage}[t]{0.325\textwidth}\centering
\begin{tikzpicture}[x=0.68cm, y=0.68cm, baseline=(current bounding box.north)]
  \marcoD{-2.7}{2.7}{-2.7}{2.7}{-2,-1,0,1,2}{-2,-1,0,1,2}
  \ejesD{-2.7}{2.7}{-2.7}{2.7}
  \clipD{-2.7}{2.7}{-2.7}{2.7}{\flechas{\x}{\y}{-2.4,-1.8,...,2.4}{-2.4,-1.8,...,2.4}{0.20}{3.4}}
\end{tikzpicture}\\[3pt]{\small (c)\; $\nab\cdot\vc{A}>0$}
\end{minipage}
\caption{Significado de la divergencia. (a)~$\vc{A}=-\vc{r}$: las flechas convergen hacia el origen (sumidero, divergencia $-2$). (b)~Campo uniforme:
lo que entra por un lado sale por el otro (divergencia nula). (c)~$\vc{A}=\vc{r}$: el campo se expande desde el origen (fuente, divergencia $+2$). El color codifica $\norm{\vc{A}}$.}
\label{fig:div3}
\end{figure}


Otra forma de entender su significado es mediante el teorema de la divergencia, $\iiint_{V}\nab\cdot\vE\,\dV=\oiint_{S}\vE\cdot\dS$: la integral de la divergencia en un volumen es el flujo a través de la superficie que lo limita. Si el flujo es positivo, sale más de lo que entra y hay fuentes encerradas;
en caso contrario hay sumideros. Dividiendo por el volumen y haciéndolo tender a cero,
\begin{equation}\label{eq:div-limite}
  \nab\cdot\vE(\vc{r})=\lim_{V\to\{\vc{r}\}}\frac{1}{V}\oiint_{\partial V}\vE\cdot\dS,
\end{equation}
la divergencia es el flujo neto por unidad de volumen, es decir, una medida de que en un punto haya fuentes o sumideros.

\begin{ejemplo}[Un caso que parece contradecir el teorema de la divergencia]\label{ej:div-r2}
Comprobar que el campo $\vc{A}=\dfrac{1}{r^{2}}\,\ur$ verifica el teorema de la divergencia en una esfera centrada en el origen.

\emph{Flujo.} En esféricas, con $\dS=r^{2}\sen\theta\,\dif\theta\,\dif\varphi\,\ur$,
\[
  \oiint_{S}\vc{A}\cdot\dS=\int_{0}^{\pi}\!\!\int_{0}^{2\pi}\frac{1}{r^{2}}\,r^{2}\sen\theta\,\dif\varphi\,\dif\theta=2\pi\int_{0}^{\pi}\sen\theta\,\dif\theta=4\pi .
\]
\emph{Divergencia, a primera vista.} $\nab\cdot\vc{A}=\frac{1}{r^{2}}\pd{}{r}\bigl(r^{2}\cdot\frac{1}{r^{2}}\bigr)=\frac{1}{r^{2}}\pd{(1)}{r}=0$, de modo que $\iiint_{V}\nab\cdot\vc{A}\,\dV=0$ y el teorema parece no cumplirse: $0\overset{?}{=}4\pi$.
Resulta extraño obtener una divergencia nula para un campo cuyas líneas, obviamente, se expanden desde el origen.

\emph{Resolución.} El cálculo olvida que el campo \emph{diverge} en $r=0$: $\nab\cdot\vc{A}=0$ en todos los puntos salvo en $r=0$, donde diverge. Como el volumen de integración contiene el origen, no se puede sustituir la divergencia por cero. Lo correcto es usar la delta de Dirac (\cref{teo:delta}):
$\nab\cdot\bigl(\ur/r^{2}\bigr)=4\pi\,\delta^{3}(\vc{r})$. Entonces
\[
  \iiint_{V}\nab\cdot\vc{A}\,\dV=\iiint_{V}4\pi\,\delta^{3}(\vc{r})\,\dV=4\pi=\oiint_{S}\vc{A}\cdot\dS\ \checkmark
\]
Es, de nuevo, el contenido del campo de una carga puntual.
\end{ejemplo}

% -------------------------------------------------------------------------
\subsection{El rotacional}\label{app:rotacional}
% -------------------------------------------------------------------------
El rotacional es una magnitud vectorial; en cartesianas,
\[
  \nab\times\vE=\Bigl(\pd{E_z}{y}-\pd{E_y}{z}\Bigr)\ux+\Bigl(\pd{E_x}{z}-\pd{E_z}{x}\Bigr)\uy+\Bigl(\pd{E_y}{x}-\pd{E_x}{y}\Bigr)\uz .
\]
Mide cuánto \emph{gira} el campo alrededor de un punto; si el campo de la \cref{fig:campos_rot} (izquierda) representa un movimiento de fluido, una pequeña rueda con paletas colocada en cualquier punto giraría, y el rotacional apunta en el eje
de giro, con el sentido que da la regla de la mano derecha (en el ejemplo, $+\uz$). Con el teorema de Stokes, $\iint_{S}(\nab\times\vE)\cdot\dS=\oint_{C}\vE\cdot\dl$: el flujo del rotacional por una superficie abierta es la \emph{suma total de remolinos}, que
se determina también observando cómo el campo sigue el contorno. A la integral de línea cerrada se la denomina \emph{circulación} del campo.

% Figura: campos con rotacional: F=-y ux + x uy (giro) y F=y ux (cizalla).
\begin{figure}[ht]
\centering
\begin{minipage}[t]{0.47\textwidth}\centering
\begin{tikzpicture}[x=0.58cm, y=0.58cm, baseline=(current bounding box.north)]
  \marcoD{-2.7}{2.7}{-2.7}{2.7}{-2,-1,0,1,2}{-2,-1,0,1,2}
  \ejesD{-2.7}{2.7}{-2.7}{2.7}
  \clipD{-2.7}{2.7}{-2.7}{2.7}{\flechas{-\y}{\x}{-2.4,-1.8,...,2.4}{-2.4,-1.8,...,2.4}{0.34}{3.4}}
\end{tikzpicture}\\[3pt]{\small $\vc{F}(\vc{r})=-y\,\ux+x\,\uy$}
\end{minipage}\hfill
\begin{minipage}[t]{0.47\textwidth}\centering
\begin{tikzpicture}[x=0.58cm, y=0.58cm, baseline=(current bounding box.north)]
  \marcoD{-2.7}{2.7}{-2.7}{2.7}{-2,-1,0,1,2}{-2,-1,0,1,2}
  \ejesD{-2.7}{2.7}{-2.7}{2.7}
  \clipD{-2.7}{2.7}{-2.7}{2.7}{\flechas{\y}{0}{-2.4,-1.8,...,2.4}{-2.4,-1.8,...,2.4}{0.34}{2.5}}
\end{tikzpicture}\\[3pt]{\small $\vc{F}(\vc{r})=y\,\ux$}
\end{minipage}
\caption{Campos con rotacional no nulo. (Izquierda) Giro rígido
$\vc{F}=-y\,\ux+x\,\uy$ (sentido antihorario): $\nab\times\vc{F}=2\,\uz$. Una
pequeña rueda colocada en cualquier punto gira con velocidad angular
constante. (Derecha) Cizalla $\vc{F}=y\,\ux$, de rotacional $-\uz$: no hay
líneas de corriente cerradas, pero la rueda gira igualmente (sentido horario)
porque el flujo es más rápido por un lado que por el otro.}
\label{fig:campos_rot}
\end{figure}


% -------------------------------------------------------------------------
\subsection{Dos identidades fundamentales}\label{app:identidades}
% -------------------------------------------------------------------------
Las aplicaciones reiteradas del operador $\nab$ llevan a dos identidades de gran interés en electromagnetismo.

\begin{proposicion}[Rotacional de un gradiente]\label{prop:rot-grad}
Para todo campo escalar $\phi$ de clase $C^{2}$, $\nab\times(\nab\phi)=\vc{0}$. Recíprocamente, si un campo vectorial es irrotacional en una región simplemente conexa, se puede expresar como el gradiente de un campo escalar.
\end{proposicion}
\begin{proof}
Por el teorema de Stokes y el del gradiente, para toda superficie $S$ de borde $C$,
\[
  \iint_{S}\bigl[\nab\times(\nab\phi)\bigr]\cdot\dS=\oint_{C}\nab\phi\cdot\dl=\oint_{C}\dif\phi=0,
\]
y como $S$ es arbitraria, el integrando es nulo. El recíproco es la proposición correspondiente del Tema~0~\cite{TEMA0}.
\end{proof}
Los campos irrotacionales se denominan \emph{conservativos} y siempre se pueden expresar como el gradiente de una función escalar. El campo electrostático es el ejemplo característico: $\nab\times\vE=\vc{0}$ y $\vE=-\nab\phi$.

\begin{proposicion}[Divergencia de un rotacional]\label{prop:div-rot}
Para todo campo vectorial $\vc{A}$ de clase $C^{2}$, $\nab\cdot(\nab\times\vc{A})=0$. Recíprocamente, un campo con divergencia nula en una región adecuada puede escribirse como rotacional de otro campo vectorial.
\end{proposicion}
\begin{proof}
Por el teorema de la divergencia, $\iiint_{V}\nab\cdot(\nab\times\vc{A})\,\dV=\oiint_{S}(\nab\times\vc{A})\cdot\dS$, con $S=\partial V$ cerrada. Cortemos $S$ por una curva cerrada $C$ en dos superficies abiertas $S_1$ y $S_2$ con el mismo borde $C$ y
orientaciones inducidas opuestas. Por Stokes, $\iint_{S_1}(\nabla\times\vc{A})\cdot\dS=\oint_{C}\vc{A}\cdot\dl$ y $\iint_{S_2}(\nabla\times\vc{A})\cdot\dS=-\oint_{C}\vc{A}\cdot\dl$; la suma es cero. Como $V$ es arbitrario, la divergencia del rotacional
es nula en todo punto.
\end{proof}
A los vectores de divergencia nula se les denomina \emph{solenoidales}. Un ejemplo característico es la inducción magnética, $\nab\cdot\vc{B}=0$, que se escribe en función del potencial vector magnético: $\vc{B}=\nab\times\vc{A}$.

% =========================================================================
\section{Formulario del tema}\label{app:formulario}
% =========================================================================
Reunimos las expresiones fundamentales del Tema~1 (cf.~\cite{FORMULARIO}).

\begin{center}
\captionof{table}{Resumen de fórmulas del Tema~1.}
\label{tab:formulario}
\small
\renewcommand{\arraystretch}{1.7}
\begin{tabular}{@{}p{4.2cm}p{10.2cm}@{}}
\toprule
\emph{Ley de Coulomb} & $\displaystyle\vF_{12}=\frac{1}{4\pi\eo}\frac{q_1q_2}{R_{12}^{2}}\,\uh{R_{12}},\quad\vc{R}_{12}=\vc{r}_2-\vc{r}_1$\\
Constantes & $K=\frac{1}{4\pi\eo}=8{,}988\cdot10^{9}\,\mathrm{N\,m^{2}/C^{2}}$; $\eo=8{,}85\cdot10^{-12}\,\mathrm{F/m}$; $c^{2}=\frac{1}{\eo\mu_0}$; $e=1{,}6\cdot10^{-19}\,\mathrm{C}$\\
Superposición & $\displaystyle\vF_q=\frac{q}{4\pi\eo}\sum_i\frac{q_i}{R_i^{2}}\,\uh{R_i}$\\
\midrule
\emph{Campo eléctrico} & $\displaystyle\vE=\lim_{q\to0}\frac{\vF_q}{q}=\frac{1}{4\pi\eo}\frac{\qp}{R^{2}}\,\uR$, \quad $\vc{R}=\vc{r}-\vc{r}\,'$\\
Distribuciones & $\displaystyle\vE=\frac{1}{4\pi\eo}\int_{V'}\rho\,\frac{\vc{R}}{R^{3}}\,\dVp$ \ (análogo con $\sigma\,\dSp$ y $\lambda\,\dlp$); $\ \rho=\frac{\dif q'}{\dVp}$, $\sigma=\frac{\dif q'}{\dSp}$, $\lambda=\frac{\dif q'}{\dlp}$\\
\midrule
\emph{Ángulo sólido} & $\displaystyle\dOm=\frac{\dS\cdot\uR}{R^{2}}$; \quad $\oint_S\dOm=4\pi$ (punto interior), $0$ (exterior)\\
\emph{Ley de Gauss} & $\displaystyle\oiint_{S}\vE\cdot\dS=\frac{\Qenc}{\eo},\qquad\nab\cdot\vE=\frac{\rho}{\eo}$\\
Delta de Dirac & $\displaystyle\nab\cdot\frac{\uR}{R^{2}}=4\pi\delta^{3}(\vc{R}),\quad\Lap\frac{1}{R}=-4\pi\delta^{3}(\vc{R}),\quad\int f(\vc{r})\delta^{3}(\vc{r}-\vc{a})\dV=f(\vc{a})$\\
\midrule
\emph{Potencial} & $\displaystyle\vE=-\nab\phi,\quad\phi=\frac{1}{4\pi\eo}\int\frac{\rho\,\dVp}{R}+\text{cte};\quad\phi(\vc{r}_2)-\phi(\vc{r}_1)=-\int_{1}^{2}\vE\cdot\dl$\\
Rotacional y circulación & $\displaystyle\nab\times\vE=\vc{0},\qquad\oint_C\vE\cdot\dl=0$\\
Poisson y Laplace & $\displaystyle\Lap\phi=-\frac{\rho}{\eo},\qquad\Lap\phi=0\ \ (\rho=0)$\\
Trabajo & $W_{AB}=q\,[\phi(B)-\phi(A)]$\\
\midrule
\emph{Condiciones de contorno} & $\displaystyle F_{2n}-F_{1n}=\lim_{h\to0}h\,\nab\cdot\vc{F},\quad\vc{F}_{2t}-\vc{F}_{1t}=\lim_{h\to0}h\,(\nab\times\vc{F})\times\un$\\
En electrostática & $\displaystyle\un\cdot(\vE_2-\vE_1)=\frac{\sigma}{\eo},\quad\un\times(\vE_2-\vE_1)=\vc{0},\quad\phi_2=\phi_1$\\
\bottomrule
\end{tabular}
\end{center}


% =========================================================================
%                              BIBLIOGRAFÍA
% =========================================================================
\printbibliography[title={Referencias}]

\end{document}
